% La conscience et l'inconscient — Philosophie Tle Tech — Cours signé OnBuch+
% Programme officiel MINESEC. Compiler avec : tectonic N04-conscience-inconscient.tex
\newcommand{\DOCMATIERE}{Philosophie}
\newcommand{\DOCNIVEAU}{Tle Tech}
\newcommand{\DOCMODULE}{Philosophie – classe de Terminale technique}
\newcommand{\DOCLECON}{N04}
\newcommand{\DOCTITRE}{La conscience et l'inconscient}
% OnBuch+ — préambule « cours signés » ENRICHI (compilé avec tectonic / XeLaTeX)
% Système visuel « Pop » d'OnBuch+ : fond crème (jamais blanc), bordures encre
% épaisses, ombres « dures » décalées, titres Archivo Black en capitales.
% Version MATHÉMATIQUES : graphes (pgfplots/tikz), boîtes pédagogiques
% (définition, propriété, méthode, attention, le savais-tu, exercice résolu,
% exercice, TP, bilan), page de garde et signature OnBuch+ ancrée en bas.
%
% Macros attendues dans main.tex AVANT \input{preamble} :
%   \DOCMATIERE  \DOCNIVEAU  \DOCMODULE  \DOCLECON (numéro)  \DOCTITRE
\documentclass[11pt]{article}

\usepackage{fontspec}
\usepackage[french]{babel}
\usepackage[a4paper,margin=2cm,top=3.4cm,bottom=2.8cm,headheight=1.4cm,footskip=1.3cm]{geometry}
\usepackage{amsmath,amssymb,amsfonts}
\usepackage{mathtools} % \xmapsto, \coloneqq... souvent utilisés par les modèles
\usepackage{cancel} % simplifications barrées (bilans de forces)
% \abs{z} (module, valeur absolue) et \norm{u} (norme), utilisés par les modèles
\providecommand{\abs}{}\let\abs\relax\DeclarePairedDelimiter{\abs}{\lvert}{\rvert}
\providecommand{\norm}{}\let\norm\relax\DeclarePairedDelimiter{\norm}{\lVert}{\rVert}
\usepackage[table]{xcolor} % [table] : fournit aussi \rowcolors (alternance de lignes)
\usepackage{pagecolor}
\usepackage{tikz}
\usetikzlibrary{arrows.meta,positioning,calc,shapes.geometric,shapes.misc,shapes.arrows,decorations.pathmorphing,decorations.pathreplacing,decorations.markings,patterns,shadows,mindmap,trees,fit,backgrounds,matrix,intersections,angles,quotes}
\usepackage{pgfplots}
\usepackage[version=4]{mhchem} % équations nucléaires \ce{^{235}_{92}U}
\usepackage[european,siunitx]{circuitikz} % schémas électriques
\pgfplotsset{compat=1.18}
\usepgfplotslibrary{fillbetween,groupplots} % aires entre courbes (calcul intégral)
\usepackage{enumitem}
\usepackage{fancyhdr}
% [most] charge toutes les bibliothèques tcolorbox (listings, minted, raster,
% documentation...) et gonfle inutilement la mémoire TeX sur un document long
% (~300 boîtes) : on ne charge que ce qui est réellement utilisé.
\usepackage[skins,breakable]{tcolorbox}
\usepackage{graphicx}
\usepackage{colortbl}
\usepackage{booktabs}
\usepackage{tabularx}
\usepackage{multirow}
\usepackage{diagbox} % case d'en-tête diagonale des tableaux à double entrée
% Colonnes extensibles centrée / gauche / droite (C, L, R), souvent utilisées par les modèles
\newcolumntype{C}{>{\centering\arraybackslash}X}
\newcolumntype{L}{>{\raggedright\arraybackslash}X}
\newcolumntype{R}{>{\raggedleft\arraybackslash}X}
\usepackage{array}
\usepackage{multicol}
\usepackage{siunitx}
% parse-numbers=false : accepte aussi \SI{2,5 \times 10^{-3}}{...}
\sisetup{locale=FR,output-decimal-marker={,},group-minimum-digits=4,parse-numbers=false}
\DeclareSIUnit\ohm{\ensuremath{\symup{\Omega}}} % U+2126 absent de la police
\DeclareSIUnit\division{div} % réglages d'oscilloscope (ms/div, V/div)
\DeclareSIUnit\tr{tr} % tours (vitesse de rotation, ex. \SI{1500}{\tr\per\minute})
\DeclareSIUnit\molar{M} % concentration molaire (ex. \SI{3}{\molar}, \SI{1}{\micro\molar})
\DeclareSIUnit\an{an} % année (le modèle confond parfois avec \year, un compteur TeX) — ex. \SI{5}{\kilo\meter\per\an}
\DeclareSIUnit\FCFA{FCFA} % franc CFA (ex. \SI{500}{\FCFA\per\kilo\gram})
\DeclareSIUnit\degreeBrix{\textdegree Brix} % teneur en sucre d'un sirop/jus (réfractométrie)
\DeclareSIUnit\month{mois} % le modèle utilise parfois \month, un compteur TeX, comme unité de durée
\DeclareSIUnit\microliter{\micro\liter} % le modèle écrit parfois \microliter en un seul token
\DeclareSIUnit\million{millions} % dénombrement (ex. \SI{15}{\million\per\mL} spermatozoïdes)
\usepackage{qrcode}
% \degree : commande fréquemment utilisée par les modèles pour un angle ou une
% température en dehors de \SI{}{} (non fournie par les paquets chargés).
\providecommand{\degree}{^{\circ}}
% \qed (fin de démonstration), utilisé par les modèles sans amsthm
\providecommand{\qed}{\hfill$\square$}
% \barre : double barre des valeurs interdites dans un tableau de variations (tkz-tab)
\providecommand{\barre}{\ensuremath{\|}}
% environnement proof (amsthm non chargé) utilisé par les modèles
\ifdefined\proof\else\newenvironment{proof}[1][Démonstration]{\par\noindent\textit{#1.}\ }{\qed\par}\fi
\usepackage{lastpage}
\usepackage{hyperref}
\hypersetup{hidelinks}

% ============================================================
% Palette « Pop » OnBuch+ — mêmes hex que la classe P (plus_ui.dart)
% ============================================================
\definecolor{poporange}{HTML}{F59321}
\definecolor{popdark}{HTML}{E07A0C}
\definecolor{popcream}{HTML}{FFF6E8}
\definecolor{popink}{HTML}{1C1714}
\definecolor{poppurple}{HTML}{7A5AE0}
\definecolor{popgreen}{HTML}{2E9E6A}
\definecolor{poppink}{HTML}{E0457B}
\definecolor{popblue}{HTML}{2F7DE1}
\definecolor{popgold}{HTML}{C98A12}
\definecolor{popmuted}{HTML}{8A7F76}
\definecolor{popline}{HTML}{EADFD2}
\definecolor{popgray}{HTML}{8A7F76}
\colorlet{popgrey}{popgray}
% Alias tolérants (le modèle nomme parfois les couleurs différemment)
\colorlet{popwhite}{white}
\colorlet{popblack}{popink}
\colorlet{popred}{poppink}
\colorlet{popyellow}{popgold}
% Teintes claires (fonds de tableaux/encadrés) — le modèle nomme parfois
% les couleurs avec un suffixe L (ex. popblueL) sans les avoir définies.
\colorlet{poporangeL}{poporange!20}
\colorlet{popdarkL}{popdark!20}
\colorlet{poppurpleL}{poppurple!20}
\colorlet{popgreenL}{popgreen!20}
\colorlet{poppinkL}{poppink!20}
\colorlet{popblueL}{popblue!20}
\colorlet{popgoldL}{popgold!20}
\colorlet{popredL}{popred!20}
\colorlet{popgrayL}{popgray!20}
% Teintes claires pour fonds de tableaux / zones
\colorlet{popblueL}{popblue!10!white}
\colorlet{poporangeL}{poporange!14!white}
\colorlet{popgreenL}{popgreen!12!white}
\colorlet{poppurpleL}{poppurple!12!white}
\colorlet{poppinkL}{poppink!12!white}
\colorlet{popgoldL}{popgold!14!white}

\pagecolor{popcream}

% ============================================================
% Typographie
% ============================================================
\defaultfontfeatures{Ligatures=TeX}
\setmainfont{PlusJakartaSans-Medium.ttf}[Path=../../../fonts/,BoldFont=PlusJakartaSans-Bold.ttf]
% Maths en sans-serif (Fira Math) avec chiffres et lettres droites de Plus
% Jakarta, pour que les formules se fondent dans le texte.
\usepackage[math-style=ISO,bold-style=ISO]{unicode-math}
\setmathfont{FiraMath-Regular.otf}[Path=../../../fonts/]
\setmathfont{PlusJakartaSans-Medium.ttf}[Path=../../../fonts/,range={up/num,up/{Latin,latin},\mathup}]
\usepackage{icomma} % virgule décimale sans espace en mode math
% Caractères mathématiques Unicode absents des polices de texte (Plus Jakarta,
% Archivo Black) : rendus par leur équivalent mathématique, en texte comme en
% formule — sinon ils s'affichent en carré vide (ex. « Arithmétique dans ℤ »).
\usepackage{newunicodechar}
\newunicodechar{ℝ}{\ensuremath{\mathbb{R}}}
\newunicodechar{ℤ}{\ensuremath{\mathbb{Z}}}
\newunicodechar{ℂ}{\ensuremath{\mathbb{C}}}
\newunicodechar{ℕ}{\ensuremath{\mathbb{N}}}
\newunicodechar{ℚ}{\ensuremath{\mathbb{Q}}}
\newunicodechar{𝔻}{\ensuremath{\mathbb{D}}}
\newunicodechar{α}{\ensuremath{\alpha}}
\newunicodechar{β}{\ensuremath{\beta}}
\newunicodechar{γ}{\ensuremath{\gamma}}
\newunicodechar{δ}{\ensuremath{\delta}}
\newunicodechar{ε}{\ensuremath{\varepsilon}}
\newunicodechar{θ}{\ensuremath{\theta}}
\newunicodechar{λ}{\ensuremath{\lambda}}
\newunicodechar{μ}{\ensuremath{\mu}}
\newunicodechar{σ}{\ensuremath{\sigma}}
\newunicodechar{τ}{\ensuremath{\tau}}
\newunicodechar{φ}{\ensuremath{\varphi}}
\newunicodechar{ψ}{\ensuremath{\psi}}
\newunicodechar{ω}{\ensuremath{\omega}}
\newunicodechar{Σ}{\ensuremath{\Sigma}}
% majuscules produites par \MakeUppercase dans les titres (α-aminés -> Α)
\newunicodechar{Α}{\ensuremath{\alpha}}
\newunicodechar{Β}{\ensuremath{\beta}}
\newunicodechar{Γ}{\ensuremath{\Gamma}}
\newunicodechar{Δ}{\ensuremath{\Delta}}
\newunicodechar{Ε}{\ensuremath{\varepsilon}}
\newunicodechar{Θ}{\ensuremath{\Theta}}
\newunicodechar{Λ}{\ensuremath{\Lambda}}
\newunicodechar{Μ}{\ensuremath{\mu}}
\newunicodechar{Τ}{\ensuremath{\tau}}
\newunicodechar{Φ}{\ensuremath{\Phi}}
\newunicodechar{Ψ}{\ensuremath{\Psi}}
\newunicodechar{Ω}{\ensuremath{\Omega}}
\newunicodechar{↦}{\ensuremath{\mapsto}}
\newunicodechar{∈}{\ensuremath{\in}}
\newunicodechar{∉}{\ensuremath{\notin}}
\newunicodechar{∧}{\ensuremath{\wedge}}
\newunicodechar{⇔}{\ensuremath{\Leftrightarrow}}
\newunicodechar{⟺}{\ensuremath{\Longleftrightarrow}}
\newunicodechar{⇒}{\ensuremath{\Rightarrow}}
\newunicodechar{⟹}{\ensuremath{\Longrightarrow}}
\newunicodechar{∘}{\ensuremath{\circ}}
\newunicodechar{∪}{\ensuremath{\cup}}
\newunicodechar{∩}{\ensuremath{\cap}}
\newunicodechar{≡}{\ensuremath{\equiv}}
\newunicodechar{⊂}{\ensuremath{\subset}}
\newunicodechar{⊥}{\ensuremath{\perp}}
\newunicodechar{∃}{\ensuremath{\exists}}
\newunicodechar{∀}{\ensuremath{\forall}}
\newunicodechar{∛}{\ensuremath{\sqrt[3]{\,}}}
\newunicodechar{∜}{\ensuremath{\sqrt[4]{\,}}}
\newunicodechar{⃗}{\ensuremath{{}^{\to}}}
\newunicodechar{ⁿ}{\ifmmode^{n}\else\textsuperscript{\ensuremath{n}}\fi}
\newunicodechar{ˣ}{\ifmmode^{x}\else\textsuperscript{\ensuremath{x}}\fi}
\newunicodechar{ᵇ}{\ifmmode^{b}\else\textsuperscript{\ensuremath{b}}\fi}
\newunicodechar{ʳ}{\ifmmode^{r}\else\textsuperscript{\ensuremath{r}}\fi}
\newunicodechar{ᵘ}{\ifmmode^{u}\else\textsuperscript{\ensuremath{u}}\fi}
\newunicodechar{ʸ}{\ifmmode^{y}\else\textsuperscript{\ensuremath{y}}\fi}
\newunicodechar{ᵃ}{\ifmmode^{a}\else\textsuperscript{\ensuremath{a}}\fi}
\newunicodechar{ᵉ}{\ifmmode^{e}\else\textsuperscript{\ensuremath{e}}\fi}
\newunicodechar{ᵏ}{\ifmmode^{k}\else\textsuperscript{\ensuremath{k}}\fi}
\newunicodechar{ᵖ}{\ifmmode^{p}\else\textsuperscript{\ensuremath{p}}\fi}
\newunicodechar{ᴺ}{\ifmmode^{N}\else\textsuperscript{\ensuremath{N}}\fi}
\newunicodechar{ᶜ}{\ifmmode^{c}\else\textsuperscript{\ensuremath{c}}\fi}
\newunicodechar{ʰ}{\ifmmode^{h}\else\textsuperscript{\ensuremath{h}}\fi}
\newunicodechar{ᵠ}{\ifmmode^{q}\else\textsuperscript{\ensuremath{q}}\fi}
\newunicodechar{ₙ}{\ifmmode_{n}\else\textsubscript{\ensuremath{n}}\fi}
\newunicodechar{ₐ}{\ifmmode_{a}\else\textsubscript{\ensuremath{a}}\fi}
\newunicodechar{ᵢ}{\ifmmode_{i}\else\textsubscript{\ensuremath{i}}\fi}
\newunicodechar{ᵣ}{\ifmmode_{r}\else\textsubscript{\ensuremath{r}}\fi}
\newunicodechar{ₖ}{\ifmmode_{k}\else\textsubscript{\ensuremath{k}}\fi}
\newunicodechar{ᵦ}{\ifmmode_{\beta}\else\textsubscript{\ensuremath{\beta}}\fi}
\newunicodechar{ₓ}{\ifmmode_{x}\else\textsubscript{\ensuremath{x}}\fi}
\newfontfamily\popdisplay{ArchivoBlack-Regular.ttf}[Path=../../../fonts/]
\newfontfamily\poplogo{SpaceGrotesk-Bold.ttf}[Path=../../../fonts/]
\newfontfamily\popsemi{PlusJakartaSans-SemiBold.ttf}[Path=../../../fonts/]
\newfontfamily\popxbold{PlusJakartaSans-ExtraBold.ttf}[Path=../../../fonts/]
\newfontfamily\popxboldls{PlusJakartaSans-ExtraBold.ttf}[Path=../../../fonts/,LetterSpace=3.0]

\setlist[enumerate]{leftmargin=1.7em,itemsep=5pt,topsep=4pt}
\setlist[itemize]{leftmargin=1.7em,itemsep=4pt,topsep=4pt,label={\color{poporange}\textbullet}}
\setlength{\parindent}{0pt}
\setlength{\parskip}{5pt}
\renewcommand{\arraystretch}{1.25}

% pgfplots : style Pop par défaut
\pgfplotsset{
  every axis/.append style={axis line style={popink,line width=1pt},tick style={popink},
    grid style={popline,line width=0.5pt},label style={font=\small},tick label style={font=\footnotesize}},
  popcurve/.style={poporange,line width=1.8pt,no markers,smooth},
}
% Même style disponible dans un \draw TikZ simple (hors pgfplots)
\tikzset{popcurve/.style={poporange,line width=1.8pt}}
\tikzset{popgrid/.style={popline,very thin}}
\colorlet{popcurve}{poporange} % le nom du style est parfois utilisé comme couleur

% ============================================================
% Pill / squiggle
% ============================================================
\newcommand{\poppill}[3]{%
  \tikz[baseline=-0.55ex]\node[draw=popink,line width=1.2pt,rounded corners=2.6pt,%
    fill=#2,inner xsep=6.5pt,inner ysep=3pt]{%
    \popxboldls\fontsize{7}{7}\selectfont\color{#3}\MakeUppercase{#1}};%
}
\newcommand{\popsquiggle}[1]{%
  \tikz[baseline=-0.4ex]{
    \draw[#1,line width=2.6pt,line cap=round]
      plot [smooth] coordinates {(0,0.09) (0.34,0.01) (0.68,0.21) (1.02,0.01) (1.36,0.10)};
  }%
}
% Mot-clé surligné (marqueur orange sous le texte)
\newcommand{\cle}[1]{{\popxbold\color{popdark}#1}}

% ============================================================
% En-tête / pied
% ============================================================
\pagestyle{fancy}
\fancyhf{}
\renewcommand{\headrulewidth}{0pt}
\renewcommand{\footrulewidth}{0pt}
\lhead{\raisebox{-0.15ex}{\poplogo\fontsize{13}{13}\selectfont\color{popink}OnBuch\color{poporange}+}}
\rhead{\poppill{\DOCMATIERE\ \textbullet\ \DOCNIVEAU\ \textbullet\ Leçon \DOCLECON}{white}{popink}}
\lfoot{\popxbold\fontsize{7.6}{7.6}\selectfont\color{popmuted}\MakeUppercase{Cours signé OnBuch+}\ \textbullet\ {\popsemi\DOCTITRE}}
\rfoot{\tikz[baseline=-0.6ex]\node[rounded corners=8.5pt,fill=popink,minimum height=17pt,minimum width=17pt,inner xsep=5pt,inner sep=0]{\popxbold\fontsize{8}{8}\selectfont\color{popcream}\thepage\,/\,\pageref*{LastPage}};}

% ============================================================
% Titres
% ============================================================
\newcommand{\chaptitre}[2]{%
  \begin{tcolorbox}[enhanced,colback=poporange,colframe=popink,boxrule=2.6pt,arc=11pt,
      drop shadow=popink,left=15pt,right=15pt,top=12pt,bottom=11pt,before skip=3pt,after skip=18pt]
    \poppill{#2}{popink}{popcream}\par\vspace{9pt}
    {\raggedright\hyphenpenalty=10000\exhyphenpenalty=10000\popdisplay\fontsize{22}{24}\selectfont\color{popcream}\MakeUppercase{#1}\par}
  \end{tcolorbox}
}
% Section numérotée : pastille orange avec le numéro + titre Archivo
\newcounter{popsec}
\newcommand{\coursec}[1]{%
  \stepcounter{popsec}\setcounter{popsub}{0}%
  \par\vspace{16pt}\noindent
  \begin{minipage}{\linewidth}
  \tikz[baseline=-0.7ex]\node[draw=popink,line width=1.6pt,fill=poporange,rounded corners=4pt,minimum size=22pt,inner sep=0]{\popdisplay\fontsize{12}{12}\selectfont\color{popcream}\thepopsec};%
  \hspace{8pt}{\popdisplay\fontsize{15}{17}\selectfont\color{popink}\MakeUppercase{#1}}\par
  \vspace{4pt}\noindent{\color{poporange}\rule{36pt}{3.6pt}}
  \end{minipage}\par\nopagebreak\vspace{8pt}
}
\newcounter{popsub}
\newcommand{\courssub}[1]{%
  \stepcounter{popsub}%
  \par\vspace{9pt}\noindent{\popxbold\fontsize{11.5}{13}\selectfont\color{popink}{\color{poporange}\thepopsec.\thepopsub}\hspace{6pt}#1}\par\nopagebreak\vspace{4pt}
}

% ============================================================
% Boîtes pédagogiques — même recette : fond blanc, bordure épaisse colorée,
% ombre dure encre, badge plein. Titre optionnel [#1].
% ============================================================
\tcbset{popbase/.style={breakable,enhanced,colback=white,boxrule=2.3pt,arc=9pt,
  shadow={3pt}{-3pt}{0pt}{popink},left=14pt,right=14pt,top=11pt,bottom=11pt,before skip=11pt,after skip=15pt,
  fontupper=\popsemi\fontsize{10.6}{14.5}\selectfont\color{popink},
  fontlower=\popsemi\fontsize{10.6}{14.5}\selectfont\color{popink}}}
% Coche dessinée (le glyphe ✓ n'existe pas dans les polices du document)
\newcommand{\popcheck}[1]{\tikz[baseline=-0.5ex]\draw[#1,line width=1.6pt,line cap=round,line join=round](-0.13,0.0)--(-0.04,-0.09)--(0.13,0.1);}
\providecommand{\coche}{\popcheck{popgreen}}
\providecommand{\resultat}[1]{\par\smallskip\noindent\fcolorbox{popgreen}{popgreenL}{\parbox{\dimexpr\linewidth-2\fboxsep-2\fboxrule\relax}{\textbf{Résultat.} #1}}\par\smallskip}
\newcommand{\popboxhead}[3]{\poppill{#1}{#2}{white}\ifx\relax#3\relax\else\hspace{6pt}{\popxbold\fontsize{10.6}{12}\selectfont\color{popink}#3}\fi\par\vspace{7pt}}

\newtcolorbox{aretenir}[1][]{popbase,colframe=popgreen,before upper={\popboxhead{À retenir}{popgreen}{#1}}}
\newtcolorbox{exemplebox}[1][]{popbase,colframe=poporange,before upper={\popboxhead{Exemple}{poporange}{#1}}}
\newtcolorbox{definition}[1][]{popbase,colframe=popblue,before upper={\popboxhead{Définition}{popblue}{#1}}}
\newtcolorbox{propriete}[1][]{popbase,colframe=poppurple,before upper={\popboxhead{Propriété}{poppurple}{#1}}}
\newtcolorbox{methode}[1][]{popbase,colframe=popgold,before upper={\popboxhead{Méthode}{popgold}{#1}}}
\newtcolorbox{attention}[1][]{popbase,colframe=poppink,before upper={\popboxhead{Attention}{poppink}{#1}}}
\newtcolorbox{experience}[1][]{popbase,colframe=popblue,colback=popblueL,before upper={\popboxhead{Expérience}{popblue}{#1}}}
% « Le savais-tu ? » : carte violette pleine (contexte, culture, Cameroun)
\newtcolorbox{savaistu}[1][]{popbase,colback=poppurple,colframe=popink,
  fontupper=\popsemi\fontsize{10.4}{14.2}\selectfont\color{white},
  before upper={\poppill{Le savais-tu\,?}{white}{poppurple}\ifx\relax#1\relax\else\hspace{6pt}{\popxbold\color{white}#1}\fi\par\vspace{7pt}}}
% Exercice résolu : énoncé en haut, \tcblower puis la solution sur fond crème
\newcounter{popexo}\newcounter{popexr}
\newtcolorbox{exoresolu}[1][]{popbase,colframe=poporange,colbacklower=popcream,
  segmentation style={popink,line width=1.2pt,dashed},
  before upper={\stepcounter{popexr}\popboxhead{Exercice résolu \thepopexr}{poporange}{#1}},
  before lower={\poppill{Solution}{popink}{popcream}\par\vspace{6pt}}}
% Exercice (non résolu en place) — la correction est dans la partie Corrigés
\newtcolorbox{exercice}[1][]{popbase,colframe=popink,
  before upper={\stepcounter{popexo}\popboxhead{Exercice \thepopexo}{popink}{#1}}}
% Corrigé
\newtcolorbox{corrige}[1][]{popbase,colframe=popgreen,colback=popgreenL,
  before upper={\popboxhead{Corrigé}{popgreen}{#1}}}
% Objectifs / prérequis / bilan
\newtcolorbox{objectifs}{popbase,colframe=popink,colback=poporangeL,
  before upper={\popboxhead{Objectifs de la leçon}{popink}{}}}
\newtcolorbox{prerequis}{popbase,colframe=popink,colback=popblueL,
  before upper={\popboxhead{Prérequis}{popblue}{}}}
\newtcolorbox{bilan}{popbase,colframe=popink,colback=white,boxrule=2.8pt,
  before upper={\popboxhead{Fiche bilan}{poporange}{L'essentiel à connaître pour l'examen}}}
% Figure encadrée avec légende
\newtcolorbox{popfigure}[1][]{enhanced,breakable=false,colback=white,colframe=popline,boxrule=1.4pt,arc=8pt,
  left=10pt,right=10pt,top=10pt,bottom=8pt,before skip=10pt,after skip=12pt,halign=center,
  lower separated=false,halign lower=center,
  fontlower=\popsemi\fontsize{9}{11.5}\selectfont\color{popmuted},
  #1}
\newcommand{\legende}[1]{\par\vspace{6pt}{\popsemi\fontsize{9}{11.5}\selectfont\color{popmuted}#1}}

% ============================================================
% Signature OnBuch+
% ============================================================
\newcommand{\poplogoinline}[1]{{\poplogo\fontsize{#1}{#1}\selectfont\color{popink}OnBuch\color{poporange}+}}
\newcommand{\signatureonbuch}{%
  \begin{tcolorbox}[enhanced,colback=popink,colframe=popink,boxrule=2.2pt,arc=10pt,
      drop shadow=poporange,left=14pt,right=14pt,top=10pt,bottom=10pt,before skip=0pt,after skip=0pt]
    \tikz[baseline=-0.7ex]\node[circle,fill=popgreen,draw=popcream,line width=1.3pt,minimum size=22pt,inner sep=0]{\popcheck{white}};%
    \hspace{9pt}\begin{minipage}[c]{0.62\linewidth}
      {\popxbold\fontsize{12}{13}\selectfont\color{popcream}Signé\ \textbullet\ {\poplogo OnBuch\color{poporange}+}}\par\vspace{2pt}
      {\raggedright\popsemi\fontsize{8}{10.5}\selectfont\color{popline}\MakeUppercase{Équipe pédagogique OnBuch+}\\\MakeUppercase{Conforme au programme officiel MINESEC}\par}
    \end{minipage}\hfill
    {\poplogo\fontsize{22}{22}\selectfont\color{popcream}OnBuch\color{poporange}+}
  \end{tcolorbox}%
}

% ============================================================
% Page de garde
% #1 titre · #2 matière · #3 niveau · #4 module · #5 n° leçon · #6 durée ·
% #7 description / objectifs (texte ou itemize)
% ============================================================
\newcommand{\pagegarde}[7]{%
  \thispagestyle{empty}
  \newgeometry{a4paper,margin=1.7cm,top=1.6cm,bottom=1.4cm}%
  \begin{tikzpicture}[remember picture,overlay]
    \fill[poporange!18] ([xshift=-3.2cm,yshift=-3.6cm]current page.north east) circle (4.6cm);
    \fill[poppurple!14] ([xshift=2.4cm,yshift=7.6cm]current page.south west) circle (3.8cm);
  \end{tikzpicture}%
  {\poplogo\fontsize{30}{30}\selectfont\color{popink}OnBuch\color{poporange}+}\hfill
  \raisebox{0.6ex}{\poppill{Cours signé \textbullet\ Premium}{popink}{popcream}}\par
  \vspace{1pt}\popsquiggle{poppurple}\par
  \vspace{8pt}
  \begin{tcolorbox}[enhanced,colback=poporange,colframe=popink,boxrule=2.8pt,arc=13pt,
      drop shadow=popink,left=18pt,right=18pt,top=12pt,bottom=13pt,after skip=10pt]
    \poppill{#2 \textbullet\ #3}{popink}{popcream}\hspace{5pt}\poppill{#4}{white}{popink}\par\vspace{8pt}
    {\popdisplay\fontsize{14}{14}\selectfont\color{popink}LEÇON #5}\par\vspace{4pt}
    {\raggedright\hyphenpenalty=10000\exhyphenpenalty=10000\popdisplay\fontsize{25}{27}\selectfont\color{popcream}\MakeUppercase{#1}\par}
    \vspace{8pt}
    {\popxbold\fontsize{9.5}{9.5}\selectfont\color{popink}\MakeUppercase{Durée conseillée : #6}}
  \end{tcolorbox}
  \begin{tcolorbox}[enhanced,colback=white,colframe=popink,boxrule=2.2pt,arc=10pt,
      drop shadow=popink,left=16pt,right=16pt,top=10pt,bottom=10pt,after skip=8pt]
    \poppill{Dans cette leçon}{poppurple}{white}\par\vspace{5pt}
    \popsemi\fontsize{9.3}{12.6}\selectfont\color{popink}#7
  \end{tcolorbox}
  \noindent\begin{minipage}[t]{0.487\linewidth}
    \begin{tcolorbox}[enhanced,colback=popgreen,colframe=popink,boxrule=2.2pt,arc=10pt,
        drop shadow=popink,left=11pt,right=11pt,top=10pt,bottom=10pt,height=3.1cm,valign=center]
      \tcbox[colback=white,colframe=popink,boxrule=1.3pt,arc=5pt,boxsep=0pt,nobeforeafter,
        left=3pt,right=3pt,top=3pt,bottom=3pt,valign=center]{\qrcode[height=1.9cm]{https://play.google.com/store/apps/details?id=com.luvvix.onbuch&hl=fr}}%
      \hspace{8pt}\begin{minipage}[c]{0.5\linewidth}
        {\popdisplay\fontsize{11}{12.5}\selectfont\color{white}\MakeUppercase{Télécharge OnBuch}}\par\vspace{4pt}
        {\raggedright\popsemi\fontsize{8.4}{11}\selectfont\color{white}Gratuit \textbullet\ Google Play\\Tuteur IA, annales, concours, résultats\par}
      \end{minipage}
    \end{tcolorbox}
  \end{minipage}\hfill
  \begin{minipage}[t]{0.487\linewidth}
    \begin{tcolorbox}[enhanced,colback=poppurple,colframe=popink,boxrule=2.2pt,arc=10pt,
        drop shadow=popink,left=11pt,right=11pt,top=10pt,bottom=10pt,height=3.1cm,valign=center]
      \tikz[baseline=-0.7ex]\node[circle,fill=white,draw=popink,line width=1.3pt,minimum size=1.6cm,inner sep=0]{\popdisplay\fontsize{24}{24}\selectfont\color{poppurple}+};%
      \hspace{8pt}\begin{minipage}[c]{0.6\linewidth}
        {\popdisplay\fontsize{11}{12.5}\selectfont\color{white}\MakeUppercase{Passe à OnBuch+}}\par\vspace{4pt}
        {\raggedright\popsemi\fontsize{8.4}{11}\selectfont\color{white}Tous les cours signés, Tuteur IA illimité, fiches PDF — onglet OnBuch+ de l'app\par}
      \end{minipage}
    \end{tcolorbox}
  \end{minipage}
  \vspace{8pt}
  \popcontenu
  \vfill
  \signatureonbuch
  \restoregeometry
  \newpage
}

% Rangée « Ce cours contient » (page de garde)
\newcommand{\poptile}[3]{% #1 couleur #2 gros texte #3 légende
  \begin{tcolorbox}[enhanced,colback=#1,colframe=popink,boxrule=2pt,arc=9pt,shadow={2.5pt}{-2.5pt}{0pt}{popink},
      left=6pt,right=6pt,top=9pt,bottom=9pt,halign=center,valign=center,height=1.75cm,width=\linewidth,nobeforeafter]
    {\popdisplay\fontsize{12.5}{13}\selectfont\color{white}\MakeUppercase{#2}}\par\vspace{3pt}
    {\centering\popsemi\fontsize{7.8}{9.5}\selectfont\color{white}#3\par}
  \end{tcolorbox}}
\newcommand{\popcontenu}{%
  \par\noindent{\popxboldls\fontsize{7.5}{7.5}\selectfont\color{popmuted}CE COURS SIGNÉ CONTIENT}\par\vspace{7pt}
  \noindent\begin{minipage}[t]{0.235\linewidth}\poptile{popblue}{Cours}{complet, illustré, pas à pas}\end{minipage}\hfill
  \begin{minipage}[t]{0.235\linewidth}\poptile{poporange}{Méthodes}{et exercices résolus}\end{minipage}\hfill
  \begin{minipage}[t]{0.235\linewidth}\poptile{poppink}{Exercices}{type examen + corrigés}\end{minipage}\hfill
  \begin{minipage}[t]{0.235\linewidth}\poptile{popgreen}{Fiche bilan}{l'essentiel à retenir}\end{minipage}\par}

% Sommaire
\newtcolorbox{sommaire}{enhanced,colback=white,colframe=popink,boxrule=2.2pt,arc=10pt,
  drop shadow=popink,left=18pt,right=18pt,top=13pt,bottom=13pt,before skip=6pt,after skip=20pt,
  before upper={\poppill{Sommaire}{poppurple}{white}\par\vspace{9pt}\popsemi\fontsize{10.6}{14.5}\selectfont\color{popink}}}

% Signature de fin de cours — poussée tout en bas de la dernière page
\newcommand{\findecours}{%
  \par\vfill\nopagebreak
  \begin{center}\popsquiggle{poporange}\end{center}\vspace{4pt}
  \signatureonbuch
}
\AtBeginDocument{\def\llbracket{\mathopen{[\mkern-3.5mu[}}\def\rrbracket{\mathclose{]\mkern-3.5mu]}}} % intervalles d'entiers (glyphe ⟦ absent de la police)
% Glyphes absents de Fira Math : redéfinis à partir de glyphes disponibles
\AtBeginDocument{%
  \def\setminus{\mathbin{\backslash}}\let\smallsetminus\setminus
  \def\vdots{\mathord{\vbox{\baselineskip3.2pt\lineskiplimit0pt\kern4pt\hbox{.}\hbox{.}\hbox{.}}}}%
  \def\ddots{\mathinner{\mkern1mu\raise6.5pt\hbox{.}\mkern2mu\raise3.6pt\hbox{.}\mkern2mu\raise0.7pt\hbox{.}\mkern1mu}}%
  \def\triangle{\mathord{\scalebox{0.9}{$\symup{\Delta}$}}}%
  \def\nabla{\mathord{\raisebox{\depth}{\scalebox{1}[-1]{$\symup{\Delta}$}}}}%
  \def\checkmark{\popcheck{popdark}}%
  \def\star{\mathbin{\ast}}\def\bigstar{\mathord{\ast}}%
  \def\diamond{\mathbin{\scalebox{0.6}{$\Diamond$}}}%
  \def\bigcup{\mathop{\mathchoice{\raisebox{-0.3em}{\scalebox{1.7}{$\cup$}}}{\scalebox{1.25}{$\cup$}}{\cup}{\cup}}\displaylimits}%
  \def\bigcap{\mathop{\mathchoice{\raisebox{-0.3em}{\scalebox{1.7}{$\cap$}}}{\scalebox{1.25}{$\cap$}}{\cap}{\cap}}\displaylimits}%
  \def\lesseqqgtr{\mathrel{\lessgtr}}\def\bigoplus{\mathop{\oplus}\displaylimits}%
}
\providecommand{\ssi}{si et seulement si} % abréviation fréquente
\AtBeginDocument{\providecommand{\wideparen}{\overparen}} % arc de cercle

\begin{document}

\pagegarde{La conscience et l'inconscient}{Philosophie}{Tle Tech}{Philosophie – classe de Terminale technique}{N04}{2 séances (fiche de progression harmonisée)}{Cette leçon étudie la conscience (nature, valeur, limites) et l'inconscient freudien, afin de déterminer dans quelle mesure l'homme est maître de lui-même. Elle évalue ensuite la valeur et les limites de la psychanalyse comme méthode de connaissance de l'esprit humain.
\vspace{4pt}
\begin{multicols}{2}\small
\begin{itemize}
\item À la fin de cette leçon, je sais définir la conscience, l'inconscient et la psychanalyse avec précision
\item Je sais distinguer conscience spontanée, conscience réfléchie et conscience morale
\item Je sais analyser la valeur de la conscience comme connaissance de soi et comme fondement de la responsabilité
\item Je sais identifier les limites de la conscience (erreurs des sens, illusions, mauvaise foi, déterminismes sociaux)
\item Je sais expliquer l'hypothèse freudienne de l'inconscient à partir d'exemples concrets (rêves, lapsus, actes manqués)
\item Je sais interpréter un cas concret de la vie courante camerounaise à l'aide des notions de conscience et d'inconscient
\end{itemize}\end{multicols}}

\begin{sommaire}
\begin{multicols}{2}
\begin{enumerate}
\item Problématique et définitions
\item Nature de la conscience
\item Valeur de la conscience
\item Limites de la conscience
\item L'inconscient freudien
\item Valeur et limites de la psychanalyse
\item Synthèse, méthode et entraînement au Bac
\item Activité d'intégration
\item Méthodes et exercices résolus
\item Exercices
\item Corrigés des exercices
\item Fiche bilan
\end{enumerate}
\end{multicols}
\end{sommaire}

\begin{objectifs}
\begin{itemize}
\item À la fin de cette leçon, je sais définir la conscience, l'inconscient et la psychanalyse avec précision
\item Je sais distinguer conscience spontanée, conscience réfléchie et conscience morale
\item Je sais analyser la valeur de la conscience comme connaissance de soi et comme fondement de la responsabilité
\item Je sais identifier les limites de la conscience (erreurs des sens, illusions, mauvaise foi, déterminismes sociaux)
\item Je sais expliquer l'hypothèse freudienne de l'inconscient à partir d'exemples concrets (rêves, lapsus, actes manqués)
\item Je sais interpréter un cas concret de la vie courante camerounaise à l'aide des notions de conscience et d'inconscient
\item Je sais évaluer la valeur et les limites de la psychanalyse comme science et comme thérapie
\item Je sais rédiger et justifier une démarche argumentée (introduction problématisée, thèse-antithèse-synthèse) sur un sujet de dissertation
\item Je sais mobiliser des références d'auteurs (Descartes, Freud, Nietzsche, Sartre) pour étayer une conclusion personnelle
\end{itemize}
\end{objectifs}

\begin{prerequis}
\begin{itemize}
\item Notion de sujet et de personne (classe de Première : la personne et l'identité personnelle)
\item Distinction entre connaissance et opinion (leçon sur la vérité)
\item Notion de liberté et de responsabilité (vie en société, éducation civique)
\item Capacité à problématiser un sujet et à construire un plan dialectique (méthode de la dissertation)
\item Notions de raison et d'expérience (leçon sur la connaissance)
\end{itemize}
\end{prerequis}

\newpage

\coursec{Problématique et définitions}

À Douala, un jeune commerçant du marché Congo jure qu'il a « oublié » de rendre la monnaie à un client pressé ; ses amis, eux, murmurent qu'il savait très bien ce qu'il faisait et qu'il a fait semblant. Au même moment, une élève de Terminale au lycée technique de New Bell fait chaque nuit le même rêve : elle arrive en retard à l'épreuve de philosophie, sa copie reste blanche, elle entend le jury rire. Elle se réveille en sueur, sans comprendre pourquoi ce scénario se répète. Sommes-nous vraiment maîtres de nos pensées, de nos paroles et de nos actes ? Que vaut le témoignage de notre propre conscience, et que cache ce que Freud appelle l'inconscient ? Cette leçon interroge la nature, la valeur et les limites de la conscience, puis évalue la portée et les limites de la psychanalyse.

\courssub{Étymologie et notion de conscience}

Le mot \cle{conscience} vient du latin \emph{conscientia}, lui-même formé sur \emph{cum-scire} : « savoir avec ». Cette étymologie révèle d'emblée la structure fondamentale : la conscience est une \textbf{connaissance réflexive}, un « savoir avec soi-même ». Elle n'est pas une simple réception passive du monde (perception), mais la présence du sujet à lui-même en tant qu'il perçoit, pense, veut ou agit.

\begin{definition}[Conscience]
La conscience est la \textbf{connaissance immédiate et réfléchie} que le sujet a de lui-même, de ses états intérieurs (pensées, sentiments, désirs) et de ses actes. Elle implique une dualité : le sujet qui connaît et l'objet connu sont identiques (le soi se saisissant soi-même).
\end{definition}

La philosophie distingue traditionnellement \textbf{trois sens} du mot conscience, qu'il est impératif de ne pas confondre dans une dissertation ou un commentaire.

\begin{popfigure}
\centering
\begin{tikzpicture}[
    node distance=0.8cm and 1.2cm,
    box/.style={rectangle, draw=popblue, thick, rounded corners=4pt, fill=popblueL, text width=4.2cm, align=center, font=\small},
    title/.style={rectangle, draw=popdark, thick, rounded corners=4pt, fill=popgold, text width=13cm, align=center, font=\bfseries\small},
    arrow/.style={->, >=Stealth, popmuted, thick}
]
% Title
\node[title] (titre) {Les trois sens du mot « conscience » avec exemples camerounais};
% Column headers
\node[box, anchor=south west] (h1) at ($(titre.south west)+(0,-0.5)$) {\textbf{Conscience psychologique}\\\textbf{(spontanée)}};
\node[box, anchor=south] (h2) at ($(titre.south)+(0,-0.5)$) {\textbf{Conscience réfléchie}\\\textbf{(moralepour certains auteurs)}};
\node[box, anchor=south east] (h3) at ($(titre.south east)+(0,-0.5)$) {\textbf{Conscience morale}\\\textbf{(jugement de valeur)}};
% Content rows
\node[box, anchor=north west] (c11) at ($(h1.south west)+(0,-0.3)$) {Présence immédiate à soi et au monde. « Je vois que je vois ». Pas de jugement, pure transparence.};
\node[box, anchor=north] (c21) at ($(h2.south)+(0,-0.3)$) {Retour du sujet sur lui-même. « Je sais que je sais ». Analyse, distance, objectivation de ses propres états.};
\node[box, anchor=north east] (c31) at ($(h3.south east)+(0,-0.3)$) {Faculté de juger le bien et le mal. Voix intérieure (\"voix de la conscience\"). Obligation morale.};
% Examples
\node[box, fill=popgreenL, draw=popgreen, anchor=north west] (c12) at ($(c11.south west)+(0,-0.2)$) {\textit{Exemple :} Un motard à Yaoundé freine brutalement devant un piéton. Il \textbf{vit} l'action sans se dire « je freine ».};
\node[box, fill=poporangeL, draw=poporange, anchor=north] (c22) at ($(c21.south)+(0,-0.2)$) {\textit{Exemple :} Le même motard, rentré chez lui, se dit : « J'ai eu peur, j'ai freiné trop fort, j'aurais pu glisser ». Il \textbf{analyse} son acte.};
\node[box, fill=poppinkL, draw=poppink, anchor=north east] (c32) at ($(c31.south east)+(0,-0.2)$) {\textit{Exemple :} Il se demande : « Ai-je bien fait de ne pas m'arrêter pour l'aider ? » Sa \textbf{conscience morale} le juge (culpabilité/fierté).};
% Arrows to show link
\draw[arrow] (c11.south) -- (c12.north);
\draw[arrow] (c21.south) -- (c22.north);
\draw[arrow] (c31.south) -- (c32.north);
\draw[arrow, dashed] (c12.east) -- ++(0.5,0) |- (c22.west) node[midway, above, sloped, font=\tiny, popmuted] {prise de distance};
\draw[arrow, dashed] (c22.east) -- ++(0.5,0) |- (c32.west) node[midway, above, sloped, font=\tiny, popmuted] {évaluation éthique};
\end{tikzpicture}
\legende{Tableau comparatif des trois sens de la conscience. La conscience psychologique est le fondement ; la réfléchie en est l'approfondissement intellectuel ; la morale en est l'exigence normative.}
\end{popfigure}

\begin{exemplebox}[Exemples concrets des trois sens au Cameroun]
\begin{itemize}
    \item \textbf{Conscience psychologique (spontanée) :} Un élève en classe sent la faim (sensation), entend le bruit de la pluie (perception), a une idée soudaine pour son devoir (pensée). Il \emph{est} ces états sans s'en faire un objet d'étude.
    \item \textbf{Conscience réfléchie :} Le même élève, après l'examen, se dit : « J'ai mal géré mon temps, j'ai paniqué sur le sujet 2, j'aurais dû commencer par le commentaire ». Il objectivité son vécu passé.
    \item \textbf{Conscience morale :} Face à une occasion de tricher (le surveillant tourne le dos), une voix intérieure dit : « C'est mal, tu perds ta dignité ». C'est le \cle{devoir moral} (Kant) ou le \cle{surmoi} (Freud) qui parle.
\end{itemize}
\end{exemplebox}

\begin{savaistu}[Le piège du double sens]
En français, le mot « conscience » désigne \textbf{à la fois} la conscience psychologique (connaissance de soi) et la conscience morale (jugement du bien et du mal). L'anglais distingue \textit{consciousness} (psychologique) et \textit{conscience} (morale). Au Bac, \textbf{ne pas préciser de quel sens on parle} est une cause majeure de hors-sujet. Exemple : « La conscience est-elle une connaissance ? » $\to$ Si on répond seulement par le devoir moral, on rate le sujet qui porte sur la structure épistémologique (psychologique/réfléchie).
\end{savaistu}

\coursec{Nature de la conscience}

\courssub{La conscience comme intentionnalité (Brentano, Husserl)}

Au-delà de la simple réflexivité, la conscience est toujours \textbf{conscience \emph{de} quelque chose}. C'est la thèse de l'\cle{intentionnalité} (Franz Brentano, repris par Husserl) : tout acte psychique vise un objet (réel ou imaginaire). Je ne « pense » pas en vide ; je pense \emph{à} cet arbre, \emph{à} mon examen, \emph{à} l'injustice subie. Cette structure intentionnelle distingue radicalement le mental du physique.

\begin{definition}[Intentionnalité]
L'intentionnalité est la propriété fondamentale de la conscience d'être toujours \textbf{tournée vers un objet}, de « viser » quelque chose qui la dépasse. Il n'y a pas de conscience « pure » sans contenu intentionnel.
\end{definition}

\courssub{Conscience spontanée vs conscience réfléchie : la genèse du « Moi »}

La conscience spontanée (vécue) est pré-réflexive : je vis ma colère. La conscience réfléchie opère un \textbf{redoublement} : je me sais en colère. C'est ce second niveau qui constitue le \cle{Moi} comme objet de connaissance pour moi-même. Mais attention : pour Sartre, la conscience réfléchie est déjà une \textbf{déformation} ; elle fige le flux vivant de la conscience spontanée en un « Moi » substantiel.

\begin{aretenir}[Nature de la conscience : trois caractères essentiels]
\begin{itemize}
    \item \textbf{Transparence (thèse cartésienne) :} Rien n'est dans la conscience que la conscience n'en ait conscience. Pas d'inconscient possible pour Descartes.
    \item \textbf{Intentionnalité :} Structure « sujet $\to$ objet ». La conscience n'est pas un contenant, mais un acte de visée.
    \item \textbf{Unité synthétique :} La conscience unifie la multiplicité des impressions en une expérience cohérente (Kant : « Je pense » doit pouvoir accompagner toutes mes représentations).
\end{itemize}
\end{aretenir}

\coursec{Valeur de la conscience}

\courssub{Valeur épistémologique : le fondement de la certitude (Descartes)}

Avec le \emph{Cogito} (\emph{Je pense, donc je suis}), la conscience devient le \textbf{premier principe} indubitable de la philosophie. Même si je doute de tout (monde, corps, Dieu), je ne peux douter que je doute, donc que je pense, donc que je suis. La conscience est \textbf{lumière naturelle} : elle se garantit elle-même.

\begin{propriete}[Évidence cartésienne]
La conscience offre une \textbf{certitude absolue} (subjective) sur l'existence du sujet pensant. C'est le point d'Archimède pour reconstruire le savoir (Dieu, monde, sciences).
\end{propriete}

\courssub{Valeur morale : fondement de la responsabilité et de la dignité (Kant)}

La conscience morale (le \emph{Gewissen}) est le tribunal intérieur. Pour Kant, elle est la \textbf{raison pratique} en moi : elle me fait sentir le devoir catégorique. Sans conscience, pas de liberté, pas d'imputation, pas de dignité humaine. L'homme est une fin en soi \emph{parce qu'il a une conscience morale}.

\begin{exemplebox}[Responsabilité pénale au Cameroun]
Le Code pénal camerounais (art. 64) prévoit l'irresponsabilité pour trouble psychique abolissant le discernement ou le contrôle des actes. La \textbf{conscience} (discernement + volonté) est la condition juridique de la culpabilité. Un somnambule qui blesse quelqu'un n'est pas « coupable » : sa conscience était absente.
\end{exemplebox}

\courssub{Valeur existentielle : l'authenticité (Sartre, Heidegger)}

La conscience (comme \emph{néantisation}, pouvoir de dire « non ») est ce qui permet à l'homme de se \textbf{déprendre} de la facticité (son passé, son milieu, son corps) pour se projeter vers l'avenir. \textbf{Valeur = Lucidité}. Fuir sa conscience (mauvaise foi), c'est se mentir à soi-même pour échapper à l'angoisse de la liberté.

\coursec{Limites de la conscience}

\courssub{L'illusion de la transparence (Nietzsche, Freud)}

« Nous ne sommes pas maîtres en notre demeure. » La conscience ne voit que la surface. Nietzsche (\emph{Généalogie de la morale}) montre que nos valeurs « nobles » masquent souvent du ressentiment. Freud systématise : le refoulement fait de la conscience un \textbf{écran} qui cache les vrais motifs (pulsions sexuelles, agressives).

\courssub{L'inconscient cognitif et les automatismes (Psychologie contemporaine)}

Bien au-delà de Freud, les sciences cognitives montrent que la majeure partie du traitement de l'information est \textbf{inconsciente} (perception, mémoire procédurale, prise de décision intuitive, priming). La conscience arrive \emph{après} le traitement (expériences de Libet : potentiel de readiness $\approx$ 550 ms avant l'intention consciente). La conscience est un « usager » tardif, pas le « PDG » du cerveau.

\courssub{Déterminismes sociaux et idéologiques (Marx, Bourdieu, Althusser)}

Ma conscience de moi-même et du monde est structurée par ma \textbf{position de classe}, mon habitus, les appareils idéologiques d'État (école, médias, famille). « Ce n'est pas la conscience des hommes qui détermine leur existence, c'est leur existence sociale qui détermine leur conscience » (Marx). La conscience est souvent \textbf{fausse conscience} (idéologie).

\begin{attention}[Limites $\neq$ Inutilité]
Dire que la conscience a des limites ne signifie pas qu'elle est \emph{illusoire} ou \emph{inutile}. Elle reste le lieu de la vérité subjective, de la délibération morale et de la construction du sens. La philosophie critique (Kant, Ricoeur) vise à \textbf{critiquer} ses prétentions à l'absolu, non à la nier.
\end{attention}

\coursec{L'inconscient freudien}

\courssub{Rupture épistémologique : de l'obscur à l'Inconscient dynamique}

Avant Freud : l'inconscient = \textbf{non-conscient} (absence de lumière : sommeil, distraction, automatismes physiologiques). C'est un \textbf{négatif}.
Avec Freud (1900, \emph{L'Interprétation des rêves}) : l'inconscient = \textbf{système psychique à part entière} (l'\emph{Inconscient} \emph{Ics}), avec ses lois (processus primaire : condensation, déplacement, intemporalité, absence de contradiction, non-négation), son énergie (pulsions) et sa dynamique (refoulement). C'est un \textbf{positif} qui \emph{agit}.

\begin{definition}[Inconscient (sens freudien strict)]
L'inconscient est l'\textbf{ensemble des contenus psychiques} (désirs refoulés, souvenirs traumatisants, pulsions sexuelles ou agressives, représentations) qui sont \textbf{inaccessibles à la conscience} par un travail actif de refoulement (\emph{Verdrängung}), mais qui \textbf{déterminent nos conduites}, nos symptômes, nos rêves et nos actes manqués.
\end{definition}

\courssub{La première topique (1900) : Conscience / Préconscient / Inconscient}

\begin{popfigure}
\centering
\begin{tikzpicture}[
    scale=1.1,
    every node/.style={font=\small},
    water/.style={fill=popblue!30, draw=popblue, thick},
    ice/.style={fill=popblue!10, draw=popblue, thick},
    ground/.style={fill=popmuted!30, draw=popmuted, thick},
    label/.style={font=\bfseries\small, text=popdark},
    arrow/.style={->, >=Stealth, poporange, thick, shorten >=2pt, shorten <=2pt}
]
% Sea level
\draw[popblue, very thick] (-4,0) -- (4,0) node[midway, above=2pt, label] {Ligne de flottaison = \textbf{Préconscient (Pcs)}};
% Conscious (above water)
\fill[ice] (-3,0) -- (3,0) -- (2.5,2.5) -- (-2.5,2.5) -- cycle;
\node[label, align=center] at (0, 1.5) {\textbf{CONSCIENCE (Cs)}\\(Clarté, logique secondaire,\\contrôle volontaire, réalité)};
\draw[arrow] (0, 1.5) -- (0, 0.3);
% Preconscious (water line zone)
\node[label, align=center, text=popblue] at (0, -0.5) {\textbf{PRÉCONSCIENT (Pcs)}\\(Disponible, souvenirs accessibles\\par effort, lien Cs/ICS)};
% Unconscious (under water)
\fill[ground] (-4,-4) rectangle (4,0);
\fill[poppurple!20, draw=poppurple, thick] (-3.5,0) -- (3.5,0) -- (3,-3.5) -- (-3,-3.5) -- cycle;
\node[label, align=center, text=poppurple] at (0, -2) {\textbf{INCONSCIENT (Ics)}\\(Refoulé, pulsions, processus primaire,\\rêves, actes manqués, intemporel)};
\draw[arrow, poppurple] (0, -2) -- (0, -0.3) node[midway, right=2pt, font=\tiny, popdark] {Pression / Retour du refoulé};
% Iceberg tip label
\node[align=center, font=\tiny, popmuted] at (3.5, 1.5) {Partie émergée\\$\sim$ 10\%};
\node[align=center, font=\tiny, popmuted] at (3.5, -2) {Partie immergée\\$\sim$ 90\%};
% Examples floating up
\node[draw=poporange, fill=poporangeL, rounded corners, font=\tiny, anchor=east, align=center] at (-4.2, 1) {Acte manqué :\\« Je te hais »\\au lieu de « Je t'aime »};
\node[draw=popgreen, fill=popgreenL, rounded corners, font=\tiny, anchor=east, align=center] at (-4.2, -1){Rêve :\\Examen raté\\= désir d'éviter l'angoisse};
\node[draw=poppink, fill=poppinkL, rounded corners, font=\tiny, anchor=east, align=center] at (-4.2, -3){Symptôme :\\Phobie, TOC,\\Lapsus};
\end{tikzpicture}
\legende{Le modèle topique freudien (1900) : l'iceberg. La conscience est la partie émergée. Le préconscient (ligne de flottaison) contient ce qui peut devenir conscient. L'inconscient (immergé) est le moteur principal, régi par le principe de plaisir et le refoulement.}
\end{popfigure}

\courssub{La seconde topique (1923) : Moi / Ça / Surmoi}

Freud affine le modèle pour rendre compte du conflit structural.

\begin{center}
\begin{tabularx}{\linewidth}{|l|X|X|X|}\hline
\rowcolor{popblueL}
\textbf{Instance} & \textbf{Origine / Fonction} & \textbf{Principes / Logique} & \textbf{Rapport à la conscience} \\ \hline
\textbf{Ça (Es/Id)} & Réservoir des pulsions (inné). Hérité, archaïque. & \textbf{Plaisir} : satisfaction immédiate, hallucinatoire. Pas de temps, pas de non, pas de contradiction. & \textbf{Entièrement inconscient}. \\ \hline
\textbf{Moi (Ego)} & Différenciation du Ça sous l'influence du monde extérieur. & \textbf{Réalité} : test de réalité, remise à plus tard, synthèse, défense. & \textbf{Partiellement conscient} (siège de la conscience), largement préconscient/inconscient (défenses). \\ \hline
\textbf{Surmoi (Über-Ich)} & Intériorisation des interdits parentaux/sociaux (Complexe d'Œdipe). Héritier du Ça. & \textbf{Perfection / Idéal} : juge, censure, culpabilité, idéal du moi. & \textbf{Largement inconscient} (conscience morale inconsciente), partiellement conscient. \\ \hline
\end{tabularx}
\end{center}

\begin{aretenir}[Conflit structural]
La vie psychique est un \textbf{théâtre de conflits} : le Moi doit satisfaire le Ça (pulsions), le Surmoi (exigences morales) et la Réalité extérieure. L'\textbf{anxiété} est le signal de danger pour le Moi. Les \textbf{mécanismes de défense} (refoulement, déplacement, rationalisation, sublimation, déni, projection...) sont des opérations \emph{inconscientes} du Moi pour réduire l'anxiété.
\end{aretenir}

\begin{attention}[Inconscient \emph{vs} Non-conscient : erreur fatale au Bac]
\begin{itemize}
    \item \textbf{Non-conscient} (ou \emph{inconscient} au sens large, non freudien) : simple absence momentanée de conscience. Exemples : le sommeil sans rêve, l'anesthésie, la distraction, les automatismes (marcher sans y penser), la digestion. C'est un \textbf{manque} de conscience.
    \item \textbf{Inconscient freudien} : une \textbf{présence active} refoulée. Ce n'est pas « ne pas savoir », c'est « savoir sans le savoir » \emph{contre sa volonté}. Le refoulement est une \textbf{défense active}. L'inconscient \emph{veut} s'exprimer (symptômes, rêves, lapsus).
    \item \textbf{Piège :} Dire « L'inconscient, c'est quand on dort » $\to$ \textbf{0 point} sur la notion. L'inconscient freudien travaille \emph{pendant} le sommeil (rêves) et \emph{pendant} la veille (actes manqués).
\end{itemize}
\end{attention}

\coursec{Valeur et limites de la psychanalyse}

\courssub{Valeur thérapeutique et herméneutique}

\begin{itemize}
    \item \textbf{Soin par la parole :} La « talking cure » libère l'énergie bloquée dans le symptôme. Lever le refoulement $\to$ récupération de l'histoire, élargissement du Moi.
    \item \textbf{Herméneutique de la suspicion (Ricoeur) :} La psychanalyse apprend à \textbf{déchiffrer} les manifestations trompeuses (rêves, lapsus, symptômes) comme des \emph{textes} codés. Elle restaure le sujet comme \emph{auteur} de son symptôme.
    \item \textbf{Éthique du désir :} « Là où c'était le Ça, il faut que le Moi advienne. » Non pas supprimer le désir, mais l'assumer, le reconnaître, en devenir responsable. Passage de la soumission à l'appropriation.
\end{itemize}

\courssub{Valeur anthropologique et culturelle}

La psychanalyse a transformé la culture occidentale : art (surréalisme), littérature, éducation (prise en compte de l'affectif, de l'enfance), droit (responsabilité atténuée), compréhension des phénomènes collectifs (fascisme, racisme comme pulsions de mort projetées, psychologie des foules).

\courssub{Limites épistémologiques et scientifiques (Popper, Grünbaum, neurosciences)}

\begin{itemize}
    \item \textbf{Non-falsifiabilité (Popper) :} La psychanalyse explique \emph{tout} et son contraire (ex. : un homme qui sauve un enfant = altruisme/sublimation ; qui ne le sauve pas = lâcheté/refoulement). Une théorie qui ne peut être réfutée n'est pas scientifique.
    \item \textbf{Preuve clinique circulaire (Grünbaum) :} Les « confirmations » viennent de la cure elle-même (transfert, suggestion de l'analyste), pas de tests indépendants.
    \item \textbf{Défi des neurosciences :} Les mécanismes freudiens (refoulement, Ça) trouvent-ils des corrélats neuronaux ? (Recherche actuelle en neuro-psychanalyse : circuits de la peur, inhibition préfrontale, mais pas de « Ça » localisé).
\end{itemize}

\courssub{Limites cliniques, éthiques et culturelles}

\begin{itemize}
    \item \textbf{Durée et coût :} Cure interminable (années), réservée à une élite. Inadaptée aux urgences psychiatriques (psychoses aiguës).
    \item \textbf{Phallocentrisme et normativité :} Critiques féministes (Simone de Beauvoir, Luce Irigaray, Judith Butler) : l'anatomie n'est pas le destin ; l'Œdipe universalise un modèle familial bourgeois viennois du XIX\textsuperscript{e} siècle.
    \item \textbf{Transfert/Contre-transfert :} Relation de pouvoir asymétrique. Risque d'emprise, de fausses souvenirs (affaire des « faux souvenirs » induits).
    \item \textbf{Indications :} Efficace sur les névroses (angoisse, phobies, TOC, dépressions névrotiques). Contre-indiquée ou insuffisante seule dans les psychoses (schizophrénie), troubles bipolaires, TCA sévères, addictions (nécessitent approche multidisciplinaire : médication, TCC, social).
\end{itemize}

\begin{aretenir}[Bilan critique : ni science exacte, ni charlatanisme]
La psychanalyse est une \textbf{herméneutique de la subjectivité} (science de l'interprétation du sens), pas une science de la nature (lois causales universelles). Sa valeur est \textbf{pragmatique} (ça soulage, ça donne du sens) et \textbf{éthique} (autonomie, lucidité), mais sa scientificité au sens poppérien reste discutée. Au Bac : distinguer \emph{validité thérapeutique} et \emph{statut épistémologique}.
\end{aretenir}

\coursec{Synthèse, méthode et entraînement au Bac}

\courssub{Tableau de synthèse : Conscience / Inconscient}

\begin{center}
\begin{tabularx}{\linewidth}{|l|X|X|}\hline
\rowcolor{popblueL}
\textbf{Dimension} & \textbf{Conscience (Idéal cartésien / Kantien)} & \textbf{Inconscient (Réalité freudienne / Critique)} \\ \hline
\textbf{Nature} & Transparence, unité, liberté, raison. & Opacité, division, déterminisme pulsionnel, processus primaire. \\ \hline
\textbf{Connaissance de soi} & Immédiate, évidente (\emph{Cogito}). & Médiatisée, difficile, infinie (analyse interminable). \\ \hline
\textbf{Liberté / Responsabilité} & Pleine imputabilité (sujet d'imputation). & Responsabilité \emph{atténuée} / \emph{partagée} (sujet divisé). \\ \hline
\textbf{Valeur} & Fondement de la morale, du droit, de la science. & Accès à la vérité profonde, soin, humanisation. \\ \hline
\textbf{Limites} & Illusion de maîtrise, ignore les racines. & Non-falsifiable, risque de déterminisme fataliste, coût, culture-dépendance. \\ \hline
\end{tabularx}
\end{center}

\courssub{Problématique centrale (rappel)}

\begin{aretenir}[Les deux pôles de la problématique]
\begin{itemize}
    \item \textbf{Pôle 1 (Épistémologie) :} \emph{La conscience suffit-elle à me connaître moi-même ?} $\to$ Non, elle est aveuglée par l'inconscient ; mais l'inconscient n'est connaissable \emph{que} par le détour de la conscience (analyse).
    \item \textbf{Pôle 2 (Éthique/Juridique) :} \emph{L'inconscient remet-il en cause ma liberté et ma responsabilité ?} $\to$ Il la \emph{relativise} (déterminisme), mais l'éthique psychanalytique vise à \emph{reconquérir} une liberté lucide (assumer son désir).
\end{itemize}
\end{aretenir}

\begin{methode}[Problématiser un sujet de dissertation (Bac Philo Tech)]
\textbf{Étape 1 : Analyser les termes.} Définir rigoureusement chaque mot du sujet (ex. : « conscience » $\to$ psychologique ? réfléchie ? morale ?).
\textbf{Étape 2 : Identifier la tension.} Repérer l'opposition implicite : \emph{Thèse A} (sens commun, évidence première) \emph{vs} \emph{Thèse B} (analyse philosophique, contre-intuitive). Ex. : « Je me connais par conscience » \emph{vs} « Je me méconnais par inconscient ».
\textbf{Étape 3 : Formuler la question problème.} Transformer l'opposition en question ouverte, nuancée, qui appelle une discussion en plusieurs temps. \emph{Mal formulée :} « Y a-t-il un inconscient ? » (Fermée, fait). \emph{Bien formulée :} « La découverte de l'inconscient détruit-elle l'idéal cartésien de maîtrise de soi ? » ou « La conscience est-elle un savoir fiable sur soi-même ? »
\textbf{Étape 4 : Annoncer le plan (en introduction).} « Nous verrons d'abord que la conscience offre une prise immédiate sur soi (I), mais que l'inconscient en révèle les illusions (II), pour enfin demander si une connaissance de soi intégrale est possible (III). »
\end{methode}

\begin{exoresolu}[Analyse d'un sujet type Probatoire/Bac Tech]
\textbf{Sujet :} \emph{« L'inconscient est-il un obstacle à la connaissance de soi ? »}

\textbf{1. Analyse des termes :}
\begin{itemize}
    \item \textbf{Inconscient :} Au sens freudien (système dynamique refoulé), pas au sens de « non-conscient » (sommeil).
    \item \textbf{Obstacle :} Ce qui barre la route, empêche l'accès. Mais aussi : ce qui oblige à détourner, à creuser (obstacle épistémologique \emph{à la Bachelard} : source de savoir nouveau).
    \item \textbf{Connaissance de soi :} Accès à ses motifs profonds, vérité du sujet (idéal socratique « Connais-toi toi-même »).
\end{itemize}

\textbf{2. Tension / Débat :}
\begin{itemize}
    \item \textbf{Oui (Thèse) :} L'inconscient \emph{cache} la vérité. Le refoulement masque les vrais motifs. Le commerçant \emph{croit} avoir oublié ; sa conscience lui ment. L'inconscient est opacité, mensonge, résistance. Il est un obstacle à la transparence cartésienne.
    \item \textbf{Non / Nuance (Antithèse) :} L'inconscient \emph{révèle} la vérité profonde. Sans lui, je reste à la surface (consciente) de moi-même. Le rêve, le lapsus, le symptôme sont des \emph{messages codés}. La psychanalyse \emph{utilise} l'inconscient comme voie d'accès (royale) à la connaissance de soi. L'obstacle devient \emph{instrument}.
    \item \textbf{Synthèse :} L'inconscient est un obstacle à la \emph{connaissance immédiate, spontanée, narcissique} (le « moi » conscient), mais la \emph{condition nécessaire} d'une \emph{connaissance médiatisée, vraie, éthique} (le « moi » analysé). Connaître soi, c'est \emph{intégrer} son inconscient, pas le supprimer.
\end{itemize}

\textbf{3. Plan détaillé (3 parties, 2--3 sous-parties chacune) :}
\begin{enumerate}
    \item \textbf{L'inconscient comme obstacle à la transparence du Cogito.}
    \begin{itemize}
        \item A. La conscience spontanée se croit transparente (Descartes) : « Je pense, donc je suis » $\to$ certitude immédiate.
        \item B. La découverte freudienne brise cette évidence : le refoulement, la censure, le moi « pas maître en sa demeure ». Ex. : l'oubli du commerçant (refoulement de la cupidité).
    \end{itemize}
    \item \textbf{L'inconscient comme voie d'accès à une vérité plus profonde.}
    \begin{itemize}
        \item A. Les formations de l'inconscient (rêves, actes manqués, symptômes) sont des \emph{compromis} qui disent la vérité déguisée. Ex. : le rêve d'examen raté = désir d'éviter l'angoisse de l'épreuve.
        \item B. La méthode psychanalytique (association libre, levée de la censure) transforme l'obstacle en outil. « Là où c'était le Ça, il faut que le Moi advienne ».
    \end{itemize}
    \item \textbf{Vers une connaissance de soi assumée : ni toute conscience, ni tout inconscient.}
    \begin{itemize}
        \item A. Limites de la psychanalyse : l'inconscient reste partiellement insaisissable (noyau impossible à symboliser), transfert, interminabilité de la cure.
        \item B. La connaissance de soi comme tâche éthique permanente : vigilance (conscience morale) + lucidité sur ses déterminismes (inconscient). Responsabilité = répondre de ses actes \emph{en connaissant} leurs racines obscures.
    \end{itemize}
\end{enumerate}

\textbf{Introduction type (schéma) :}
\begin{quote}
\emph{« Connais-toi toi-même »} ordonnait l'oracle de Delphes. Descartes fonda la philosophie moderne sur la certitude du \emph{Cogito} : la conscience est transparente à elle-même. Mais au tournant du XX\textsuperscript{e} siècle, Freud découvre un continent caché : l'inconscient. Le commerçant qui « oublie » de rendre la monnaie, l'élève qui rêve d'échec : leurs consciences les trompent. L'inconscient est-il dès lors un mur infranchissable dressé contre la connaissance de soi, ou bien le détour obligé pour atteindre une vérité plus authentique ? Nous montrerons que l'inconscient ruine l'illusion d'une transparence immédiate (I), qu'il devient paradoxalement la voie royale vers une connaissance médiatisée (II), pour conclure que la vraie connaissance de soi n'est ni dans le tout-conscient ni dans le tout-inconscient, mais dans l'assomption lucide de sa division (III).
\end{quote}
\end{exoresolu}

\begin{exercice}[Entraînement Bac - Sujets probables]
\begin{enumerate}
    \item « La conscience est-elle un savoir fiable sur soi-même ? » (Dissertation)
    \item « L'inconscient nous excuse-t-il ? » (Dissertation)
    \item Commentaire de texte : Extrait de Freud, \emph{Introduction à la psychanalyse} (Conférence XVII : « Signification du symptôme ») ou Sartre, \emph{L'Être et le Néant} (Introduction : « La poursuite de l'être » / Mauvaise foi).
    \item « Peut-on être responsable sans le savoir ? » (Dissertation - articulation conscience morale / inconscient).
\end{enumerate}
\end{exercice}

\begin{corrige}[Exercice n°1 - Pistes de dissertation]
\textbf{Sujet :} \emph{« La conscience est-elle un savoir fiable sur soi-même ? »}

\textbf{Analyse :} « Savoir fiable » = certitude, exhaustivité, objectivité. « Sur soi-même » = motifs, désirs, identité profonde.
\textbf{Plan :}
\begin{enumerate}
    \item \textbf{La conscience comme savoir premier et certain (Thèse).} A. Évidence cartésienne (Cogito) : je ne peux douter que je pense. B. Conscience morale (Kant) : accès immédiat au devoir. C. Phénoménologie (Husserl) : description rigoureuse du vécu.
    \item \textbf{Les illusions de la conscience : l'inconscient comme limite structurelle (Antithèse).} A. Freud : refoulement, Ça, Surmoi inconscient $\to$ la conscience ment (rationalisation). B. Marx/Bourdieu : conscience fausse, habitus, déterminismes sociaux. C. Neurosciences : décision inconsciente avant intention consciente (Libet).
    \item \textbf{Vers une fiabilité critique : la conscience comme tâche (Synthèse).} A. La fiabilité n'est pas donnée, elle se \emph{construit} (examen de conscience, psychanalyse, vigilance épistémologique). B. Distinction : fiabilité sur l'\emph{acte} (je sais que j'agis) vs fiabilité sur le \emph{motif} (je sais \emph{pourquoi} j'agis). C. La conscience fiable est une \emph{conscience critique}, lucide sur ses propres zones d'ombre.
\end{enumerate}
\end{corrige}

\coursec{Nature de la conscience}

Dans la section précédente, nous avons posé la problématique de la leçon et défini la conscience comme la faculté de se représenter soi-même et le monde. Mais cette définition reste générale : elle ne dit pas encore \emph{ce qu'est} la conscience dans sa nature profonde. La conscience est-elle une simple présence au monde, semblable à celle de l'animal ? Est-elle un pouvoir de retour sur soi ? Est-elle une voix qui juge ? Pour répondre, il faut distinguer plusieurs \emph{degrés} ou \emph{formes} de la conscience, puis interroger ses caractères essentiels : l'intentionnalité, le langage, la temporalité. Telle est la tâche de cette section.

\courssub{La conscience spontanée : présence immédiate à soi et au monde}

Commençons par l'expérience la plus simple. Vous marchez sous le soleil de Maroua en pleine saison sèche : vous sentez la chaleur sur votre peau, vous entendez les motos, vous voyez la poussière de la route. Vous n'y \emph{pensez} pas : vous le \emph{vivez}. Cette présence immédiate, directe, sans distance, est la \cle{conscience spontanée}.

\begin{definition}[Conscience spontanée]
La \cle{conscience spontanée} (ou conscience immédiate) est la présence directe et non réfléchie du sujet à lui-même et au monde. Elle accompagne toutes nos perceptions, sensations et actions sans que nous ayons besoin d'y penser. Elle ne se formule pas en paroles : elle se vit.
\end{definition}

Ses caractères sont les suivants :
\begin{itemize}
\item \textbf{L'immédiateté} : aucun intermédiaire, aucun raisonnement ne sépare le sujet de son vécu. La douleur d'une brûlure est sentie avant d'être pensée.
\item \textbf{L'adhésion au monde} : dans la conscience spontanée, je suis « absorbé » par ce que je fais. Le chauffeur de taxi-brousse expérimenté conduit sans se répéter chaque geste.
\item \textbf{L'universalité} : elle est partagée par tous les hommes, et sans doute par les animaux supérieurs, qui perçoivent et réagissent à leur milieu.
\end{itemize}

\begin{exemplebox}[Le footballeur de quartier]
Un joueur, lors d'un match à Yaoundé, dribble deux adversaires et marque. Interrogé après le match, il dit : « Je n'ai pas réfléchi, c'est venu tout seul. » Pendant l'action, il était pleinement conscient — présent à son corps, au ballon, à l'espace — mais sans aucune réflexion sur ce qu'il faisait. C'est la conscience spontanée dans sa pureté.
\end{exemplebox}

La conscience spontanée est donc le \emph{sol} de toute vie consciente. Mais si l'homme n'avait qu'elle, il ne se distinguerait guère de l'animal. Ce qui fait la grandeur et la spécificité humaine, c'est le passage à un second degré.

\courssub{La conscience réfléchie : le retour de l'esprit sur lui-même}

\begin{definition}[Conscience réfléchie]
La \cle{conscience réfléchie} est le retour de l'esprit sur lui-même : au lieu d'être seulement tournée vers le monde, la conscience se prend elle-même pour objet. Je ne sens plus seulement la chaleur : je \emph{sais que je la sens}. Je ne suis plus seulement en colère : je \emph{m'observe} en colère.
\end{definition}

Le mot « réfléchie » vient du latin \emph{reflectere}, « plier en arrière » : comme un miroir renvoie la lumière, l'esprit se renvoie à lui-même sa propre image. Cette capacité de dédoublement — être à la fois celui qui agit et celui qui s'observe agir — est proprement humaine.

C'est ici que s'insère la démarche du philosophe français René \textbf{Descartes} (1596-1650). Dans le \emph{Discours de la méthode} (1637), Descartes entreprend de douter de tout : les sens peuvent nous tromper, les raisonnements peuvent être faux, et peut-être même que je rêve en ce moment. Mais au bout de ce doute radical, une vérité résiste absolument.

\begin{definition}[Le cogito cartésien]
Le \cle{cogito} est la formule de Descartes : « \emph{Je pense, donc je suis} » (en latin \emph{cogito ergo sum}), énoncée dans le \emph{Discours de la méthode} (1637). Même si je doute de tout, même si un malin génie me trompe sur tout, je ne peux pas douter que je doute, donc que je pense ; et pour penser, il faut que je sois, que j'existe. La conscience de soi est ainsi la \textbf{première vérité indubitable}, le fondement sur lequel tout le reste du savoir peut être reconstruit.
\end{definition}

Le raisonnement peut se reconstruire ainsi, pas à pas :
\begin{enumerate}
\item Je doute de tout ce qui peut être mis en doute (sensations, savoirs, existence du monde).
\item Mais douter est encore une manière de \emph{penser}.
\item Or, pour penser, il faut exister : on ne peut pas penser sans être.
\item Donc, au moment même où je pense, il est absolument certain que j'existe.
\end{enumerate}

Descartes en tire une conclusion sur la \emph{nature} de la conscience : je suis une « chose qui pense », une \cle{substance pensante} (\emph{res cogitans}). L'essence de l'esprit humain, c'est la pensée, c'est-à-dire la conscience. La conscience n'est donc pas un accident ou une propriété secondaire : elle est ce que je suis de plus certain et de plus fondamental.

\begin{attention}[Erreur fréquente au Bac]
N'attribuez jamais le cogito à Platon, à Socrate ou à Sartre. « Je pense, donc je suis » est la formule de \textbf{Descartes} (\emph{Discours de la méthode}, 1637). Platon est le philosophe de l'Idée et du mythe de la caverne ; Sartre est un existentialiste du XX\textsuperscript{e} siècle. Confondre les auteurs est une faute grave qui coûte des points à l'examen.
\end{attention}

\courssub{La conscience morale : la voix intérieure qui juge nos actes}

Il existe une troisième forme de conscience, qui n'est ni simple présence au monde ni connaissance de soi, mais \emph{jugement} sur soi.

\begin{definition}[Conscience morale]
La \cle{conscience morale} est la faculté par laquelle le sujet juge ses propres actes et intentions au regard du bien et du mal. Elle se manifeste comme une « voix intérieure » qui approuve (satisfaction, fierté) ou condamne (remords, honte, culpabilité) ce que nous faisons ou projetons de faire.
\end{definition}

\begin{exemplebox}[Le remords après la tricherie]
Un élève de Terminale technique recopie la réponse de son voisin pendant une épreuve du Baccalauréat blanc. Il obtient une bonne note, mais le soir, il ne dort pas bien : une voix intérieure lui répète qu'il a mal agi, qu'il s'est trompé lui-même avant de tromper le correcteur. Ce \emph{remords} n'a été commandé par personne : c'est sa propre conscience morale qui juge son acte. Rousseau voyait dans cette voix intérieure un instinct divin du bien ; Kant y voyait la raison pratique qui se donne à elle-même sa loi.
\end{exemplebox}

\begin{exemplebox}[Un indice d'universalité : sondage simulé en classe]
Dans un sondage simulé réalisé dans une classe de Terminale technique de 40 élèves, 32 élèves, soit 80~\%, déclarent avoir déjà éprouvé un remords après une mauvaise action (mensonge, vol, tricherie, parole blessante). Ce résultat, bien que fictif et local, illustre l'\textbf{universalité de la conscience morale} : quelles que soient la culture, la religion ou la région d'origine, tout homme semble porter en lui un juge intérieur. Reste à savoir — question ouverte pour la dissertation — si le \emph{contenu} de ce juge (ce qu'il condamne) est le même partout ou s'il dépend de l'éducation.
\end{exemplebox}

\begin{propriete}[Conscience réfléchie, connaissance de soi et responsabilité]
La conscience réfléchie est le fondement de la connaissance de soi et de la responsabilité morale. En effet : (1) seul un être capable de retour sur soi peut se connaître, identifier ses motifs et ses défauts ; (2) seul un être capable de s'observer et de se juger peut être tenu pour responsable de ses actes, car la responsabilité suppose que le sujet pouvait s'examiner, comparer son acte à la règle et choisir. C'est pourquoi le fou ou l'enfant très jeune, dont la conscience réfléchie est absente ou défaillante, ne sont pas tenus pour pleinement responsables devant la loi. (Propriété admise, non démontrée ici.)
\end{propriete}

\begin{popfigure}
\begin{center}
\begin{tikzpicture}[node distance=1.6cm and 2.2cm,
  etape/.style={rectangle, rounded corners=6pt, draw=popblue, fill=popblueL, thick, text width=4.6cm, align=center, minimum height=1.6cm, font=\small},
  fleche/.style={-{Stealth[length=3mm]}, very thick, poporange}]
\node[etape, align=center] (sp){\textbf{Conscience spontanée}\\ Présence immédiate à soi et au monde\\ \emph{« Je sens la chaleur »}};
\node[etape, right=of sp, draw=poppurple, fill=poppurpleL, align=center] (ref){\textbf{Conscience réfléchie}\\ Retour de l'esprit sur lui-même\\ \emph{« Je sais que je sens »}};
\node[etape, right=of ref, draw=popgreen, fill=popgreenL, align=center] (mor){\textbf{Conscience morale}\\ Jugement de mes actes\\ \emph{« Je juge que j'ai mal agi »}};
\draw[fleche] (sp) -- (ref) node[midway, above, font=\small\itshape, popdark]{approfondissement};
\draw[fleche] (ref) -- (mor) node[midway, above, font=\small\itshape, popdark]{jugement};
\end{tikzpicture}
\end{center}
\legende{Les trois degrés de la conscience : de la simple présence au monde (spontanée) au retour sur soi (réfléchie), puis au jugement de soi (morale). Chaque degré suppose le précédent et ajoute une profondeur nouvelle.}
\end{popfigure}

\begin{savaistu}[L'examen de conscience, une sagesse universelle]
La pratique du retour sur soi n'est pas propre à la philosophie occidentale. Dans plusieurs traditions de sagesse africaine, notamment chez les anciens du Cameroun, la journée se clôt par un \emph{examen de conscience du soir} : l'ancien repasse ses paroles et ses actes, demande pardon à qui il a offensé et se corrige pour le lendemain. Cette pratique, proche de la conscience réfléchie décrite par les philosophes, montre que le souci de se connaître et de se juger est un héritage de toute l'humanité, et non le monopole d'une culture.
\end{savaistu}

\courssub{La conscience comme intentionnalité : toute conscience est conscience de quelque chose}

Nous avons décrit les degrés de la conscience. Mais quelle est sa \emph{structure} fondamentale ? Le philosophe allemand Edmund \textbf{Husserl} (1859-1938), reprenant une idée de son maître Brentano, répond : la conscience est essentiellement \cle{intentionnalité}.

\begin{definition}[Intentionnalité]
L'\cle{intentionnalité} est la propriété essentielle de la conscience d'être toujours \emph{conscience de quelque chose} : elle vise, elle « tend vers » un objet. Je ne perçois pas « rien » : je perçois un arbre ; je n'aime pas « en général » : j'aime quelqu'un ; je ne me souviens pas dans le vide : je me souviens de mon enfance à Bafoussam. La conscience n'est pas une boîte fermée contenant des idées : elle est une flèche, un mouvement vers le monde.
\end{definition}

Cette thèse a une conséquence capitale : elle renverse l'image d'une conscience enfermée en elle-même, spectatrice passive d'un film intérieur. Si toute conscience est conscience de quelque chose, alors la conscience est \emph{ouverture au monde}, relation vivante entre un sujet et un objet. Sartre, héritier de Husserl, dira que la conscience est « dehors », dans le monde, comme le vent : elle n'a pas d'intérieur.

\begin{popfigure}
\begin{center}
\begin{tikzpicture}[node distance=4.5cm,
  pole/.style={rectangle, rounded corners=6pt, thick, text width=3.2cm, align=center, minimum height=1.4cm, font=\small},
  fleche/.style={-{Stealth[length=3.5mm]}, very thick, poporange}]
\node[pole, draw=popblue, fill=popblueL, align=center] (s){\textbf{Sujet}\\ (conscience)};
\node[pole, draw=poppurple, fill=poppurpleL, right=of s, align=center] (o){\textbf{Objet}\\ (chose visée)};
\draw[fleche] (s) -- (o) node[midway, above, font=\small\itshape, popdark]{vise (intention)};
\node[below=0.9cm of s, font=\small, align=center, text width=4.2cm] {percevoir, imaginer,\\ aimer, se souvenir...};
\node[below=0.9cm of o, font=\small, align=center, text width=4.2cm] {un arbre, un ami,\\ un souvenir, un projet...};
\end{tikzpicture}
\end{center}
\legende{Schéma de l'intentionnalité (Husserl) : la conscience n'existe pas seule ; elle est toujours un mouvement du sujet vers un objet visé. Sans objet visé, pas d'acte de conscience ; sans sujet, pas de visée.}
\end{popfigure}

\courssub{Conscience et langage : dire « je »}

Comment savons-nous qu'un être possède la conscience de soi ? Un indice puissant est le langage. Prononcer le mot « \cle{je} » n'est pas anodin : celui qui dit « je » se désigne lui-même comme sujet, se distingue des autres (« tu », « il ») et du monde. Dire « je » suppose donc déjà une conscience réfléchie de soi.

Réciproquement, le langage \emph{structure} la conscience : grâce aux mots, je peux nommer mes émotions (« je suis triste », « j'ai honte »), raconter ma vie, formuler mes intentions. Un enfant qui apprend à parler apprend en même temps à se penser. La conscience et le langage avancent donc ensemble : le « je » parlé est le signe et l'instrument du « je » pensant. C'est pourquoi Descartes voyait dans l'usage raisonnable de la parole le signe distinctif de l'homme par rapport à l'animal et à la machine.

\courssub{Conscience et temporalité : mémoire, anticipation, projet}

Dernier caractère essentiel : la conscience n'est pas enfermée dans l'instant présent. Elle \emph{relie le passé et l'avenir}.

\begin{itemize}
\item Par la \textbf{mémoire}, la conscience retient le passé : je me souviens de mes années au collège, de mes réussites et de mes échecs. Cette mémoire constitue mon \emph{identité} : je suis le même que celui qui, enfant, allait à l'école primaire de mon village.
\item Par l'\textbf{anticipation}, la conscience se projette dans l'avenir : j'imagine mon examen de fin d'année, je prévois, je crains ou j'espère.
\item Par le \textbf{projet}, la conscience organise le présent en vue de l'avenir : l'élève de Terminale technique qui révise chaque soir relie son effort présent à son but futur — obtenir son Baccalauréat et devenir technicien.
\end{itemize}

\begin{exemplebox}[Le projet professionnel comme œuvre de la conscience]
Une élève de Terminale F (génie civil) décide de devenir conductrice de travaux. Ce projet suppose : la mémoire (elle se souvient de son goût pour les chantiers de son quartier), l'anticipation (elle s'imagine dans cinq ans sur un chantier à Douala) et la volonté consciente (elle organise ses révisions en conséquence). Aucun animal ne construit ainsi son existence : la temporalité vécue est une marque propre de la conscience humaine.
\end{exemplebox}

Ainsi, la conscience tisse la trame d'une vie : sans elle, nous ne serions qu'une suite d'instants sans lien, sans identité ni destin. Saint Augustin, dans les \emph{Confessions}, voyait déjà dans cette « distension » de l'âme entre mémoire et attente le propre du temps humain.

\courssub{Exercice résolu et pièges}

\begin{exoresolu}[Identifier les formes de la conscience]
\textbf{Énoncé.} Pour chacune des situations suivantes, indiquez la forme de conscience dominante (spontanée, réfléchie, morale) et justifiez en une phrase.
\begin{enumerate}
\item Un mécanicien de Nkongsamba répare un moteur « les yeux fermés », sans se poser de questions.
\item Le soir, une élève écrit dans son journal intime : « Aujourd'hui, j'ai compris pourquoi je me mets si vite en colère. »
\item Après avoir menti à son père, un lycéen ne se sent pas en paix et décide d'avouer.
\item En regardant le mont Cameroun, un randonneur se dit : « Je suis en train de voir le mont Cameroun. »
\end{enumerate}
\tcblower
\textbf{Solution détaillée.}
\begin{enumerate}
\item \textbf{Conscience spontanée} : le mécanicien est pleinement présent à sa tâche, sans retour réflexif sur lui-même ; son savoir-faire s'exerce immédiatement.
\item \textbf{Conscience réfléchie} : l'élève se prend elle-même pour objet d'analyse ; elle cherche à se connaître, ce qui est le propre du retour de l'esprit sur lui-même.
\item \textbf{Conscience morale} : le trouble intérieur du lycéen est un jugement porté sur son acte au regard du bien ; le désir d'avouer manifeste la voix intérieure qui condamne le mensonge.
\item \textbf{Conscience réfléchie} : le randonneur ne se contente pas de voir ; il \emph{sait qu'il voit}. Le « je suis en train de... » est la marque du dédoublement réflexif (et illustre au passage le cogito : dans cet acte de pensée, il se saisit comme existant).
\end{enumerate}
\end{exoresolu}

\begin{attention}[Pièges de rédaction]
\begin{itemize}
\item Ne confondez pas \emph{conscience spontanée} et \emph{inconscience} : le footballeur absorbé par son match est pleinement conscient, mais non réfléchi.
\item Ne réduisez pas la conscience morale à la peur du gendarme ou de la sanction : le remords existe même quand personne n'a vu la faute. C'est précisément ce qui prouve qu'il vient de l'intérieur.
\item Ne présentez pas le cogito comme un raisonnement compliqué : c'est une \emph{évidence saisie dans l'acte même de penser}, et non une déduction abstraite.
\end{itemize}
\end{attention}

\begin{aretenir}[L'essentiel de la section]
\begin{itemize}
\item La \cle{conscience spontanée} est la présence immédiate à soi et au monde ; la \cle{conscience réfléchie} est le retour de l'esprit sur lui-même ; la \cle{conscience morale} est le jugement intérieur de nos actes.
\item Le \cle{cogito} de Descartes (« Je pense, donc je suis », 1637) établit la conscience de soi comme première vérité indubitable et l'esprit comme substance pensante.
\item Selon Husserl, la conscience est \cle{intentionnalité} : elle est toujours conscience \emph{de} quelque chose.
\item Dire « je » suppose la conscience de soi ; la conscience relie passé (mémoire) et avenir (anticipation, projet), fondant ainsi l'identité personnelle.
\end{itemize}
\end{aretenir}

Cette analyse de la nature de la conscience pose déjà une question : si la conscience est présence, réflexion, jugement et projet, quelle est sa \emph{valeur} pour la connaissance, l'action et la liberté ? C'est ce que nous examinerons dans la section suivante.

\coursec{Valeur de la conscience}

\noindent Après avoir analysé la \emph{nature} de la conscience — sa structure intentionnelle, sa réflexivité et son unité —, nous devons maintenant interroger sa \emph{valeur}. Pourquoi la conscience est-elle précieuse ? Que nous permet-elle de faire que nous ne pourrions pas faire sans elle ? La réponse s'articule autour de quatre piliers : la connaissance de soi, la liberté, la responsabilité (morale et juridique) et la vie morale et sociale. Sans conscience, l'homme ne serait qu'un automate biologique, incapable de vérité, de choix, de devoir ni de parole donnée.

\courssub{La conscience comme connaissance de soi : « Connais-toi toi-même »}

\noindent La première valeur de la conscience est \emph{épistémologique} : elle rend le sujet présent à lui-même. L'injonction inscrite au fronton du temple de Delphes, « \emph{Gnôthi seauton} » (Connais-toi toi-même), reprise par \textbf{Socrate}, ne saurait s'accomplir sans la conscience réflexive. C'est parce que je suis conscient de mes pensées, de mes désirs, de mes motifs, que je peux les examiner, les critiquer, les ordonner.

\begin{definition}[Connaissance de soi par la conscience]
La connaissance de soi est l'acte par lequel le sujet, par la conscience réflexive, saisit ses propres états mentaux (pensées, sentiments, volontés) comme \emph{siens}, les distingue les uns des autres et en discerne les motifs et les enjeux.
\end{definition}

\noindent Contrairement à la perception extérieure qui me donne le monde, la conscience réflexive me donne \emph{moi-même} comme objet. Mais attention : cette transparence n'est pas totale. Comme nous le verrons dans l'étude de l'inconscient, la psychanalyse freudienne vient limiter cette prétention à une pleine lucidité. Pour l'instant, retenons que la conscience est la \emph{condition de possibilité} de toute sagesse, de tout examen de vie, de tout progrès moral.

\begin{exemplebox}[L'examen de conscience socratique]
Un élève de Terminale technique hésite entre deux filières post-bac. Au lieu de suivre l'avis de ses parents ou la mode, il pratique l'examen de conscience : il interroge ses goûts réels, ses aptitudes, ses valeurs profondes. Ce dialogue intérieur, rendu possible par la conscience réflexive, lui permet un choix \emph{authentique} et non subi.
\end{exemplebox}

\begin{popfigure}
\begin{tikzpicture}[
  mindmap,
  grow cyclic,
  every node/.style={concept, execute at begin node=\hskip0pt},
  concept color=popblue!80,
  level 1/.style={level distance=4.5cm, sibling angle=90},
  level 2/.style={level distance=3.2cm, sibling angle=45},
  text=white,
  font=\small
]
\node[concept, minimum size=2.2cm] (center) {Conscience}
  child[concept color=popgreen!70] {node[concept, minimum size=1.6cm, align=center]{Connaissance\\ de soi}
    child {node[concept, minimum size=1.1cm, align=center]{Transparence\\ réflexive}}
    child {node[concept, minimum size=1.1cm, align=center]{Examen\\ de vie}}
    child {node[concept, minimum size=1.1cm, align=center]{Authenticité\\ des choix}}
  }
  child[concept color=poporange!70] {node[concept, minimum size=1.6cm] {Liberté}
    child {node[concept, minimum size=1.1cm] {Délibération}}
    child {node[concept, minimum size=1.1cm, align=center]{Choix\\ éclairé}}
    child {node[concept, minimum size=1.1cm] {Autonomie}}
  }
  child[concept color=poppurple!70] {node[concept, minimum size=1.6cm] {Responsabilité}
    child {node[concept, minimum size=1.1cm, align=center]{Moralité\\ de l'acte}}
    child {node[concept, minimum size=1.1cm, align=center]{Imputabilité\\ juridique}}
    child {node[concept, minimum size=1.1cm, align=center]{Sanction\\ juste}}
  }
  child[concept color=poppink!70] {node[concept, minimum size=1.6cm, align=center]{Vie morale\\ et sociale}
    child {node[concept, minimum size=1.1cm, align=center]{Reconnaissance\\ d'autrui}}
    child {node[concept, minimum size=1.1cm, align=center]{Parole donnée\\ promesse}}
    child {node[concept, minimum size=1.1cm, align=center]{Contrat\\ social}}
  };
\end{tikzpicture}
\legende{Carte mentale des valeurs de la conscience : quatre branches principales issues de la conscience réflexive.}
\end{popfigure}

\courssub{La conscience comme fondement de la liberté}

\noindent Deuxième valeur, \emph{ontologique} celle-là : la conscience est la condition \emph{sine qua non} de la liberté. Être libre, ce n'est pas faire « ce qu'on veut » au gré des pulsions ; c'est pouvoir \emph{délibérer}, c'est-à-dire se représenter des motifs, anticiper des conséquences, et \emph{choisir} en connaissance de cause. Sans conscience, pas de représentation des possibles, pas de mise à distance des déterminismes, pas de décision.

\begin{propriete}[La conscience est la condition de la liberté]
Choisir suppose de se représenter les motifs et les conséquences de l'action. La conscience réflexive permet cette mise à distance : je ne \emph{subis} pas mon désir, je le \emph{contemple} et je décide de le suivre ou de le réprimer. \textbf{Justification} : l'animal agit par instinct (déterminisme naturel) ; l'homme conscient peut dire « non » à son penchant immédiat au nom d'une valeur supérieure (ex. : jeûne, courage, persévérance dans l'étude).
\end{propriete}

\noindent Cette thèse est au cœur du \emph{cogito} cartésien et de la philosophie des Lumières : la liberté naît de la lumière de la raison consciente. Mais attention : la conscience \emph{n'est pas} la liberté ; elle en est la \emph{condition}. On peut être conscient et non libre (contrainte extérieure, aliénation). En revanche, on ne peut être libre \emph{sans} conscience.

\begin{attention}[Piège fréquent : confondre conscience et liberté]
Ne pas écrire : « La conscience \emph{est} la liberté. » Écrire : « La conscience \emph{rend possible} la liberté en permettant la délibération. » La distinction est cruciale en dissertation : la conscience est \emph{nécessaire} mais non \emph{suffisante} à la liberté.
\end{attention}

\begin{exoresolu}[Conscience et liberté dans un choix technique]
\textbf{Énoncé :} Un technicien en maintenance découvre une anomalie sur une machine. Il peut la signaler (arrêt de production, perte financière) ou la taire (risque d'accident). Montrez en quoi la conscience est le lieu de sa liberté.
\tcblower
\textbf{Solution détaillée :}
\begin{enumerate}
\item \textbf{Représentation des options :} La conscience permet au technicien de se représenter deux scénarios : (A) signaler, avec un coût immédiat mais la sécurité assurée ; (B) taire, avec une production continue mais un risque différé.
\item \textbf{Mise à distance de la pression :} La contrainte hiérarchique (« ne pas arrêter la chaîne ») s'impose à lui comme un \emph{fait}. Mais la conscience le rend \emph{présent à lui-même} face à ce fait : il peut l'accepter ou le refuser.
\item \textbf{Délibération morale :} Il évalue les conséquences à la lumière de ses valeurs (responsabilité professionnelle, sécurité des collègues). Ce dialogue intérieur est l'exercice même de la liberté.
\item \textbf{Choix :} Il décide de signaler. Ce choix n'est pas déterminé par l'instinct de conservation (qui pousserait au silence) ni par la contrainte seule. Il émane d'une \emph{volonté consciente}.
\end{enumerate}
\emph{Conclusion :} Sans conscience réflexive, le technicien serait un rouage de plus, déterminé par la contrainte la plus forte. La conscience fait de lui un \emph{agent moral libre}.
\end{exoresolu}

\courssub{La conscience comme fondement de la responsabilité morale et juridique}

\noindent Troisième valeur, \emph{normative} : la conscience fonde la responsabilité. On ne demande des comptes qu'à qui a agi \emph{en connaissance de cause}. L'imputabilité — morale comme juridique — suppose que l'auteur de l'acte en a compris le sens, la portée et la valeur. C'est ce que le droit appelle le \emph{discernement}.

\begin{definition}[Responsabilité]
Obligation de répondre de ses actes accomplis en connaissance de cause. La responsabilité suppose trois conditions : (1) l'acte est volontaire (liberté), (2) l'agent en comprend la signification (conscience et discernement), (3) l'agent peut en être tenu pour auteur (imputabilité).
\end{definition}

\noindent En droit camerounais comme dans la plupart des systèmes juridiques modernes, la distinction entre acte volontaire et acte involontaire est fondatrice. Le Code pénal camerounais (loi n°2016/007 du 12 juillet 2016) prévoit l'\emph{irresponsabilité pénale} lorsque l'auteur, au moment de l'acte, était atteint d'un trouble mental ayant aboli sa capacité de \emph{discernement} ou le contrôle de ses actes. La loi reconnaît ainsi explicitement le lien entre conscience (discernement) et responsabilité.

\begin{exemplebox}[Application : le Code pénal camerounais et l'abolition du discernement]
\textbf{Principe légal :} N'est pas pénalement responsable la personne qui, au moment des faits, était atteinte d'un trouble mental ayant aboli son discernement ou le contrôle de ses actes.

\textbf{Exemple concret :} Un conducteur, victime d'une crise de délire aigu (bouffée délirante), renverse un piéton. L'expertise psychiatrique établit que, au moment des faits, il ne comprenait pas la nature de son acte. Le tribunal prononce l'irresponsabilité pénale (mais peut ordonner des mesures de sûreté, comme une hospitalisation d'office). \emph{La conscience (discernement) manquait, donc pas de responsabilité pénale.}
\end{exemplebox}

\begin{popfigure}
\begin{tikzpicture}[
  node distance=1.2cm and 1.5cm,
  decision/.style={diamond, draw=popdark, thick, aspect=2, text width=2.8cm, align=center, fill=poporange!15, font=\small},
  process/.style={rectangle, draw=popblue, thick, rounded corners, text width=3cm, align=center, fill=popblue!10, font=\small},
  startstop/.style={rectangle, draw=popgreen, thick, rounded corners, text width=3cm, align=center, fill=popgreen!15, font=\small},
  arrow/.style={->, thick, popdark},
  lbl/.style={font=\footnotesize, popmuted}
]
\node[startstop, align=center] (start){Acte commis\\ (fait matériel)};
\node[process, below of=start, align=center] (qualif){Qualification juridique\\ de l'acte};
\node[decision, below of=qualif, align=center] (discern){L'auteur avait-il\\ son discernement ?};
\node[process, left of=discern, xshift=-3.5cm, align=center] (non){Irresponsabilité\\ pénale};
\node[process, right of=discern, xshift=3.5cm, align=center] (oui){Responsabilité\\ pénale engagée};
\node[decision, below of=oui, align=center] (volonte){L'acte était-il\\ volontaire ?};
\node[process, left of=volonte, xshift=-3.5cm, align=center] (invol){Responsabilité\\ atténuée\\ (involontaire)};
\node[process, right of=volonte, xshift=3.5cm, align=center] (vol){Responsabilité\\ pleine\\ (volontaire)};
\node[startstop, below of=volonte, yshift=-1.2cm, align=center] (fin){Sanction proportionnée\\ (peine, réparation)};

\draw[arrow] (start) -- (qualif);
\draw[arrow] (qualif) -- (discern);
\draw[arrow] (discern) -- node[lbl, above] {Non} (non);
\draw[arrow] (discern) -- node[lbl, above] {Oui} (oui);
\draw[arrow] (oui) -- (volonte);
\draw[arrow] (volonte) -- node[lbl, above] {Non} (invol);
\draw[arrow] (volonte) -- node[lbl, above] {Oui} (vol);
\draw[arrow] (non) |- (fin);
\draw[arrow] (invol) |- (fin);
\draw[arrow] (vol) -- (fin);
\end{tikzpicture}
\legende{Organigramme de décision judiciaire : l'appréciation de la responsabilité passe par le discernement (conscience) puis par la volontarité de l'acte.}
\end{popfigure}

\begin{aretenir}[Conscience et responsabilité : le lien nécessaire]
\begin{itemize}
\item \textbf{Moralement :} Je ne suis blâmable ou louable que pour ce que j'ai fait \emph{sciemment}. L'erreur invincible excuse (ex. : je donne un produit toxique en pensant que c'est un remède).
\item \textbf{Juridiquement :} Le droit positif fait du discernement (conscience de la portée de l'acte) une condition de l'imputabilité. Trouble mental, contrainte physique, force irrésistible entraînent l'abolition ou l'altération du discernement, donc l'irresponsabilité ou une responsabilité atténuée.
\end{itemize}
\end{aretenir}

\courssub{La conscience morale comme guide de l'action : le devoir (Kant)}

\noindent Quatrième valeur, \emph{éthique} : la conscience n'est pas seulement lumière (connaissance) ; elle est \emph{voix}. La \emph{conscience morale} est cette instance intérieure qui me dit « tu dois » ou « tu ne dois pas ». Elle distingue le bien du mal, non par calcul d'intérêt, mais par exigence inconditionnée.

\noindent Pour \textbf{Emmanuel Kant}, la conscience morale est le \emph{témoin intérieur} de la loi morale. Elle est la voix de la raison pratique en moi : ce qui élève l'homme au-dessus de la nature, c'est qu'il peut agir \emph{par devoir}, c'est-à-dire par respect pour la loi morale qu'il se donne à lui-même en tant qu'être raisonnable.

\begin{savaistu}[Kant : la conscience morale, dignité de l'homme]
Kant affirme que la conscience morale est ce qui élève l'homme au-dessus des animaux. L'animal suit ses instincts ; l'homme conscient du devoir peut leur \emph{résister} au nom d'un « impératif catégorique » (ex. : « Ne mens pas », même si le mensonge m'avantage). Cette autonomie morale — se donner à soi-même sa propre loi — est le fondement de la \emph{dignité} humaine : la personne a une valeur infinie, elle n'a pas de prix.
\end{savaistu}

\begin{exemplebox}[Cas de conscience : le technicien et le rapport de conformité]
Un technicien chargé du contrôle qualité doit certifier la conformité d'une installation. Le directeur le presse de signer « pour ne pas retarder le chantier ». Le technicien sait qu'un défaut mineur mais réel subsiste. Sa conscience morale lui intime : « Ne certifie pas ce qui n'est pas conforme. » Il résiste à la pression hiérarchique. C'est l'expérience du \emph{devoir} comme exigence intérieure librement acceptée.
\end{exemplebox}

\courssub{La conscience comme condition de la vie sociale}

\noindent Cinquième valeur, \emph{politique et anthropologique} : la conscience rend possible la vie en société. La reconnaissance d'autrui comme \emph{sujet conscient} (et non comme objet) est le fondement de l'éthique de la rencontre. La \emph{parole donnée}, la \emph{promesse}, le \emph{contrat} supposent que je suis conscient de m'engager, que l'autre est conscient de recevoir mon engagement, et que nous sommes tous deux conscients de l'obligation qui en naît.

\begin{propriete}[La conscience fonde l'intersubjectivité et le contrat]
Sans conscience réflexive, pas de \emph{sujet de droit} capable de s'obliger par sa parole. La promesse (« Je m'engage à... ») suppose : (1) je me représente l'avenir, (2) je me représente l'autre comme destinataire de ma parole, (3) je me représente l'obligation morale qui me liera. Tout cela est acte de conscience.
\end{propriete}

\noindent C'est ce que l'on appelle la \emph{personnalité juridique} : la personne de droit est une construction qui repose sur la personne physique consciente et responsable. La société camerounaise, comme toute société, repose sur ce pacte implicite : « Je te reconnais comme conscience libre et responsable, tu me reconnais de même. »

\begin{attention}[Ne pas confondre conscience morale et conformisme social]
La conscience morale n'est pas la simple intériorisation des normes sociales. Elle garde une \emph{distance critique} : je peux juger qu'une loi ou une coutume est injuste (ex. : pratiques discriminatoires) et m'y opposer \emph{au nom de ma conscience}. C'est la distinction entre le \emph{légal} et le \emph{légitime}.
\end{attention}

\courssub{Valeur théorique : la conscience, instrument de toute connaissance}

\noindent Sixième valeur, \emph{gnoséologique} : la conscience est l'instrument universel de la connaissance. Toute science, toute vérité, tout savoir présuppose un sujet \emph{conscient} qui connaît. Comme l'écrit \textbf{Descartes} dans les \emph{Méditations} : « Je pense, donc je suis » — le \emph{cogito} est la première certitude parce qu'il est l'acte même de la conscience se saisissant elle-même en train de penser.

\begin{aretenir}[Conscience et vérité]
\begin{itemize}
\item \textbf{Science :} Le savant observe, formule des hypothèses, déduit, vérifie. Toutes ces opérations sont des actes de conscience intentionnelle (conscience \emph{de} quelque chose).
\item \textbf{Technique :} L'ingénieur conçoit, modélise, teste. La conscience projective (anticipation, modélisation) est au cœur de l'acte technique.
\item \textbf{Vérité :} La vérité n'est pas « dans le monde » comme un objet ; elle est l'adéquation entre un \emph{jugement conscient} et la réalité. Pas de jugement sans conscience.
\end{itemize}
\end{aretenir}

\courssub{Synthèse : sans conscience, ni liberté, ni morale, ni connaissance ne sont possibles}

\noindent Résumons. La conscience n'est pas un « luxe » ni un épiphénomène ; elle est le \emph{socle} de l'humanité de l'homme :

\begin{center}
\begin{tabularx}{\linewidth}{|l|X|X|}\hline
\rowcolor{popblueL}\textbf{Domaine} & \textbf{Apport de la conscience} & \textbf{Si la conscience est absente} \\ \hline
Connaissance de soi & Transparence, examen de vie, authenticité & Opacité, hétéronomie, aliénation \\ \hline
Liberté & Délibération, choix, autonomie & Déterminisme pur, automatisme \\ \hline
Responsabilité & Imputabilité, justice, sanction juste & Irresponsabilité générale, arbitraire \\ \hline
Morale & Devoir, dignité, résistance au mal & Conformisme, loi du plus fort \\ \hline
Vie sociale & Promesse, contrat, reconnaissance d'autrui & Guerre de tous contre tous (Hobbes) \\ \hline
Connaissance et science & Vérité, méthode, objectivité & Opinion, préjugé, hasard \\ \hline
\end{tabularx}
\end{center}

\begin{methode}[Valoriser une notion dans une dissertation : la méthode du « rendre possible »]
Pour montrer la \emph{valeur} d'une notion (ici : la conscience) dans une dissertation :
\begin{enumerate}
\item \textbf{Identifiez les grandes capacités humaines} qu'elle rend possibles : liberté, responsabilité, morale, science, vie sociale, dignité.
\item \textbf{Montrez le lien de conditionnalité} : « Sans conscience, pas de X » (condition nécessaire).
\item \textbf{Nuancez} : la conscience est nécessaire mais non toujours suffisante (ex. : conscience + contrainte extérieure = pas de liberté effective).
\item \textbf{Illustrez par des exemples concrets} (juridiques, techniques, quotidiens) et des références d'auteurs (Socrate, Descartes, Kant ; Freud pour la limite).
\item \textbf{Concluez sur le statut anthropologique} : la conscience est ce qui fait de l'homme un \emph{sujet} et non un objet.
\end{enumerate}
\end{methode}

\begin{exercice}[Entraînement Bac — Sujet type]
\textbf{Sujet :} « La conscience est-elle la condition de la responsabilité ? »

\textbf{Consignes :} Problématisez, proposez un plan détaillé (3 parties, 2 sous-parties chacune), rédigez l'introduction complète.
\end{exercice}

\begin{corrige}[Exercice n°1 — Éléments de corrigé]
\textbf{Problématique :} La responsabilité suppose-t-elle \emph{uniquement} la conscience (discernement) ou aussi la liberté effective ? La conscience est-elle condition \emph{nécessaire et suffisante} ?

\textbf{Plan :}
\begin{enumerate}
\item \textbf{La conscience comme condition nécessaire de la responsabilité} : (a) discernement moral et examen de ses actes ; (b) imputabilité juridique — l'abolition du discernement entraîne l'irresponsabilité pénale.
\item \textbf{Mais la conscience ne suffit pas : la liberté effective est requise} : (a) contrainte physique et force irrésistible ; (b) déterminismes sociaux et psychiques — la responsabilité peut être atténuée.
\item \textbf{La responsabilité comme articulation indissociable de la conscience et de la liberté} : (a) le sujet de droit et l'autonomie kantienne ; (b) la sanction comme reconnaissance de la dignité de la personne.
\end{enumerate}

\textbf{Introduction type :} « "Chacun est responsable de ses actes" : cette évidence morale et juridique repose sur un présupposé anthropologique : l'homme est un être conscient et libre. Mais la conscience suffit-elle à fonder la responsabilité ? Si le droit excuse celui dont le discernement est aboli, il exige aussi, pour la responsabilité pleine, que l'acte soit volontaire. Nous montrerons que la conscience est la condition \emph{sine qua non} de la responsabilité, mais qu'elle doit s'articuler à la liberté effective pour fonder une imputabilité juste. »
\end{corrige}

\coursec{Limites de la conscience}

Nous venons de le voir : la conscience est une \cle{valeur} précieuse. Elle nous rend libres, responsables, capables de réflexion et de choix. Descartes en faisait le fondement indubitable de toute connaissance : « Je pense, donc je suis ». Mais faut-il pour autant lui accorder une confiance absolue ? La conscience dit-elle toujours la vérité ? Saisit-elle tout ce qui se passe en nous et autour de nous ? L'observation la plus simple montre que non : la conscience peut se tromper, se mentir à elle-même, et ignorer les forces qui la déterminent. C'est ce que nous allons établir dans cette section : la conscience, si précieuse soit-elle, possède des \cle{limites} qu'il faut connaître pour ne pas devenir dupe de soi-même.

\courssub{Les illusions des sens : la conscience peut se tromper}

Commençons par une expérience que chacun peut faire. Plongez un bâton à moitié dans un seau d'eau : il paraît \emph{brisé}, coudé à la surface de l'eau. Retirez-le : il est parfaitement droit. Votre conscience, qui s'appuie sur la vue, vous a trompé. De même, la lune paraît plus grosse à l'horizon qu'au zénith, les rails du chemin de fer semblent se rejoindre au loin, et le voyageur dans le train arrêté croit avancer quand le train voisin démarre.

\begin{exemplebox}[L'illusion du bâton brisé]
Un pêcheur de la Wouri voit son harpon se « casser » à la surface de l'eau. Sa conscience perceptive lui affirme une chose fausse. La cause est physique : la lumière change de direction en passant de l'eau à l'air (réfraction). Mais la conscience, qui ne connaît pas ce mécanisme, juge d'après l'apparence. \textbf{Leçon} : la conscience perçoit les apparences, non les choses mêmes. C'est pourquoi Descartes lui-même, dans les \emph{Méditations}, recommande de se méfier des sens : « il est de la prudence de ne se fier jamais entièrement à ceux qui nous ont une fois trompés ».
\end{exemplebox}

Ces erreurs ne sont pas des accidents rares : elles sont \emph{normales} et \emph{constantes}. La conscience ne reçoit pas le monde tel qu'il est, mais une reconstruction que notre appareil sensoriel fabrique. Première limite donc : la conscience est \cle{faillible} dans la connaissance du monde extérieur.

\courssub{L'illusion de transparence : la mauvaise foi}

On pourrait croire que si la conscience se trompe sur le monde, elle se connaît au moins elle-même parfaitement. « Je sais ce que je pense, je sais pourquoi j'agis », disons-nous. C'est ce qu'on appelle l'\cle{illusion de transparence} : l'idée que la conscience serait transparente à elle-même, comme une vitre sans tache. Or l'expérience montre le contraire : nous nous mentons à nous-mêmes.

\begin{definition}[La mauvaise foi (Sartre)]
La \cle{mauvaise foi} est, selon Jean-Paul Sartre (\emph{L'Être et le Néant}, 1943), l'attitude par laquelle le sujet \textbf{se ment à lui-même} pour fuir sa liberté et sa responsabilité. Le sujet en mauvaise foi se cache la vérité qu'il connaît pourtant : il se fait passer pour ce qu'il n'est pas, ou il prétend n'être que ce que les circonstances ont fait de lui. La mauvaise foi est un mensonge où le menteur et le menti sont la même personne.
\end{definition}

Sartre donne l'exemple célèbre du garçon de café qui joue à être garçon de café : ses gestes sont trop rapides, trop précis, trop empressés. Il se prend pour une chose, un « être-de-café », comme s'il n'avait pas choisi ce métier et ne pouvait pas en changer. Autre exemple, plus proche de nous : l'élève qui échoue et déclare « de toute façon, je ne suis pas fait pour les études », alors qu'il n'a pas travaillé. Il se ment : il transforme un choix (ne pas travailler) en destin (être incapable), pour échapper au reproche et à l'effort.

\begin{attention}[Ne pas confondre mauvaise foi et mensonge ordinaire]
Dans le mensonge ordinaire, je trompe \emph{autrui} et je connais la vérité. Dans la mauvaise foi, je me trompe \emph{moi-même} : je suis à la fois celui qui cache la vérité et celui à qui on la cache. C'est précisément ce paradoxe qui montre que la conscience n'est pas une : une partie de moi sait, une autre partie refuse de savoir. La conscience n'est donc pas transparente à elle-même.
\end{attention}

\courssub{Les déterminismes sociaux : la conscience conditionnée par le milieu}

Troisième limite, plus profonde encore : ce que je crois penser librement est souvent le produit de mon \cle{milieu social}. Karl Marx l'a formulé dans une phrase célèbre de \emph{L'Idéologie allemande} (1846) :

\begin{aretenir}[La thèse de Marx]
« Ce n'est pas la conscience qui détermine la vie, mais la vie qui détermine la conscience. » Autrement dit : nos idées, nos valeurs, nos goûts, nos jugements ne naissent pas d'une réflexion pure et libre ; ils sont façonnés par nos conditions d'existence — notre famille, notre quartier, notre milieu professionnel, notre époque. La conscience croit être la cause de nos choix ; elle en est souvent l'effet.
\end{aretenir}

Prenons un exemple camerounais concret. Un jeune de Yaoundé veut devenir artiste musicien. Sa famille, son quartier, son entourage lui répètent : « fais médecine, fais le droit, c'est ça la réussite ». Il finit par choisir la médecine en disant : « c'est mon choix, j'ai réfléchi ». Mais ce choix « libre » a été orienté par la pression familiale, le prestige social du métier, la peur de décevoir. La conscience du jeune s'approprie après coup une décision que le milieu a en grande partie fabriquée. De même, la coutume, la langue, la religion du groupe s'installent en nous dès l'enfance, avant toute réflexion : nous les tenons pour évidentes alors qu'elles sont acquises.

\begin{exemplebox}[L'expérience d'Asch (1951) : le groupe fait taire le jugement]
Le psychologue Solomon Asch place un sujet dans un groupe de complices. On montre des lignes de longueurs visiblement différentes et on demande lesquelles sont égales. Les complices donnent à voix haute une réponse manifestement fausse. Résultat : environ \textbf{75 \% des sujets se conforment au moins une fois à l'avis faux du groupe}, contre leur propre témoignage visuel. Interrogés ensuite, beaucoup croient sincèrement avoir bien jugé. \textbf{Leçon} : la pression sociale modifie le jugement conscient, et la conscience ne s'en aperçoit même pas. Elle croit décider seule alors que le groupe décide en elle.
\end{exemplebox}

\begin{savaistu}[Le poisson ne voit pas l'eau]
Un proverbe dit : « le poisson ne voit pas l'eau ». Le poisson vit dans l'eau, l'eau le porte, le nourrit, le conditionne — mais justement parce qu'elle est partout, il ne la remarque jamais. C'est exactement la thèse de Marx : notre milieu social nous détermine \emph{sans que nous en ayons conscience}, précisément parce qu'il nous est trop familier. Le citadin ne « voit » plus la publicité qui oriente ses achats ; le jeune ne « voit » plus les codes de son quartier qui dictent sa façon de parler et de s'habiller. Le premier pas vers la liberté est de voir l'eau.
\end{savaistu}

\courssub{Les déterminismes biologiques et psychologiques}

Quatrième limite : la conscience dépend du \cle{corps}. Nous le constatons chaque jour sans en tirer la leçon.

\begin{itemize}
\item Après une nuit sans sommeil, le même élève devient irritable, pessimiste, incapable de se concentrer : ses jugements changent, alors que seule la \emph{fatigue} a changé.
\item La faim, la maladie, la fièvre modifient l'humeur et les décisions. Un malade juge le monde plus hostile qu'un bien-portant.
\item Les \cle{pulsions} — colère, peur, désir — s'imposent à la conscience et la font déraisonner : « je n'étais pas moi-même », dit-on après coup. Cette expression populaire est philosophiquement exacte : quelque chose en moi, qui n'est pas ma conscience, a parlé et agi à ma place.
\item Certains produits (alcool, drogues) altèrent radicalement la conscience, preuve qu'elle dépend d'états physiques du cerveau.
\end{itemize}

Spinoza l'avait déjà dit : les hommes se croient libres parce qu'ils ont conscience de leurs désirs, mais ils \emph{ignorent les causes} qui les déterminent à désirer. La conscience enregistre le résultat (« je veux ceci ») sans connaître le mécanisme (« mon corps, ma fatigue, mon passé me poussent à le vouloir »).

\courssub{Nietzsche : la conscience n'est que la surface de l'esprit}

Le philosophe Friedrich Nietzsche radicalise ce constat. Pour lui, la conscience n'est qu'une mince pellicule à la surface d'une vie psychique immense et profonde, faite d'instincts, de pulsions, de forces que nous ne contrôlons pas.

\begin{aretenir}[La thèse de Nietzsche]
Pour Nietzsche, la conscience est « la dernière et la plus tardive » des acquisitions de l'esprit : elle ne contient qu'une infime partie de ce qui se passe en nous. Pire : elle est un \emph{mensonge} sur nos véritables mobiles. Elle nous raconte de belles raisons (« j'agis par devoir, par générosité ») alors que les vraies causes sont des instincts (la peur, l'orgueil, la volonté de puissance). La conscience serait comme un interprète qui traduit mal et embellit ce que les profondeurs de l'esprit décident réellement.
\end{aretenir}

Exemple : un commerçant de Douala fait un don à une école et croit agir par pure générosité. Nietzsche demanderait : n'y a-t-il pas aussi le désir d'être vu, admiré, respecté ? La conscience donne une version flatteuse de l'acte ; les mobiles réels restent dans l'ombre. Cela ne signifie pas que le don est mauvais, mais que la conscience \emph{ne sait pas tout} de ce qui la fait agir.

\begin{popfigure}
\begin{center}
\begin{tikzpicture}[scale=0.95]
% Surface : conscience
\fill[popblueL] (-5.2,0) rectangle (5.2,1.2);
\draw[popblue, thick] (-5.2,0) rectangle (5.2,1.2);
\node[popdark, font=\bfseries] at (0,0.6) {CONSCIENCE (surface) : ce que je crois penser et vouloir};
% Profondeurs
\fill[poporangeL] (-5.2,-1.4) rectangle (5.2,0);
\draw[poporange, thick] (-5.2,-1.4) rectangle (5.2,0);
\node[popdark] at (0,-0.7) {Déterminismes sociaux : famille, quartier, coutume, groupe (Marx)};
\fill[popgreenL] (-5.2,-2.8) rectangle (5.2,-1.4);
\draw[popgreen, thick] (-5.2,-2.8) rectangle (5.2,-1.4);
\node[popdark] at (0,-2.1) {Déterminismes biologiques : corps, fatigue, maladie, pulsions (Spinoza)};
\fill[poppurpleL] (-5.2,-4.2) rectangle (5.2,-2.8);
\draw[poppurple, thick] (-5.2,-4.2) rectangle (5.2,-2.8);
\node[popdark] at (0,-3.5) {Vie psychique cachée : instincts, désirs refoulés (Nietzsche, Freud)};
% Flèches d'influence vers le haut
\draw[-{Stealth}, very thick, poppurple] (-3.5,-4.0) -- (-3.5,1.0);
\draw[-{Stealth}, very thick, popgreen] (0,-2.6) -- (0,1.0);
\draw[-{Stealth}, very thick, poporange] (3.5,-1.2) -- (3.5,1.0);
\node[popmuted, font=\small, anchor=west] at (5.4,0.6) {visible};
\node[popmuted, font=\small, anchor=west] at (5.4,-3.5) {invisible};
\draw[-{Stealth}, popmuted] (5.9,0.2) -- (5.9,-3.2);
\end{tikzpicture}
\end{center}
\legende{Schéma en couches de la vie psychique : la conscience n'est que la surface visible ; les déterminismes sociaux, biologiques et inconscients agissent en profondeur et remontent vers elle (flèches) sans qu'elle le sache.}
\end{popfigure}

\courssub{Les actes inexpliqués : indices d'une vie psychique cachée}

Dernière observation, décisive : il existe des actes que la conscience ne peut pas s'expliquer elle-même.

\begin{itemize}
\item Le \cle{lapsus} : je veux dire « monsieur le principal » et je dis « monsieur le prisonnier ». Qui a parlé ? Pas ma conscience, puisqu'elle voulait dire autre chose.
\item L'\cle{oubli} : j'oublie précisément le rendez-vous qui m'angoissait, le nom de la personne que je n'aime pas. L'oubli n'est pas toujours un accident : il semble obéir à une intention cachée.
\item Le \cle{rêve} : chaque nuit, des images, des histoires se produisent en moi sans que je les aie voulues ni comprises. D'où viennent-elles ?
\item L'\cle{acte manqué} : je casse « maladroitement » le cadeau d'une personne que je jalouse, je rate « par hasard » le bus qui me menait à un examen redouté.
\end{itemize}

Ces faits banals sont des \emph{indices} : ils suggèrent qu'il existe en nous une vie psychique qui échappe à la conscience et qui pourtant agit, parle, décide. La conscience n'est donc pas toute la vie mentale. C'est la porte ouverte à l'hypothèse de l'\cle{inconscient}, que nous étudierons avec Freud dans la section suivante.

\begin{exoresolu}[Analyse d'une situation concrète]
\textbf{Énoncé.} Amina, élève de Terminale à Bafoussam, affirme : « J'ai librement choisi la série F parce que j'aime la mécanique. » Or son père est mécanicien, ses frères sont en série F, et son quartier admire les techniciens. De plus, la veille de son choix, elle a fait un rêve étrange où elle réparait la moto de son père. Montrez, en vous appuyant sur deux limites de la conscience, que le choix d'Amina n'est peut-être pas aussi libre qu'elle le croit.
\tcblower
\textbf{Solution détaillée.}
\begin{enumerate}
\item \emph{Déterminisme social (Marx).} Le choix d'Amina est conditionné par son milieu : le métier du père, l'exemple des frères, le prestige des techniciens dans le quartier. Sa conscience s'approprie (« j'aime la mécanique ») une orientation que la vie familiale et sociale a fabriquée. Comme le poisson qui ne voit pas l'eau, Amina ne voit pas le milieu qui la détermine.
\item \emph{Vie psychique cachée (indice du rêve).} Le rêve de la moto montre que des désirs et des images liés au père travaillent Amina sans qu'elle les maîtrise. Sa conscience ignore ces mobiles profonds.
\item \emph{Conclusion nuancée.} On ne peut pas dire que le choix d'Amina ne vaut rien : sa conscience a pesé, comparé, décidé. Mais ce choix est \emph{conditionné} par des causes sociales et psychiques qu'elle ignore. La conscience est nécessaire à la décision, mais non suffisante pour l'expliquer entièrement.
\end{enumerate}
\end{exoresolu}

\begin{popfigure}
\begin{center}
\begin{tikzpicture}
\matrix[column sep=0.4cm]{
\node[fill=popblueL, draw=popblue, thick, rounded corners, minimum width=6.6cm, minimum height=0.9cm, font=\bfseries] {Ce que la conscience croit}; &
\node[fill=poppurpleL, draw=poppurple, thick, rounded corners, minimum width=6.6cm, minimum height=0.9cm, font=\bfseries] {Ce qui détermine réellement l'acte}; \\[0.25cm]
\node[draw=popline, rounded corners, minimum width=6.6cm, minimum height=1.5cm, align=center, text width=6.2cm] {« J'ai choisi médecine\\par passion » (jeune de Yaoundé)}; &
\node[draw=popline, rounded corners, minimum width=6.6cm, minimum height=1.5cm, align=center, text width=6.2cm] {Pression de la famille, prestige social du métier, peur de décevoir (Marx)}; \\[0.25cm]
\node[draw=popline, rounded corners, minimum width=6.6cm, minimum height=1.5cm, align=center, text width=6.2cm] {« Je refuse ce marché\\par honnêteté » (commerçant de Douala)}; &
\node[draw=popline, rounded corners, minimum width=6.6cm, minimum height=1.5cm, align=center, text width=6.2cm] {Peur du scandale, souci de sa réputation : la conscience embellit le mobile (Nietzsche)}; \\[0.25cm]
\node[draw=popline, rounded corners, minimum width=6.6cm, minimum height=1.5cm, align=center, text width=6.2cm] {« J'ai oublié le rendez-vous\\par hasard » (élève de Bafoussam)}; &
\node[draw=popline, rounded corners, minimum width=6.6cm, minimum height=1.5cm, align=center, text width=6.2cm] {Angoisse refoulée devant l'épreuve : l'oubli obéit à une intention cachée (Freud)}; \\
};
\end{tikzpicture}
\end{center}
\legende{Tableau comparatif : trois exemples camerounais montrant l'écart entre ce que la conscience croit et ce qui détermine réellement l'acte.}
\end{popfigure}

\courssub{Conclusion de section : une conscience précieuse mais limitée}

\begin{propriete}[Thèse à retenir]
La conscience est \cle{limitée} parce qu'elle ignore les déterminations — sociales, biologiques et inconscientes — qui la conditionnent. Elle se trompe sur le monde (illusions des sens), elle se trompe sur elle-même (mauvaise foi), et elle ignore les causes profondes de ses propres actes (déterminismes sociaux, corporels, pulsionnels). (Notion admise, aucune démonstration n'est exigible au Bac.)
\end{propriete}

\begin{attention}[Piège de dissertation : ne pas dévaloriser totalement la conscience]
L'erreur fréquente consiste à conclure : « la conscience ne vaut rien, tout est déterminé ». C'est excessif et contradictoire : pour découvrir ses limites, il a fallu... la conscience elle-même (réflexion, science, philosophie). La bonne conclusion est nuancée : la conscience est \textbf{nécessaire mais non suffisante}. Elle est nécessaire, car sans elle pas de liberté, pas de responsabilité, pas de connaissance ; elle n'est pas suffisante, car elle n'épuise pas la vie mentale et ne maîtrise pas tout ce qui la détermine. La lucidité consiste à connaître ses limites pour mieux les repousser.
\end{attention}

Ainsi, la conscience n'est pas toute la vie mentale : lapsus, rêves, oublis, actes manqués en témoignent. Mais alors, qu'y a-t-il \emph{sous} la conscience ? Comment nommer et étudier cette vie psychique cachée qui agit en nous à notre insu ? C'est la révolution opérée par Sigmund Freud avec la découverte de l'\cle{inconscient}, que nous allons maintenant examiner.

\coursec{L'inconscient freudien}

Nous venons de voir, dans la section précédente, que la conscience a des \cle{limites} : elle ne connaît pas tout d'elle-même, elle se trompe, elle est soumise à l'illusion. Mais faut-il aller plus loin ? Et s'il existait en nous une véritable \emph{autre scène}, où se joueraient des pensées, des désirs, des souvenirs que nous ignorons totalement, mais qui agissent sur nous ? C'est l'hypothèse audacieuse que formule Sigmund Freud à la fin du XIX\textsuperscript{e} siècle : l'hypothèse de l'\cle{inconscient}. Cette section présente cette hypothèse, son architecture (les deux « topiques »), son mécanisme central (le refoulement) et ses manifestations (rêves, lapsus, actes manqués, symptômes), avant d'en évaluer la portée philosophique et les limites épistémologiques.

\courssub{Le contexte : Freud et la naissance de la psychanalyse}

\textbf{Un médecin viennois face aux maladies sans cause visible.} Sigmund \cle{Freud} (1856-1939) est un médecin neurologue viennois. À la fin du XIX\textsuperscript{e} siècle, il reçoit dans son cabinet des patients, surtout des femmes, atteints d'\emph{hystérie} : paralysies, aveuglements, toux, angoisses... sans qu'aucune lésion organique ne soit décelable. La médecine de l'époque est impuissante. Freud fait alors une découverte décisive : lorsque le patient, invité à parler librement (c'est la méthode des \emph{associations libres}), retrouve le souvenir douloureux oublié à l'origine du trouble, le symptôme s'atténue ou disparaît. Conclusion révolutionnaire : le symptôme a une \emph{cause psychique}, mais cette cause est \emph{inconsciente} — le patient l'ignore.

\textbf{La fondation de la psychanalyse.} De cette pratique naît la \cle{psychanalyse}, à la fois méthode d'investigation du psychisme, méthode thérapeutique (la « cure par la parole ») et théorie de l'appareil psychique. Son ouvrage fondateur est \emph{L'Interprétation des rêves} (1900).

\begin{exemplebox}[Un travail de longue haleine]
Freud consacre plus de \textbf{600 pages} de \emph{L'Interprétation des rêves} (1900) à l'analyse de ses propres rêves et de ceux de ses patients. Il y montre, rêve après rêve, que chaque rêve a un sens caché : il est la réalisation déguisée d'un désir refoulé. Le rêve n'est donc pas un non-sens, mais un texte à déchiffrer.
\end{exemplebox}

\courssub{L'hypothèse de l'inconscient}

\begin{definition}[L'inconscient]
L'\cle{inconscient} désigne, chez Freud, l'ensemble des contenus psychiques (désirs, souvenirs, représentations, conflits) qui échappent à la conscience mais qui \emph{agissent} sur elle : ils déterminent nos conduites, nos émotions, nos symptômes, sans que nous le sachions.
\end{definition}

\textbf{L'intuition d'abord.} Avez-vous déjà dit un mot « à la place d'un autre » et rougi en vous demandant pourquoi ? Avez-vous déjà éprouvé une colère disproportionnée pour un détail insignifiant ? Oublié précisément le rendez-vous que vous redoutiez ? Dans tous ces cas, il semble que « quelque chose en nous » sache ce que nous ne voulons pas savoir. Telle est l'intuition de départ : \emph{nous ne sommes pas transparents à nous-mêmes}.

\textbf{L'énoncé rigoureux.} L'inconscient n'est pas une simple absence de conscience (comme le sommeil ou l'évanouissement) : c'est une \emph{activité} psychique dotée de ses propres lois. Ses contenus ont été activement \emph{tenus à distance} de la conscience, et ils cherchent sans cesse à y revenir, sous des formes déguisées. Freud résume sa thèse par une formule célèbre : \og le moi n'est pas maître dans sa propre maison \fg{}. Autrement dit, je ne suis pas le maître absolu de mes pensées et de mes actes : une partie de ma vie psychique m'échappe.

\begin{attention}[Piège fréquent]
Ne présentez jamais l'inconscient comme un \emph{lieu} du cerveau, une « pièce » ou un « organe » localisé quelque part dans la tête. L'inconscient est un \textbf{concept fonctionnel} : il désigne un mode de fonctionnement du psychisme (des contenus actifs mais inaccessibles à la conscience), non une région anatomique que le chirurgien pourrait montrer du doigt.
\end{attention}

\courssub{La première topique : conscience, préconscient, inconscient}

Freud appelle \cle{topique} (du grec \emph{topos}, lieu) une représentation de l'appareil psychique en « lieux » ou systèmes. La première topique (vers 1900) distingue trois systèmes :

\begin{itemize}
\item La \cle{conscience} : ce que je sais de moi à l'instant présent — mes perceptions, mes pensées actuelles. C'est la partie émergée de l'iceberg.
\item Le \cle{préconscient} : ce qui n'est pas actuellement conscient mais peut le devenir facilement — mes souvenirs disponibles, mes connaissances (par exemple, la capitale du Cameroun, que je ne pense pas en ce moment mais que je peux rappeler à volonté).
\item L'\cle{inconscient} : ce qui est tenu hors de la conscience par une force active et ne peut y accéder directement — désirs inacceptables, souvenirs traumatiques, conflits refoulés.
\end{itemize}

Entre l'inconscient et la conscience se tient une \cle{censure}, comparable à un portier qui examine ce qui demande à passer et refoule ce qui est inacceptable. C'est elle qui explique que les contenus inconscients ne réapparaissent que \emph{déguisés}.

\courssub{La seconde topique : le ça, le moi, le surmoi}

Vers 1920, Freud propose un second modèle, plus dynamique, qui décrit le psychisme comme le théâtre d'un \emph{conflit} entre trois instances :

\begin{itemize}
\item Le \cle{ça} : le réservoir des \emph{pulsions} (désirs sexuels, agressivité, besoins). Il obéit au \emph{principe de plaisir} : il veut la satisfaction immédiate, sans tenir compte de la réalité ni de la morale. Le nourrisson est presque tout entier « ça » : il crie dès qu'il a faim.
\item Le \cle{moi} : l'instance de la \emph{réalité}. Il obéit au \emph{principe de réalité} : il négocie entre les exigences du ça, les interdits du surmoi et les possibilités du monde extérieur. C'est un médiateur, souvent épuisé, qui doit « servir trois maîtres ».
\item Le \cle{surmoi} : l'instance des \emph{interdits} et des idéaux. Il est la morale intériorisée : les « il faut », « il ne faut pas » transmis par les parents, l'école, la religion, la société. Il juge le moi, culpabilise, interdit.
\end{itemize}

\begin{popfigure}
\begin{center}
\begin{tikzpicture}[scale=1]
% Le ça
\draw[fill=poppinkL,draw=poppink,thick] (0,0) ellipse (3.2 and 1.1);
\node[poppink,font=\bfseries] at (0,-0.75) {Le \c{c}a (pulsions, principe de plaisir)};
% Le surmoi
\draw[fill=poppurpleL,draw=poppurple,thick] (0,2.6) ellipse (2.4 and 0.9);
\node[poppurple,font=\bfseries] at (0,3.15) {Le surmoi (interdits, morale)};
% Le moi
\draw[fill=popblueL,draw=popblue,very thick] (0,1.3) circle (0.85);
\node[popblue,font=\bfseries] at (0,1.3) {Le moi};
\node[popblue,font=\footnotesize] at (0,1.05) {(r\'ealit\'e)};
% Flèches de conflit
\draw[-{Stealth},very thick,popink] (0,0.6) -- (0,0.42);
\draw[{Stealth}-{Stealth},very thick,poporange] (-0.5,0.45) -- (-0.6,2.0);
\draw[{Stealth}-{Stealth},very thick,poporange] (0.5,0.45) -- (0.6,2.0);
\node[poporange,font=\footnotesize,align=center] at (2.6,1.3) {conflit : d\'esirs\\contre interdits};
\end{tikzpicture}
\end{center}
\legende{La seconde topique freudienne. Le moi (au centre) est pris entre les pulsions du ça (en bas) et les interdits du surmoi (en haut) : la vie psychique est un conflit permanent que le moi doit arbitrer.}
\end{popfigure}

\begin{exemplebox}[Le conflit psychique en situation camerounaise]
Un élève de Terminale, la veille d'une évaluation, a très envie d'aller au match de football du quartier (le \emph{ça} : plaisir immédiat). Une voix intérieure lui dit : « Un bon fils travaille, tu décevras ta famille » (le \emph{surmoi} : interdit intériorisé). Le \emph{moi} négocie : « Je révise deux heures, puis je regarde le résumé du match. » La vie psychique quotidienne est faite de tels arbitrages.
\end{exemplebox}

\courssub{Le refoulement : le mécanisme central}

Comment un désir devient-il inconscient ? Par le \cle{refoulement}, mécanisme fondamental de la psychanalyse.

\begin{definition}[Le refoulement]
Le \cle{refoulement} est un mécanisme de défense par lequel le moi maintient hors de la conscience les représentations, les souvenirs et les désirs incompatibles avec lui-même — c'est-à-dire ceux qui blesseraient l'image que je me fais de moi ou qui violeraient les interdits du surmoi.
\end{definition}

\textbf{Établissement du raisonnement.} Partons d'une observation clinique : les patientes hystériques de Freud avaient « oublié » des scènes pénibles de leur passé, mais cet oubli n'était pas une simple disparition — le souvenir continuait d'agir (sous forme de symptôme) et pouvait être retrouvé par la parole. Il faut donc admettre une force qui a \emph{chassé} le contenu (la censure) et une énergie qui le maintient actif à distance. Le refoulé ne disparaît pas : il est comme un ballon maintenu sous l'eau — dès que la pression se relâche (sommeil, fatigue, distraction), il remonte à la surface, mais déformé, déguisé.

\begin{propriete}[Le retour déguisé du refoulé]
L'inconscient ne se manifeste jamais directement, mais toujours de manière \emph{déguisée} (rêves, lapsus, symptômes), parce que ses contenus, refoulés par la censure, ne peuvent franchir la barrière de la conscience qu'en se camouflant. \emph{Justification} : si un contenu refoulé apparaissait tel quel, la censure le rejetterait à nouveau ; il doit donc se transformer (par condensation, déplacement, symbolisation) pour devenir acceptable — c'est précisément ce camouflage qui rend l'interprétation nécessaire.
\end{propriete}

\begin{popfigure}
\begin{center}
\begin{tikzpicture}[scale=1,font=\small]
% Boîtes
\node[draw=popgreen,fill=popgreenL,rounded corners,minimum width=2.4cm,minimum height=0.9cm,align=center] (desir) at (0,0) {D\'esir ou souvenir\\inacceptable};
\node[draw=popink,fill=poppinkL,rounded corners,minimum width=2.2cm,minimum height=0.9cm,align=center] (censure) at (4.2,0) {Barri\`ere de\\la censure};
\node[draw=poppurple,fill=poppurpleL,rounded corners,minimum width=2.2cm,minimum height=0.9cm,align=center] (inc) at (8.4,0) {Inconscient\\(refoul\'e)};
\node[draw=poporange,fill=poporangeL,rounded corners,minimum width=2.8cm,minimum height=1.1cm,align=center] (retour) at (8.4,-2.4) {Retour d\'eguis\'e :\\r\^eve, lapsus,\\sympt\^ome};
% Flèches
\draw[-{Stealth},very thick,popdark] (desir) -- (censure) node[midway,above,font=\footnotesize]{refoulement};
\draw[-{Stealth},very thick,popdark] (censure) -- (inc);
\draw[-{Stealth},very thick,poporange] (inc.south) -- (retour.north) node[midway,right,font=\footnotesize]{quand la censure se rel\^ache};
\draw[-{Stealth},dashed,thick,popmuted] (retour.west) .. controls (4.2,-3.2) .. (0,-0.6) node[midway,below,font=\footnotesize,align=center]{appara\^it dans la conscience\\sans \^etre reconnu};
\end{tikzpicture}
\end{center}
\legende{Le circuit du refoulement. Le désir inacceptable est chassé par la censure dans l'inconscient, d'où il revient sous des formes déguisées (rêve, lapsus, symptôme) que la conscience ne reconnaît pas.}
\end{popfigure}

\courssub{Les manifestations de l'inconscient}

\textbf{Les rêves.} Pour Freud, le rêve est \og la voie royale vers l'inconscient \fg{}. Pendant le sommeil, la censure se relâche : les désirs refoulés remontent, mais encore déguisés. Freud distingue le \emph{contenu manifeste} (le rêve tel qu'on le raconte, souvent bizarre) et le \emph{contenu latent} (le désir caché que le rêve réalise symboliquement). Le travail du rêve est un travail de transformation : condensation (plusieurs idées fondues en une image) et déplacement (l'émotion glisse d'un objet important vers un objet anodin).

\textbf{Les lapsus et les actes manqués.} Dans \emph{Psychopathologie de la vie quotidienne} (1901), Freud montre que nos « erreurs » quotidiennes ne sont pas des hasards.

\begin{definition}[Le lapsus]
Le \cle{lapsus} est une erreur de parole ou d'action (mot substitué, oubli, objet égaré, geste raté) qui trahit un désir ou une pensée inconsciente que le sujet voulait cacher — y compris à lui-même.
\end{definition}

\textbf{Les symptômes névrotiques.} Phobies, obsessions, angoisses sans objet apparent : le symptôme est un \emph{compromis} entre le désir refoulé qui veut s'exprimer et la censure qui l'interdit. Il a un sens, que la cure analytique aide à déchiffrer.

\begin{exoresolu}[Analyser deux exemples à la manière de Freud]
\textbf{Énoncé.} (1) Un élève rêve, la veille du Baccalauréat technique, qu'il arrive en salle d'examen et que sa feuille reste désespérément blanche. (2) Un chef de famille, présentant son épouse à des invités, dit : « Je vous présente ma mère. » Analysez ces deux cas avec les notions freudiennes.
\tcblower
\textbf{Solution.} (1) Le rêve d'examen est le \emph{contenu manifeste}. Son \emph{contenu latent} est une \emph{angoisse refoulée} : l'élève, conscient de l'enjeu, n'ose pas s'avouer sa peur de l'échec (peur incompatible avec l'image d'un élève confiant — le moi — et avec l'exigence familiale de réussite — le surmoi). Pendant le sommeil, la censure se relâche et l'angoisse remonte, déguisée en scène d'examen. Le rêve révèle donc un conflit inconscient. (2) Le lapsus du chef de famille substitue « ma mère » à « mon épouse » : il trahit une pensée inconsciente — par exemple, le fait qu'il attend de son épouse les mêmes soins maternels que de sa mère, ou qu'il perçoit inconsciemment son épouse comme une figure d'autorité. L'erreur n'est pas un hasard : elle dit une vérité que le sujet n'ose pas formuler. Dans les deux cas, l'inconscient se manifeste de manière déguisée, conformément au mécanisme du refoulement.
\end{exoresolu}

\courssub{Conséquence philosophique : la remise en cause du cogito}

La thèse freudienne a une portée philosophique immense. Depuis Descartes, la philosophie moderne reposait sur le \cle{cogito} : « je pense, donc je suis » — la conscience se connaît elle-même de façon immédiate et certaine ; elle est transparente à elle-même. Freud renverse cette certitude : si une partie de ma psyché m'échappe et me détermine, alors \emph{je ne suis pas ce que je crois être}. La conscience n'est plus le fondement, mais une surface ; le sujet n'est plus maître, mais habité par des forces qui le dépassent. D'où la formule : \og le moi n'est pas maître dans sa propre maison \fg{}.

\begin{savaistu}[Les trois blessures narcissiques de l'humanité]
Freud remarquait que la science a infligé trois grandes humiliations à l'orgueil humain, trois « blessures narcissiques » : \textbf{Copernic} a montré que la Terre n'est pas le centre de l'univers (blessure cosmologique) ; \textbf{Darwin} a montré que l'homme descend de l'animal (blessure biologique) ; et \textbf{Freud} lui-même a montré que le moi n'est pas maître chez lui (blessure psychologique). L'inconscient achève ainsi de détrôner l'homme : il n'est ni au centre du monde, ni au-dessus de la nature, ni même au centre de lui-même.
\end{savaistu}

\courssub{Valeur et limites de la psychanalyse}

\courssub{La valeur heuristique et thérapeutique}

La psychanalyse a opéré une révolution dans la compréhension de l'humain. Sa \cle{valeur} est triple :
\begin{itemize}
\item \textbf{Thérapeutique :} elle a inventé la « cure par la parole », montrant que la mise en mots (l'association libre, l'interprétation des rêves, du transfert) peut lever des symptômes corporels là où la médecine organique échoue. Elle a fondé les psychothérapies modernes.
\item \textbf{Heuristique (explicative) :} elle fournit une grille de lecture puissante pour des phénomènes restés obscurs : l'irrationnel dans le comportement, l'art, la religion, les mythes, les névroses collectives, les slips de la vie quotidienne. Elle a irriguer l'anthropologie (Lévi-Strauss), la littérature, l'histoire, la sociologie.
\item \textbf{Anthropologique :} elle replace l'homme dans sa condition de \emph{sujet divisé}, tiraillé entre nature (pulsions) et culture (interdits), fondant une éthique de la lucidité : « Là où était le ça, le moi doit advenir ».
\end{itemize}

\begin{exemplebox}[Le transfert, moteur de la cure]
Le \cle{transfert} est le phénomène par lequel le patient reporte sur l'analyste des affects, des désirs, des attentes liés à des figures du passé (parents). Bien géré (contre-transfert maîtrisé), il devient le levier thérapeutique : revivre la scène originelle en présence d'un « autre » bienveillant permet de l'élaborer différemment.
\end{exemplebox}

\courssub{Les limites épistémologiques et critiques}

Malgré sa fécondité, la psychanalyse fait l'objet de critiques majeures qu'il faut connaître pour le Bac :

\begin{itemize}
\item \textbf{Statut scientifique (Karl Popper) :} La psychanalyse serait une \emph{pseudo-science} car ses thèses ne sont pas \emph{réfutables} (falsifiables). Toute observation (guérison, échec, résistance du patient, silence) est interprétée comme confirmant la théorie. Exemple : si le patient nie l'interprétation, c'est de la « résistance » ; s'il l'accepte, c'est une confirmation. La théorie devient irréfutable, donc non scientifique.
\item \textbf{Validité des données cliniques (Adolf Grünbaum) :} Les « preuves » freudiennes reposent sur des études de cas uniques, non reproductibles, sans groupe témoin, où le thérapeute suggère involontairement les réponses attendues (biais de confirmation). L'efficacité spécifique de la psychanalyse (vs placebo, vs soutien amical, vs thérapies brèves) reste difficile à établir scientifiquement.
\item \textbf{Determinisme et responsabilité :} En expliquant l'acte par des causes inconscientes échappant au sujet, la psychanalyse risque de dissoudre la \cle{responsabilité} morale et juridique. « C'est mon inconscient qui a fait ça » devient une excuse commode.
\item \textbf{Biais culturels et sexués :} Construite sur une clinique viennoise bourgeoise de 1900, la théorie universalise des normes patriarcales (complexe d'Œdipe phallocentrique, « envie du pénis », place assignée à la femme). Des penseuses comme Simone de Beauvoir, Luce Irigaray ou Judith Butler ont dénoncé un discours qui naturalise la domination masculine.
\item \textbf{Surinterprétation sexuelle :} Freud réduit souvent la pulsion à la sexualité infantile (polymorphe perverse), négligeant d'autres moteurs : la recherche de reconnaissance, l'attachement, la volonté de puissance, le sens (Viktor Frankl, logothérapie).
\end{itemize}

\begin{attention}[Nuance attendue au Bac]
Dire « la psychanalyse n'est pas scientifique » est une affirmation de Popper, pas un fait établi indiscutable. Au Bac, on \textbf{problématise} : \og En quel sens la psychanalyse est-elle une science ? Une herméneutique ? Une éthique ? \fg{} Paul Ricœur propose de la lire comme une \emph{herméneutique de la suspicion} (avec Marx et Nietzsche) : elle déchiffre des significations cachées, non comme une physique du psychisme. Cette distinction « sciences de la nature / sciences de l'esprit » (Dilthey) est cruciale pour une dissertation réussie.
\end{attention}

\begin{savaistu}[Ricœur et l'archéologie du sujet]
Pour Paul Ricœur (\emph{De l'interprétation}, 1965), Freud n'est pas un physicien de l'âme mais un \emph{herméneute}. L'inconscient n'est pas un « objet » qu'on observe au microscope, mais un \emph{sens} qu'on déchiffre à travers des signes (rêves, lapsus). La psychanalyse est une « archéologie » : elle reconstitue l'histoire profonde du sujet à partir de traces fragmentées. Le sujet freudien n'est pas une substance, c'est un \emph{sujet attribué} : je me constitue en m'attribuant mes actes, y compris ceux qui m'échappent.
\end{savaistu}

\begin{aretenir}[L'essentiel de la section]
\begin{itemize}
\item Freud, médecin viennois, fonde la psychanalyse (\emph{L'Interprétation des rêves}, 1900) : une partie de notre psyché nous échappe mais agit sur nous.
\item Première topique : conscience / préconscient / inconscient, séparés par la censure.
\item Seconde topique : le ça (pulsions), le moi (réalité), le surmoi (interdits) — le psychisme est un conflit.
\item Le refoulement maintient hors de la conscience les désirs inacceptables ; le refoulé revient déguisé : rêves (« voie royale »), lapsus, actes manqués, symptômes.
\item Portée philosophique : remise en cause du cogito — « le moi n'est pas maître dans sa propre maison ».
\item Valeur : révolution thérapeutique (cure par la parole), grille de lecture de la culture, éthique de la lucidité.
\item Limites : statut scientifique discuté (Popper, Grünbaum), risque de déterminisme absolvant, biais patriarcaux, réductionnisme sexuel.
\end{itemize}
\end{aretenir}

\courssub{Méthode et entraînement au Bac}

\begin{methode}[Dissertation : traiter un sujet sur l'inconscient]
\begin{enumerate}
\item \textbf{Analyser le sujet} : identifier les termes clés (inconscient, responsabilité, liberté, rêve, savoir, maître), les oppositions implicites (conscient/inconscient, maîtrise/aliénation, nature/culture), la thèse implicite à discuter.
\item \textbf{Problématiser} : transformer le sujet en question philosophique. Exemple : « L'inconscient est-il une excuse ou une explication ? » $\rightarrow$ « La découverte de l'inconscient détruit-elle la responsabilité morale ou la refonde-t-elle sur une lucidité nouvelle ? »
\item \textbf{Construire le plan dialectique} (Thèse / Antithèse / Synthèse) ou thématique (Psychologique / Épistémologique / Éthique).
\item \textbf{Rédiger} : Introduction (accroche, définitions, problème, annonce du plan), Développement (arguments étayés par auteurs \textbf{et} exemples concrets), Conclusion (réponse synthétique, ouverture).
\end{enumerate}
\end{methode}

\begin{methode}[Commentaire de texte : expliquer un extrait freudien ou critique]
\begin{enumerate}
\item \textbf{Identifier} : Auteur, œuvre, date, contexte (Freud, \emph{L'Interprétation des rêves}, 1900 ; ou Ricœur, \emph{De l'interprétation}, 1965).
\item \textbf{Structurer} : Découper le texte en 2 ou 3 parties logiques (thèses successives).
\item \textbf{Expliquer} : Pour chaque partie, reformuler la thèse, expliciter les concepts techniques (refoulement, censure, condensation, déplacement, hermèneutique), montrer le raisonnement (arguments, exemples du texte).
\item \textbf{Évaluer (Discussion critique)} : Confronter la thèse à d'autres auteurs (Sartre : « mauvaise foi » vs inconscient ; Popper : scientificité ; Ricœur : herméneutique vs science). Conclure sur la portée et les limites de la pensée de l'auteur.
\end{enumerate}
\end{methode}

\begin{exercice}[Sujets types Probatoire / Baccalauréat Technique]
\begin{enumerate}
\item \textbf{Dissertation :} « L'inconscient nous excuse-t-il d'être responsables ? »
\item \textbf{Dissertation :} « Le rêve est-il la voie royale vers l'inconscient ? »
\item \textbf{Commentaire :} Extrait de Freud, \emph{Introduction à la psychanalyse} (1916-1917), Conférence XVII : « Le sens des symptômes » (sur le compromis entre désir et défense).
\item \textbf{Commentaire :} Extrait de Paul Ricœur, \emph{De l'interprétation} (1965) : « La psychanalyse comme herméneutique de la suspicion ».
\end{enumerate}
\end{exercice}

\begin{corrige}[Exercice 1 : Dissertation « L'inconscient nous excuse-t-il d'être responsables ? »]
\textbf{Plan dialectique suggéré :}
\textbf{Introduction.} Accroche : formule de Freud « le moi n'est pas maître dans sa propre maison ». Définition : l'inconscient comme ensemble de déterminants psychiques échappant à la conscience. Problème : si mes actes viennent de l'inconscient, suis-je encore auteur et responsable ? Annonce : 1) L'inconscient menace la responsabilité (thèse) ; 2) Mais la responsabilité peut se redéfinir par la lucidité (antithèse) ; 3) Vers une responsabilité assumée, non malgré l'inconscient, mais \emph{grâce} à son élucidation (synthèse).
\textbf{Développement.}
I. \textbf{La menace : l'inconscient comme cause étrangère.} Argument : déterminisme psychique (Freud), l'acte a une cause que je ne choisis pas (pulsion refoulée, traumatisme). Exemple : l'acte manqué, le symptôme, le crime « pulsionnel ». Conséquence : irresponsabilité pénale (alinéa 1 de l'article 122-1 du Code pénal camerounais / français : abolition du discernement), irresponsabilité morale (je ne suis pas l'auteur).
II. \textbf{La résistance : la responsabilité comme appropriation.} Argument : Sartre (\emph{L'Être et le Néant}) : l'inconscient est une « mauvaise foi », je me cache ma liberté. Ricœur : le sujet se constitue en s'attribuant ses actes (\emph{Soimême comme un autre}). La cure analytique vise l'autonomie : « Là où était le ça, le moi doit advenir ». Exemple : l'aveu en justice, la réparation symbolique en thérapie.
III. \textbf{La synthèse : responsabilité lucide.} L'inconnu ne m'excuse pas, il m'oblige. Responsabilité \emph{élargie} : je suis responsable de \emph{connaître} mes mobiles profonds (devoir de vérité sur soi). Responsabilité \emph{partagée} : la société doit soigner (psychiatrie) plutôt que seulement punir. La responsabilité n'est pas la toute-puissance du cogito, c'est la capacité de répondre de ses actes \emph{en connaissant ses failles}.
\textbf{Conclusion.} L'inconscient ne supprime pas la responsabilité, il la complexifie. Elle passe de « je suis maître de tout » à « je réponds de ce que je suis, y compris de ce qui m'échappe ». Ouverture : enjeu éthique de l'IA et des algorithmes prédictifs : qui est responsable quand la « boîte noire » décide ?
\end{corrige}

\coursec{Valeur et limites de la psychanalyse}

Nous avons vu dans la section précédente que Freud postule l'existence d'un \cle{inconscient}, instance psychique dont les contenus refoulés déterminent une grande partie de nos conduites. Mais cette théorie est-elle scientifiquement fondée ? A-t-elle une utilité réelle ? Ne risque-t-elle pas, au contraire, de dissoudre la liberté humaine ? Toute théorie, aussi séduisante soit-elle, doit être soumise à l'épreuve de la critique. C'est le sens de la question qui guide cette section : \emph{quelle est la valeur de la psychanalyse, et quelles sont ses limites ?} Nous montrerons d'abord sa triple valeur (théorique, thérapeutique, culturelle), puis ses limites (scientifique, méthodologique, philosophique), avant de proposer un bilan nuancé, conforme à l'exigence dialectique de la dissertation.

\courssub{La valeur de la psychanalyse}

\textbf{1. Valeur théorique : élargir la connaissance de l'homme}\par

L'intuition de départ est simple : avant Freud, la psychologie occidentale identifiait l'esprit à la conscience. Être un sujet, c'était être transparent à soi-même, comme le voulait Descartes. Or l'observation quotidienne contredit cette transparence : lapsus, oublis, actes manqués, rêves, symptômes névrotiques sont autant de conduites que la conscience seule ne peut expliquer. La psychanalyse a donc une première valeur : elle \textbf{élargit la connaissance de l'homme au-delà de la conscience}. Elle rend intelligibles des comportements jusque-là jugés absurdes ou mystérieux.

\begin{exemplebox}[L'acte manqué expliqué]
Un élève de Terminale, la veille de l'examen, « oublie » sa convocation. La psychologie classique dira : simple négligence. La psychanalyse propose une interprétation : cet oubli exprime un désir inconscient (la peur de l'échec, le refus de l'épreuve). L'acte manqué devient ainsi un \emph{message} déchiffrable, non un accident insignifiant. Même si l'on conteste cette lecture, il faut reconnaître à Freud le mérite d'avoir posé une question neuve : et si nos conduites avaient un sens caché ?
\end{exemplebox}

\textbf{2. Valeur thérapeutique : la cure par la parole}\par

La psychanalyse n'est pas seulement une théorie : elle est d'abord une \textbf{pratique de soin}. Freud, partant des travaux de Breuer sur l'hystérie, met au point la \cle{cure par la parole} : le patient, allongé sur le divan, dit tout ce qui lui vient à l'esprit (\emph{association libre}), et l'analyste interprète ses propos, ses rêves, ses résistances. En rendant conscient le conflit refoulé, la parole soulage le symptôme.

\begin{exemplebox}[Le cas d'Anna O.]
Bertha Pappenheim, dite « Anna O. », patiente de Breuer (1880-1882), souffrait de paralysies et de troubles de la vue sans cause organique. En racontant, sous hypnose puis librement, les souvenirs liés à la maladie de son père, elle vit ses symptômes disparaître. Elle baptisa elle-même cette méthode « cure par la parole » (\emph{talking cure}). Ce cas fondateur montre la valeur thérapeutique de l'approche : des troubles que la médecine du corps ne guérissait pas furent soulagés par la parole et la remémoration.
\end{exemplebox}

\textbf{3. Valeur culturelle : une influence considérable}\par

Au-delà de la clinique, la psychanalyse a profondément marqué la culture du XX\textsuperscript{e} siècle :
\begin{itemize}
\item \textbf{Littérature} : l'exploration des profondeurs du moi (le « courant de conscience » chez Joyce ou Woolf), l'analyse psychologique des personnages.
\item \textbf{Cinéma} : les films d'Hitchcock (\emph{La Maison du docteur Edwardes}, \emph{Psychose}) mettent en scène refoulement, traumatismes et rêves.
\item \textbf{Éducation} : l'attention portée à l'enfance, aux désirs de l'enfant, à la qualité du lien parents-enfants.
\item \textbf{Interprétation des mythes} : Freud lit le mythe d'{\OE}dipe ou les totems comme expressions de conflits psychiques universels.
\end{itemize}

\courssub{Les limites de la psychanalyse}

\textbf{1. La limite scientifique : la critique de Karl Popper}\par

Le philosophe des sciences Karl \cle{Popper} (1902-1994) a formulé la critique la plus célèbre. Pour lui, une théorie n'est scientifique que si elle accepte de \emph{risquer} quelque chose devant l'expérience.

\begin{definition}[Réfutabilité (Popper)]
La \cle{réfutabilité} (ou falsifiabilité) est le critère selon lequel une théorie n'est scientifique que si elle peut être \textbf{mise en échec par l'expérience} : il doit exister des observations possibles qui, si elles se produisaient, prouveraient que la théorie est fausse. Une théorie compatible avec tous les faits imaginables n'affirme rien de précis et n'est donc pas scientifique au sens strict.
\end{definition}

Or, selon Popper, la psychanalyse \emph{explique tout} : si le patient accepte l'interprétation de l'analyste, cela confirme la théorie ; s'il la refuse, ce refus est interprété comme une « résistance »... qui confirme encore la théorie ! Une théorie qui ne risque rien ne peut pas être réfutée : elle n'est donc pas une science expérimentale, au même titre que l'astrologie, que Popper oppose à la physique d'Einstein (qui, elle, prenait le risque d'être démentie par l'observation de l'éclipse de 1919).

\textbf{2. La limite méthodologique}\par

Deux reproches complètent cette critique :
\begin{itemize}
\item Freud a construit sa théorie générale de l'esprit humain à partir d'un \textbf{nombre restreint de cas cliniques}, issus d'un milieu social particulier (la bourgeoisie viennoise de 1900). Peut-on légitimement généraliser à toute l'humanité ?
\item L'interprétation psychanalytique comporte un \textbf{risque d'arbitraire} : un même rêve peut recevoir des lectures différentes selon les analystes, sans procédure claire pour trancher entre elles.
\end{itemize}

\textbf{3. Les critiques philosophiques : Sartre et la liberté}\par

Jean-Paul \cle{Sartre} (1905-1980) reproche à l'inconscient freudien de \textbf{dissoudre la liberté et la responsabilité} humaines. Si mes actes sont déterminés par des forces inconscientes, alors je ne suis plus responsable de ce que je fais : « ce n'est pas moi, c'est mon inconscient ». Pour Sartre, l'homme est « condamné à être libre » : il se choisit à chaque instant. Ce que Freud attribue à l'inconscient, Sartre l'explique par la \cle{mauvaise foi} : le sujet se ment à lui-même, se cache sa propre liberté, mais il n'est pas pour autant divisé en deux instances. La mauvaise foi rend compte des mêmes phénomènes \emph{sans postuler un inconscient}.

\textbf{4. Les critiques contemporaines : les neurosciences}\par

Les \cle{neurosciences} contemporaines proposent d'autres modèles du fonctionnement mental : les processus non conscients (perception implicite, mémoire automatique, prise de décision rapide) y sont expliqués par le fonctionnement cérébral, mesurable par l'imagerie médicale, sans référence au refoulement ni au complexe d'{\OE}dipe. Ces approches, expérimentales et réfutables, concurrencent le modèle freudien, même si certains chercheurs tentent un dialogue (« neuropsychanalyse »).

\begin{attention}[Erreur fréquente à éviter]
N'écris jamais dans une copie que Freud a « \textbf{prouvé} » ou « \textbf{démontré} » l'existence de l'inconscient. L'inconscient freudien est une \textbf{hypothèse interprétative}, non un résultat de démonstration expérimentale. On ne peut pas observer l'inconscient au microscope ni le mesurer : on le \emph{postule} pour rendre compte de certains faits (rêves, lapsus, névroses). La formulation correcte est : « Freud \emph{postule} l'existence de l'inconscient » ou « la psychanalyse \emph{propose une interprétation} des conduites humaines ».
\end{attention}

\begin{propriete}[Bilan sur le statut de la psychanalyse]
La psychanalyse possède une \textbf{valeur interprétative et thérapeutique reconnue} (elle éclaire les conduites, soulage certains troubles, féconde la culture), mais son \textbf{statut scientifique est contesté} : ses hypothèses sont difficilement vérifiables et réfutables au sens de Popper.
\end{propriete}

\begin{popfigure}
\begin{center}
\begin{tikzpicture}[scale=1, every node/.style={font=\small}]
% En-têtes
\fill[popblue!25] (-0.2,0) rectangle (3.4,0.9);
\fill[popgreen!25] (3.4,0) rectangle (7.4,0.9);
\fill[poporange!30] (7.4,0) rectangle (11.4,0.9);
\node[font=\small\bfseries] at (1.6,0.45) {Plan};
\node[font=\small\bfseries] at (5.4,0.45) {Valeur de la psychanalyse};
\node[font=\small\bfseries] at (9.4,0.45) {Limites};
% Ligne 1
\fill[popblue!8] (-0.2,-1.5) rectangle (3.4,0);
\fill[popgreen!8] (3.4,-1.5) rectangle (7.4,0);
\fill[poporange!10] (7.4,-1.5) rectangle (11.4,0);
\node[font=\small\bfseries, align=center] at (1.6,-0.75) {Théorique};
\node[align=center, text width=3.6cm] at (5.4,-0.75) {Élargit la connaissance de l'homme au-delà de la conscience};
\node[align=center, text width=3.6cm] at (9.4,-0.75) {Hypothèses peu vérifiables ; risque d'interprétation arbitraire};
% Ligne 2
\fill[popblue!8] (-0.2,-3) rectangle (3.4,-1.5);
\fill[popgreen!8] (3.4,-3) rectangle (7.4,-1.5);
\fill[poporange!10] (7.4,-3) rectangle (11.4,-1.5);
\node[font=\small\bfseries, align=center] at (1.6,-2.25) {Thérapeutique};
\node[align=center, text width=3.6cm] at (5.4,-2.25) {Cure par la parole : soulage les troubles névrotiques};
\node[align=center, text width=3.6cm] at (9.4,-2.25) {Efficacité difficile à mesurer ; cures longues et coûteuses};
% Ligne 3
\fill[popblue!8] (-0.2,-4.5) rectangle (3.4,-3);
\fill[popgreen!8] (3.4,-4.5) rectangle (7.4,-3);
\fill[poporange!10] (7.4,-4.5) rectangle (11.4,-3);
\node[font=\small\bfseries, align=center] at (1.6,-3.75) {Scientifique};
\node[align=center, text width=3.6cm] at (5.4,-3.75) {Modèle cohérent qui rend intelligibles rêves, lapsus, symptômes};
\node[align=center, text width=3.6cm] at (9.4,-3.75) {Irréfutable (Popper) ; concurrencée par les neurosciences};
% Cadre
\draw[popdark, thick] (-0.2,0.9) rectangle (11.4,-4.5);
\draw[popdark] (3.4,0.9) -- (3.4,-4.5);
\draw[popdark] (7.4,0.9) -- (7.4,-4.5);
\draw[popdark] (-0.2,0) -- (11.4,0);
\draw[popdark] (-0.2,-1.5) -- (11.4,-1.5);
\draw[popdark] (-0.2,-3) -- (11.4,-3);
\end{tikzpicture}
\end{center}
\legende{Tableau à double entrée : valeur et limites de la psychanalyse sur les plans théorique, thérapeutique et scientifique.}
\end{popfigure}

\begin{popfigure}
\begin{center}
\begin{tikzpicture}[scale=1, every node/.style={font=\small}]
\draw[-{Stealth[length=3mm]}, popdark, very thick] (0,0) -- (12.5,0);
% 1900
\draw[popblue, thick] (1.2,-0.15) -- (1.2,0.15);
\fill[popblue] (1.2,0) circle (0.09);
\node[above, align=center, text width=2.6cm, popblue] at (1.2,0.25) {\textbf{1900}\\ \emph{L'Interprétation des rêves} (Freud)};
% 1923
\draw[popblue, thick] (4.2,-0.15) -- (4.2,0.15);
\fill[popblue] (4.2,0) circle (0.09);
\node[below, align=center, text width=2.6cm, popblue] at (4.2,-0.25) {\textbf{1923}\\ \emph{Le Moi et le Ça} : seconde topique};
% Popper
\draw[poporange, thick] (7.2,-0.15) -- (7.2,0.15);
\fill[poporange] (7.2,0) circle (0.09);
\node[above, align=center, text width=3cm, poporange] at (7.2,0.25) {\textbf{Années 1930-1950}\\ Critique de Popper : la psychanalyse est irréfutable};
% Sartre
\draw[poppurple, thick] (9.2,-0.15) -- (9.2,0.15);
\fill[poppurple] (9.2,0) circle (0.09);
\node[below, align=center, text width=2.4cm, poppurple] at (9.2,-0.25) {\textbf{1943}\\ Sartre : la mauvaise foi contre l'inconscient};
% Neurosciences
\draw[popgreen, thick] (11.4,-0.15) -- (11.4,0.15);
\fill[popgreen] (11.4,0) circle (0.09);
\node[above, align=center, text width=2.4cm, popgreen] at (11.4,0.25) {\textbf{Aujourd'hui}\\ Neurosciences : autres modèles du mental};
\end{tikzpicture}
\end{center}
\legende{Frise chronologique : de la fondation de la psychanalyse à ses critiques contemporaines.}
\end{popfigure}

\courssub{Bilan nuancé : une herméneutique féconde}

Comment conclure sans tomber dans deux excès symétriques — l'admiration aveugle ou le rejet total ? Le bilan le plus solide consiste à reconnaître que la psychanalyse n'est pas une \textbf{science expérimentale} mais une \cle{herméneutique}, c'est-à-dire une \textbf{méthode d'interprétation} des conduites, des rêves et des discours humains. Comme telle, elle est féconde : elle a appris à l'humanité que l'homme n'est pas pure conscience transparente à elle-même. Mais l'inverse est également vrai : l'homme n'est pas non plus le pur jouet de forces inconscientes, car il conserve une capacité de choix, de réflexion et de responsabilité. La vérité sur le sujet humain se situe donc dans une tension : \emph{je ne suis pas totalement maître chez moi, mais je ne suis pas non plus totalement étranger à moi-même}.

\begin{methode}[Conclure une dissertation dialectique]
Face à un sujet du type « La psychanalyse est-elle une science ? » ou « L'inconscient nous détermine-t-il ? », procède ainsi :
\begin{enumerate}
\item \textbf{Thèse} : expose la valeur de la position (ici : la psychanalyse éclaire et soigne).
\item \textbf{Antithèse} : expose les objections les plus fortes (Popper, Sartre, neurosciences).
\item \textbf{Synthèse nuancée} : dépasse l'opposition en précisant les \textbf{conditions de validité} de chaque position. Exemple de formulation : « La psychanalyse est légitime comme méthode d'interprétation et de soin, mais ne peut prétendre au statut de science expérimentale ; l'inconscient éclaire une part de nos conduites sans abolir notre liberté. »
\item Termine par une ouverture ou une formule d'équilibre, jamais par un « tout est vrai » vague.
\end{enumerate}
\end{methode}

\begin{savaistu}[Rêves et traditions africaines]
Au Cameroun comme ailleurs en Afrique, l'interprétation traditionnelle des rêves occupe une place importante : songes prémonitoires, messages des ancêtres, avertissements transmis pendant le sommeil. Cette tradition offre un parallèle culturel intéressant avec l'herméneutique freudienne : dans les deux cas, le rêve est tenu pour un \emph{message à déchiffrer}, non pour un simple délire nocturne. La différence porte sur l'origine du message : pour Freud, il vient de notre propre inconscient ; dans les traditions africaines, il vient d'un au-delà ou d'une communauté d'ancêtres. Mobiliser ce parallèle dans un développement, avec prudence et respect, peut donner à ta copie du Bac une originalité remarquée.
\end{savaistu}

\begin{exoresolu}[Sujet de réflexion : « La psychanalyse est-elle une science ? »]
Construis en quelques lignes le plan dialectique d'une réponse à ce sujet.
\tcblower
\textbf{Corrigé.}
\begin{enumerate}
\item \textbf{Thèse} : la psychanalyse présente des traits scientifiques — elle part de l'observation clinique (cas d'Anna O.), formule des hypothèses générales (inconscient, refoulement) et obtient des résultats thérapeutiques vérifiables (disparition de symptômes).
\item \textbf{Antithèse} : selon le critère de Popper, elle n'est pas scientifique, car irréfutable : tout fait peut être interprété en sa faveur (l'accord du patient confirme, son désaccord est une « résistance »). De plus, ses généralisations reposent sur peu de cas et ses interprétations risquent l'arbitraire.
\item \textbf{Synthèse} : la psychanalyse n'est pas une science expérimentale au sens de la physique, mais une \emph{herméneutique} : une méthode d'interprétation du sens des conduites humaines, dotée d'une valeur thérapeutique et culturelle réelle. Sa légitimité est donc réelle mais d'un autre ordre que celle des sciences de la nature.
\end{enumerate}
\end{exoresolu}

\begin{exercice}[Pour t'entraîner]
\begin{enumerate}
\item Définis en deux phrases la \cle{réfutabilité} et explique pourquoi, selon Popper, la psychanalyse ne satisfait pas ce critère.
\item En quoi la notion sartrienne de \cle{mauvaise foi} permet-elle de rendre compte des mêmes faits que l'inconscient freudien, sans diviser le sujet ?
\item Rédige un paragraphe argumenté (10 à 15 lignes) répondant à la question : « L'inconscient détruit-il la responsabilité morale ? »
\end{enumerate}
\end{exercice}

\begin{corrige}[Exercice : Pour t'entraîner]
\begin{enumerate}
\item La réfutabilité est le critère selon lequel une théorie n'est scientifique que si des observations possibles peuvent la contredire. Or la psychanalyse interprète tout en sa faveur : l'accord du patient confirme la théorie, son désaccord est lu comme une « résistance » qui la confirme encore. Aucun fait ne peut donc l'invalider : elle est irréfutable, donc non scientifique au sens de Popper.
\item Pour Sartre, le sujet qui se trompe sur lui-même n'est pas victime d'une instance cachée : il se ment \emph{librement} à lui-même, il se cache sa propre liberté par mauvaise foi. Le sujet reste un et responsable ; il n'y a pas de « ça » qui agirait à sa place.
\item Éléments attendus : thèse (si mes actes sont déterminés par l'inconscient, je ne suis plus responsable : « ce n'est pas moi ») ; antithèse (Sartre : l'homme est condamné à être libre, la mauvaise foi explique l'auto-tromperie sans supprimer la responsabilité) ; synthèse (l'inconscient éclaire des motivations sans abolir la capacité de choix : connaître ses déterminations peut au contraire \emph{élargir} la liberté et la responsabilité).
\end{enumerate}
\end{corrige}

\begin{aretenir}[L'essentiel de la section]
\begin{itemize}
\item \textbf{Valeur} : théorique (l'homme au-delà de la conscience), thérapeutique (cure par la parole), culturelle (littérature, cinéma, éducation, mythes).
\item \textbf{Limites} : irréfutabilité (Popper), généralisations hâtives, interprétations arbitraires, dissolution de la liberté (Sartre), concurrence des neurosciences.
\item \textbf{Bilan} : la psychanalyse est une herméneutique féconde, non une science expérimentale ; l'homme n'est ni pure conscience ni pur jouet de l'inconscient.
\end{itemize}
\end{aretenir}

\coursec{Synthèse, méthode et entraînement au Bac}

Nous avons vu, dans la section précédente, que la psychanalyse freudienne possède une valeur indéniable — elle élargit notre connaissance de la vie psychique et soigne certaines souffrances — mais aussi des limites réelles : ses hypothèses sont difficilement vérifiables et ses interprétations parfois arbitraires. Il est temps maintenant de rassembler tous les fils de la leçon, de répondre clairement à notre problématique de départ, et surtout d'apprendre à mobiliser ces savoirs dans les deux épreuves du Baccalauréat technique : la dissertation et l'explication de texte.

\courssub{Carte mentale de synthèse}

La leçon entière tient en une image. La conscience se définit par sa nature (connaissance immédiate de soi et du monde), sa valeur (fondement de la liberté, de la morale et de la connaissance) et ses limites (elle peut se tromper, elle n'atteint pas les motifs profonds de nos actes). Face à elle, l'inconscient se présente comme une hypothèse (Freud) qui se manifeste dans les rêves, les lapsus, les actes manqués et les symptômes. La psychanalyse, enfin, est la méthode qui prétend explorer cet inconscient : elle a une valeur thérapeutique et herméneutique, mais des limites scientifiques.

\begin{popfigure}
\begin{center}
\begin{tikzpicture}[mindmap, grow cyclic,
  every node/.style={concept, font=\small},
  root concept/.append style={concept color=popblue, font=\bfseries},
  level 1/.append style={level distance=3.4cm, sibling angle=60},
  level 2/.append style={level distance=2.2cm, sibling angle=45, font=\scriptsize}]
\node[root concept]{Conscience et inconscient}
  child[concept color=popgreen] { node {Nature de la conscience}
    child { node {Se connaître soi-même} }
    child { node {Connaître le monde} }
  }
  child[concept color=popgreen!70] { node {Valeur de la conscience}
    child { node {Liberté} }
    child { node {Morale} }
    child { node {Connaissance} }
  }
  child[concept color=poporange] { node {Limites de la conscience}
    child { node {Illusions} }
    child { node {Motifs cachés} }
  }
  child[concept color=poppurple] { node {Hypothèse de l'inconscient}
    child { node {Freud} }
    child { node {Pulsions refoulées} }
  }
  child[concept color=poppurple!70] { node {Manifestations}
    child { node {Rêves} }
    child { node {Lapsus} }
    child { node {Actes manqués} }
  }
  child[concept color=popgold] { node {Psychanalyse}
    child { node {Valeur : cure, sens} }
    child { node {Limites : invérifiable} }
  };
\end{tikzpicture}
\end{center}
\legende{Carte mentale de la leçon : six branches relient la conscience (nature, valeur, limites), l'inconscient (hypothèse, manifestations) et la psychanalyse (valeur, limites).}
\end{popfigure}

\begin{aretenir}[L'essentiel de la leçon]
La conscience est pouvoir de se représenter soi-même et le monde : elle fonde la liberté et la responsabilité morale. Mais elle n'est pas toute la vie psychique : l'inconscient freudien, fait de désirs refoulés, agit en nous à notre insu et se trahit dans les rêves, les lapsus et les actes manqués. La psychanalyse propose d'explorer ce continent obscur, avec une réelle fécondité, mais sans atteindre la rigueur d'une science expérimentale.
\end{aretenir}

\courssub{Réponse argumentée à la problématique}

Notre problématique de départ était : la conscience suffit-elle à rendre compte de la vie psychique et de ce qu'est l'homme ? Nous pouvons maintenant y répondre de façon nuancée, en trois temps.

\textbf{Premier temps.} La conscience est bien le fondement de la liberté et de la morale. Sans conscience de soi, pas de délibération : l'élève qui hésite entre tricher et travailler pèse des raisons, et c'est cette transparence à soi qui le rend responsable de son choix. Descartes avait raison : le \emph{je pense} est le premier solide sur lequel tout s'appuie.

\textbf{Deuxième temps.} Mais la conscience n'est pas toute la vie psychique. L'observation des rêves, des lapsus, des oublis significatifs montre que des contenus psychiques agissent sans être connus. Freud compare la conscience à la partie émergée d'un iceberg : la plus grande partie de la vie mentale reste sous l'eau.

\textbf{Troisième temps.} L'inconscient ne supprime pas la conscience : il l'oblige à une connaissance de soi plus humble et plus profonde. Je ne suis pas transparent à moi-même ; me connaître devient un travail, non une évidence. Mais ce travail reste possible : prendre conscience de ses motivations cachées, c'est précisément élargir le domaine de la liberté. L'inconscient ne détrône pas la conscience, il la met au défi de grandir.

\courssub{Méthode de la dissertation}

La dissertation de philosophie n'est pas une récitation du cours : c'est une réponse argumentée et personnelle à une question précise.

\begin{methode}[Rédiger une introduction en cinq étapes]
\begin{enumerate}
\item \textbf{Amener le sujet} par une situation concrète ou une constatation générale (exemple : « Chacun croit se connaître lui-même... »).
\item \textbf{Définir les termes} importants du sujet (conscience, inconscient, responsabilité...).
\item \textbf{Dégager la contradiction} ou la difficulté que le sujet contient (c'est le cœur de l'introduction).
\item \textbf{Poser la problématique} sous forme d'une question claire.
\item \textbf{Annoncer le plan} : les deux ou trois grandes étapes du devoir.
\end{enumerate}
\end{methode}

\begin{popfigure}
\begin{center}
\begin{tikzpicture}[node distance=0.35cm,
  etape/.style={draw=popblue, fill=popblueL, rounded corners, text width=10.5cm, align=left, font=\small, inner sep=5pt},
  fleche/.style={-{Stealth[length=2.5mm]}, thick, popblue}]
\node[etape] (e1) {\textbf{1. Accroche} : situation concrète ou idée reçue qui amène le sujet};
\node[etape, below=of e1] (e2) {\textbf{2. Analyse des termes} : définitions provisoires des mots-clés};
\node[etape, below=of e2] (e3) {\textbf{3. Difficulté} : la contradiction ou le problème caché dans le sujet};
\node[etape, below=of e3] (e4) {\textbf{4. Problématique} : la question directrice, formulée en une phrase};
\node[etape, below=of e4] (e5) {\textbf{5. Annonce du plan} : « Nous verrons d'abord..., puis..., enfin... »};
\draw[fleche] (e1) -- (e2); \draw[fleche] (e2) -- (e3);
\draw[fleche] (e3) -- (e4); \draw[fleche] (e4) -- (e5);
\end{tikzpicture}
\end{center}
\legende{Schéma-méthode : les cinq étapes de l'introduction de dissertation, de l'accroche à l'annonce du plan.}
\end{popfigure}

\begin{methode}[Justifier une conclusion en trois gestes]
\begin{enumerate}
\item \textbf{Rappeler le chemin parcouru} : résumer en deux phrases les étapes du devoir.
\item \textbf{Trancher avec nuance} : répondre clairement à la problématique, en montrant la part de vérité de chaque thèse examinée.
\item \textbf{Ouvrir sur une question nouvelle} : prolonger la réflexion (exemple : « Reste à savoir si les neurosciences confirmeront l'hypothèse de l'inconscient. »).
\end{enumerate}
\end{methode}

\begin{attention}[L'erreur qui coûte le plus cher]
Réciter le cours sans répondre au sujet posé. Le correcteur ne demande pas « tout ce que tu sais sur la conscience », mais une réponse à \emph{cette} question. Chaque paragraphe doit faire avancer la réponse à la problématique : avant d'écrire un paragraphe, demande-toi : « En quoi ceci répond-il à la question du sujet ? »
\end{attention}

\begin{exemplebox}[Ce que vaut une bonne problématique au Bac]
Au Baccalauréat technique, la dissertation est notée sur 20. Les correcteurs valorisent en priorité quatre qualités : une introduction complète (accroche, définitions, problématique, annonce du plan), un plan clair et progressif, une argumentation appuyée sur des références exactes, et une langue soignée. Une problématique claire et un plan annoncé valent donc à eux seuls plusieurs points : c'est l'investissement le plus rentable du devoir.
\end{exemplebox}

\begin{savaistu}[Un classique des sessions camerounaises]
Les sujets portant sur la conscience et l'inconscient figurent parmi les plus fréquents des sessions de philosophie au Cameroun depuis plus de vingt ans, tant au Probatoire qu'au Baccalauréat. « La conscience suffit-elle à définir l'homme ? », « Peut-on tout savoir de soi-même ? », « L'inconscient excuse-t-il ? » reviennent régulièrement sous des formulations voisines. Maîtriser cette leçon, c'est se préparer à l'un des sujets les plus probables de l'examen.
\end{savaistu}

\courssub{Sujet type 1 : construction guidée du plan}

\begin{exoresolu}[Sujet : « La conscience suffit-elle à définir l'homme ? »]
Construire le plan détaillé de cette dissertation.
\tcblower
\textbf{Analyse du sujet.} « Suffire » signifie : être à elle seule suffisant. La question demande donc si la conscience épuise ce que l'homme est, ou s'il y a en l'homme autre chose que la conscience.

\textbf{Problématique.} Si l'homme se définit par la conscience (penser, se connaître, être libre), comment expliquer qu'il se conduise souvent sans savoir pourquoi ? La conscience suffit-elle à définir l'homme, ou faut-il admettre en lui une part obscure ?

\textbf{Plan dialectique en trois parties.}
\begin{itemize}
\item \textbf{I. Thèse : la conscience définit l'homme.} Descartes : le \emph{cogito}, « je pense, donc je suis » ; la conscience rend l'homme libre, responsable et capable de connaissance. Exemple : l'élève qui délibère avant un choix d'orientation.
\item \textbf{II. Antithèse : la conscience ne suffit pas.} Freud : l'inconscient, ses manifestations (rêves, lapsus, actes manqués) ; l'homme n'est pas transparent à lui-même. Exemple : oublier le nom d'une personne que l'on n'aime pas.
\item \textbf{III. Synthèse : l'homme est un être de conscience traversé par l'inconscient.} La conscience reste le sommet de l'humain, mais la connaissance de soi devient un travail ; la psychanalyse élargit la liberté au lieu de la nier.
\end{itemize}
\textbf{Conclusion possible.} La conscience définit l'homme comme être libre, mais non comme être transparent : se connaître soi-même est une tâche infinie.
\end{exoresolu}

\courssub{Sujet type 2 : rédaction guidée de l'introduction}

\begin{exoresolu}[Sujet : « L'inconscient nous délivre-t-il de la responsabilité ? »]
Rédiger l'introduction complète de cette dissertation.
\tcblower
\textbf{Accroche.} Devant le juge, l'accusé déclare parfois : « Je ne sais pas ce qui m'a pris. » De même, l'élève surpris en train de tricher invoque un geste qu'il « n'a pas voulu ».

\textbf{Définitions.} L'inconscient désigne, chez Freud, l'ensemble des contenus psychiques (désirs, souvenirs refoulés) qui agissent en nous sans que nous en ayons conscience. La responsabilité est le fait de devoir répondre de ses actes, parce qu'on les a voulus en connaissance de cause.

\textbf{Difficulté.} Si nos actes sont déterminés par des forces inconscientes, nous ne les choisissons pas vraiment : comment alors nous en tenir pour responsables ? Mais si l'inconscient excuse tout, toute morale et toute justice deviennent impossibles.

\textbf{Problématique.} La découverte de l'inconscient supprime-t-elle la responsabilité humaine, ou au contraire la déplace-t-elle et l'approfondit-elle ?

\textbf{Annonce du plan.} Nous verrons d'abord que l'inconscient semble détruire la responsabilité, puis qu'il ne saurait tout excuser, enfin qu'il invite à une responsabilité nouvelle : celle de se connaître soi-même.
\end{exoresolu}

\courssub{Explication de texte : Freud, « Le moi n'est pas maître dans sa propre maison »}

\begin{methode}[Lire un texte philosophique en quatre questions]
\begin{enumerate}
\item \textbf{Thèse} : que soutient l'auteur, en une phrase ?
\item \textbf{Arguments} : comment le démontre-t-il (exemples, raisonnements, comparaisons) ?
\item \textbf{Portée} : quelles conséquences pour notre conception de l'homme ?
\item \textbf{Critique} : la thèse est-elle solide ? Quelles objections peut-on lui opposer ?
\end{enumerate}
\end{methode}

\textbf{Le texte.} Freud écrit que l'humanité a subi trois grandes blessures narcissiques : Copernic a montré que la Terre n'est pas le centre de l'univers ; Darwin, que l'homme descend de l'animal ; la psychanalyse, enfin, montre que « le moi n'est pas maître dans sa propre maison ».

\textbf{Thèse.} L'homme n'est pas transparent à lui-même : une partie de sa vie psychique lui échappe, et la conscience n'est pas souveraine dans l'esprit.

\textbf{Arguments.} Freud s'appuie sur les manifestations de l'inconscient (rêves, lapsus, symptômes) qui prouvent que des contenus psychiques agissent sans notre savoir. La comparaison avec Copernic et Darwin donne à sa découverte une portée historique : comme eux, il décentre l'homme de sa prétendue position de maître.

\textbf{Portée.} L'orgueil humain en prend un coup : je ne suis pas totalement ce que je crois être. Mais cette humiliation est aussi un progrès de la connaissance de soi.

\textbf{Critique.} On peut objecter que l'inconscient est une hypothèse invérifiable (on ne peut pas l'observer directement), et que reconnaître des forces inconscientes ne supprime pas la maîtrise : prendre conscience de ses déterminations, c'est précisément commencer à les gouverner. Le moi n'est peut-être pas maître absolu, mais il peut devenir un meilleur locataire de sa maison.

\courssub{Application aux situations concrètes}

\begin{exemplebox}[Trois situations de la vie courante]
\begin{itemize}
\item \textbf{La responsabilité de l'élève tricheur.} Un élève copie pendant l'examen et prétend qu'il « n'a pas pu s'en empêcher ». L'inconscient peut expliquer son geste (angoisse, peur de décevoir ses parents), mais ne l'excuse pas : il reste capable de réfléchir et de choisir. La sanction vise justement à réveiller cette conscience.
\item \textbf{Traditions et groupe dans les choix individuels.} Un jeune choisit un métier ou un conjoint sous la pression de la famille et du village, en croyant choisir librement. Freud et la psychologie sociale montrent que des désirs et des modèles intégrés sans examen orientent nos « choix ». Prendre conscience de ces influences, c'est rendre sa liberté plus réelle.
\item \textbf{Réseaux sociaux et image de soi.} Sur les réseaux sociaux, chacun construit une image idéalisée de lui-même. Ce décalage entre le moi réel et le moi affiché illustre le jeu du refoulement et du narcissisme : l'inconscient cherche à fuir ce qui déplaît. La conscience critique doit apprendre à distinguer l'image de la réalité.
\end{itemize}
\end{exemplebox}

\begin{exercice}[Pour s'entraîner]
Sujet de dissertation : « Peut-on tout savoir de soi-même ? » Rédiger l'introduction complète (cinq étapes) et le plan détaillé en trois parties.
\end{exercice}

\begin{corrige}[Exercice]
\textbf{Introduction.} Accroche : chacun croit être la personne qu'il connaît le mieux. Définitions : « soi-même » désigne la vie psychique propre ; « tout savoir » signifie une connaissance totale et transparente. Difficulté : je me connais de l'intérieur, et pourtant mes rêves, mes lapsus, mes oublis me surprennent. Problématique : la conscience de soi permet-elle une connaissance totale de soi, ou l'inconscient en fixe-t-il la limite ? Annonce du plan : nous verrons d'abord que la conscience semble tout pouvoir savoir de soi, puis que l'inconscient limite ce pouvoir, enfin que la connaissance de soi est un progrès sans fin.

\textbf{Plan.} I. Oui : la conscience est connaissance immédiate de soi (Descartes) ; la réflexion permet de progresser dans la connaissance de soi. II. Non : l'inconscient freudien échappe par principe à la conscience directe ; il ne se révèle que par des signes interprétés. III. Nuance : on ne peut pas tout savoir de soi, mais on peut toujours savoir davantage ; la connaissance de soi est un travail sans fin, condition d'une liberté plus grande.
\end{corrige}

\newpage

\coursec{Activité d'intégration}

\coursec{La conscience et l'inconscient}

\courssub{Objectifs de l'activité}

Cette activité d'intégration mobilise l'ensemble des savoirs de la leçon~\emph{La conscience et l'inconscient} dans une situation concrète de la vie courante. À l'issue de ce travail pratique philosophique, l'élève doit être capable de~:

\begin{itemize}
\item \cle{Appliquer} les notions de conscience, d'inconscient, d'acte manqué et de mauvaise foi à des conduites réelles et observables~;
\item \cle{Distinguer} un acte conscient et volontaire d'un acte influencé par le milieu social et d'une manifestation possible de l'inconscient~;
\item \cle{Construire} un tableau d'interprétation et un schéma (le diagramme «~iceberg~») pour organiser une analyse~;
\item \cle{Rédiger et justifier} une démarche argumentée~: une conclusion individuelle d'une quinzaine de lignes répondant à la question «~Sommes-nous responsables de tout ce que nous faisons~?~»~;
\item \cle{Confronter} son interprétation à celle des autres groupes lors d'une restitution orale, exercice directement préparatoire à la dissertation du Probatoire et du Baccalauréat technique.
\end{itemize}

\begin{aretenir}[Fil conducteur]
La question centrale du TP reprend la célèbre formule de Freud~: \emph{«~Le moi n'est pas maître dans sa propre maison.~»} Il s'agit de vérifier, sur des cas concrets, si cette thèse est fondée, et quelles conséquences elle a sur la \cle{responsabilité} morale et juridique de l'individu.
\end{aretenir}

\courssub{Situation}

\textbf{TP philosophique~: «~Le moi est-il maître dans sa propre maison~?~»}\par

Dans le cadre de la semaine philosophique du lycée technique de la ville, le professeur de Philosophie remet à chaque groupe de quatre élèves un \emph{dossier d'enquête} composé de quatre documents. La classe compte 40~élèves, répartis en 10~groupes~; chaque groupe dispose de 30~minutes d'analyse, puis la classe entière consacre 20~minutes à la restitution orale (2~minutes par groupe) et 10~minutes au débat final.

\textbf{Document 1~— Le commerçant de Douala.}\par
M.~Ekane tient une boutique de vivre frais au marché Mboppi à Douala. Un soir, après une journée fatigante, une cliente lui achète pour \num{2500}~FCFA de marchandises et paie avec un billet de \num{5000}~FCFA. M.~Ekane lui rend \num{1500}~FCFA au lieu de \num{2500}~FCFA. La cliente, distraite, s'éloigne. Interrogé le lendemain par la cliente revenue se plaindre, M.~Ekane jure qu'il n'a «~rien remarqué~». Or, ce même jour, il avait eu une vive dispute avec cette cliente au sujet d'une dette ancienne de \num{10000}~FCFA. Ses voisins de boutique affirment qu'il est «~honnête d'habitude~».

\textbf{Document 2~— Le rêve répétitif.}\par
Ariane, élève en Terminale~technique, rêve chaque nuit, depuis une semaine, qu'elle arrive en retard à l'épreuve de Philosophie~: la salle est vide, sa feuille reste blanche, et le correcteur est son père qui la regarde sans parler. Éveillée, Ariane affirme~: «~Je ne suis pas du tout stressée, tout va bien.~» Son carnet montre pourtant qu'elle a manqué trois devoirs ce trimestre et que son père lui a annoncé qu'il cesserait de payer ses études en cas d'échec au Baccalauréat.

\textbf{Document 3~— Freud, extrait de la \emph{Psychopathologie de la vie quotidienne} (1901).}\par
«~Les actes manqués — oublis, lapsus, erreurs — ne sont pas des accidents sans signification. Ils sont l'expression de désirs refoulés qui trouvent, malgré la censure de la conscience, un chemin détourné pour se manifester. L'oubli d'un rendez-vous désagréable, le lapsus qui trahit une pensée cachée, révèlent que le moi n'est pas maître dans sa propre maison.~»

\textbf{Document 4~— Sartre, extrait de \emph{L'Être et le Néant} (1943).}\par
«~La mauvaise foi consiste à se mentir à soi-même~: l'homme se cache sa propre liberté en se faisant passer pour une chose déterminée. Le garçon de café qui joue à être garçon de café, l'amoureux qui prétend «~ne pas savoir~» ce qu'il sait, fuient leur responsabilité. Nous sommes condamnés à être libres~: il n'y a pas d'inconscient qui décide à notre place.~»

\textbf{Travail attendu.} Chaque groupe doit~: (a) classer chaque conduite décrite dans les documents 1 et 2 selon trois catégories~— \emph{acte conscient et volontaire}, \emph{acte influencé par le milieu}, \emph{manifestation possible de l'inconscient}~; (b) construire un tableau d'interprétation croisant les quatre documents~; (c) dessiner et annoter un diagramme «~iceberg~» de la psyché~; (d) rédiger individuellement une conclusion argumentée d'environ quinze lignes répondant à la question~: \emph{«~Sommes-nous responsables de tout ce que nous faisons~?~»}

\courssub{Consignes}

\begin{enumerate}
\item \textbf{Analyse du document 1.} Relevez dans le récit de M.~Ekane trois indices qui permettent de douter que l'«~oubli~» soit purement accidentel. Ce geste peut-il être interprété comme un acte manqué au sens de Freud~? Justifiez avec le document 3.
\item \textbf{Analyse du document 2.} Pourquoi le discours éveillé d'Ariane («~je ne suis pas stressée~») entre-t-il en contradiction avec son rêve~? Que révèle, selon Freud, le contenu symbolique du rêve (retard, feuille blanche, père-correcteur)~?
\item \textbf{Confrontation des thèses.} À l'aide du document 4, formulez l'objection que Sartre adresserait à l'interprétation freudienne des deux cas. Que dirait Sartre de M.~Ekane et d'Ariane~?
\item \textbf{Classification.} Complétez le tableau d'interprétation à trois colonnes (acte conscient et volontaire / acte influencé par le milieu / manifestation possible de l'inconscient) en y plaçant chaque conduite avec sa justification.
\item \textbf{Schématisation.} Dessinez l'iceberg de la psyché~: partie émergée (conscience), ligne de flottaison (préconscient), partie immergée (inconscient). Annotez chaque zone avec un exemple tiré des documents.
\item \textbf{Synthèse rédigée.} Rédigez individuellement, en une quinzaine de lignes, une réponse argumentée à la question~: «~Sommes-nous responsables de tout ce que nous faisons~?~» Votre réponse doit comporter une thèse, un argument, un exemple tiré du dossier et une objection discutée.
\item \textbf{Restitution.} Présentez oralement en 2~minutes la position de votre groupe et écoutez celles des autres~: notez une convergence et une divergence.
\end{enumerate}

\courssub{Ressources à mobiliser}

\begin{aretenir}[Notions utiles]
\begin{itemize}
\item \textbf{Conscience}~: connaissance immédiate que le sujet a de lui-même et du monde~; chez Descartes, le \emph{cogito} en fait le fondement certain de toute connaissance. Sa \emph{valeur}~: elle rend possible la volonté, la liberté et la responsabilité. Ses \emph{limites}~: elle est sélective, faillible, trompée par les sens, les habitudes et le milieu social.
\item \textbf{Inconscient (Freud)}~: ensemble des désirs, souvenirs et conflits \cle{refoulés}, inaccessibles à la conscience, qui s'expriment de manière déguisée dans les \emph{rêves}, les \emph{lapsus} et les \emph{actes manqués}. Topique~: conscience / préconscient / inconscient.
\item \textbf{Psychanalyse}~: méthode d'exploration de l'inconscient (libre association, interprétation des rêves). \emph{Valeur}~: elle explique des conduites incompréhensibles et soigne certains troubles. \emph{Limites}~: ses interprétations sont difficilement vérifiables scientifiquement~; elle risque d'excuser toutes les fautes.
\item \textbf{Mauvaise foi (Sartre)}~: attitude de celui qui se ment à soi-même pour fuir sa liberté et sa responsabilité~; objection majeure à l'inconscient.
\item \textbf{Responsabilité}~: obligation de répondre de ses actes~; elle suppose la conscience et la liberté, mais le droit admet des circonstances atténuantes (trouble psychique, contrainte).
\end{itemize}
\end{aretenir}

\courssub{Démarche et corrigé commenté}

\begin{exoresolu}[Corrigé commenté du TP]
Voir le dossier d'enquête ci-dessus (Documents 1 à 4) et les consignes 1 à 6.
\tcblower
\textbf{Étape 1 — Analyse du document 1 (consigne 1).} Trois indices troublent la thèse de l'accident~: (i) l'erreur porte exactement sur la cliente avec qui M.~Ekane s'est disputé le jour même~; (ii) l'erreur lui est \emph{favorable} (il garde \num{1000}~FCFA)~; (iii) le motif existe (la dette de \num{10000}~FCFA et le ressentiment). Selon Freud (document 3), un oubli n'est jamais neutre~: il exprime un désir refoulé. Ici, le désir inconscient de «~se faire payer autrement~» la dette aurait détourné le calcul. C'est un \emph{acte manqué} typique~: la conscience croit calculer, l'inconscient règle ses comptes.

\textbf{Étape 2 — Analyse du document 2 (consigne 2).} Le discours éveillé d'Ariane («~tout va bien~») est contredit par son rêve répétitif. Pour Freud, le rêve est la «~voie royale~» vers l'inconscient~: il réalise de façon déguisée ce que la conscience refuse d'avouer. Le retard symbolise la peur de l'échec~; la feuille blanche, l'angoisse de l'impuissance~; le père-correcteur, la menace réelle (fin du financement des études) que la conscience d'Ariane refoule pour se protéger. Le rêve dit la vérité que la conscience tait.

\textbf{Étape 3 — Objection sartrienne (consigne 3).} Sartre répondrait que l'inconscient est une invention commode pour se décharger de sa liberté. M.~Ekane \emph{sait} au fond qu'il a mal rendu la monnaie~: son «~je n'ai rien remarqué~» serait de la mauvaise foi, un mensonge à soi-même. De même, Ariane choisit de se dire «~pas stressée~» pour ne pas affronter la situation~: elle est responsable de son déni. Pour Sartre, «~nous sommes condamnés à être libres~»~: aucune force cachée ne décide à notre place.

\textbf{Étape 4 — Classification (consigne 4).}
\begin{center}
\begin{tabularx}{\linewidth}{|l|X|X|}\hline
\rowcolor{popblueL} \textbf{Catégorie} & \textbf{Conduite} & \textbf{Justification} \\ \hline
Acte conscient et volontaire & Ariane affirme «~tout va bien~» & Parole délibérée, contrôlée (ou mauvaise foi selon Sartre) \\ \hline
Acte influencé par le milieu & Fatigue de M.~Ekane après le marché~; pression du père d'Ariane & Conditions sociales et matérielles qui pèsent sur la conscience \\ \hline
Manifestation possible de l'inconscient & Erreur de monnaie~; rêve répétitif & Acte manqué et rêve~: contenus refoulés qui ressortent déguisés \\ \hline
\end{tabularx}
\end{center}

\textbf{Étape 5 — Iceberg annoté (consigne 5).}
\begin{popfigure}
\begin{tikzpicture}[scale=0.85]
% Ligne de flottaison
\draw[popblue, thick] (-3.5,0) -- (3.5,0) node[right, popdark] {Niveau de la conscience};
% Partie émergée (Conscience)
\draw[popblue!60, fill=popblue!10, thick] (-2.5,0) .. controls (-3,0.8) and (-3,1.8) .. (0,2.2) .. controls (3,1.8) and (3,0.8) .. (2.5,0);
% Partie immergée (Inconscient)
\draw[popblue!60, fill=popblue!25, thick] (-2.5,0) .. controls (-3,-1) and (-3,-3) .. (0,-4) .. controls (3,-3) and (3,-1) .. (2.5,0);
% Étiquettes zones
\node[popdark, font=\bfseries\small] at (0, 1.3) {Conscience (émergée)};
\node[popdark, font=\bfseries\small] at (0, -0.4) {Préconscient (flottaison)};
\node[popdark, font=\bfseries\small] at (0, -2.5) {Inconscient (immergé)};
% Exemples côté gauche
\node[popdark, font=\small, align=center, text width=2.8cm] at (-3.2, 1.5) {Calculs, paroles volontaires\\ (Ariane~: «~tout va bien~»)}; 
\node[popdark, font=\small, align=center, text width=2.8cm] at (-3.2, -0.5) {Souvenirs accessibles\\ (Dispute M.~Ekane)}; 
\node[popdark, font=\small, align=center, text width=2.8cm] at (-3.2, -2.5) {Désirs refoulés\\ (Ressentiment M.~Ekane)}; 
% Exemples côté droit
\node[popdark, font=\small, align=center, text width=2.8cm] at (3.2, 1.5) {Perception immédiate\\ (Client, monnaie)}; 
\node[popdark, font=\small, align=center, text width=2.8cm] at (3.2, -0.5) {Connaissances latentes\\ (Règles de calcul)}; 
\node[popdark, font=\small, align=center, text width=2.8cm] at (3.2, -2.5) {Angoisses refoulées\\ (Échec Ariane, père)}; 
% Flèches de liaison
\draw[->, poporange, thick, shorten >=2pt] (-3.2, 1.5) -- (-1.8, 1.0);
\draw[->, poporange, thick, shorten >=2pt] (-3.2, -0.5) -- (-1.5, -0.3);
\draw[->, poporange, thick, shorten >=2pt] (-3.2, -2.5) -- (-1.2, -1.8);
\draw[->, poporange, thick, shorten >=2pt] (3.2, 1.5) -- (1.8, 1.0);
\draw[->, poporange, thick, shorten >=2pt] (3.2, -0.5) -- (1.5, -0.3);
\draw[->, poporange, thick, shorten >=2pt] (3.2, -2.5) -- (1.2, -1.8);
\end{tikzpicture}
\legende{Diagramme «~iceberg~» de la topique freudienne annoté d'après le dossier.}
\end{popfigure}

\textbf{Étape 6 — Modèle de conclusion rédigée (consigne 6).}\par
«~Sommes-nous responsables de tout ce que nous faisons~? Non, pas de \emph{tout}~: le dossier montre que certaines conduites échappent en partie à la volonté consciente. L'erreur de monnaie de M.~Ekane et le rêve d'Ariane s'interprètent, avec Freud, comme des manifestations de désirs et d'angoisses refoulés que le sujet ne contrôle pas directement. La conscience, sélective et faillible, n'est pas maîtresse absolue dans sa propre maison. Cependant, l'objection de Sartre interdit d'en conclure à l'irresponsabilité totale~: invoquer l'inconscient peut devenir une mauvaise foi, une ruse pour fuir sa liberté. M.~Ekane reste tenu de réparer son erreur une fois qu'elle est révélée, et Ariane doit affronter son angoisse au lieu de la nier. Nous concluons donc que l'inconscient \emph{atténue} la responsabilité sans l'\emph{abolir}~: être responsable, ce n'est pas avoir tout voulu, c'est assumer ce que l'on découvre de soi et y répondre.~»

\textbf{Justification de la démarche~:} la réponse suit le schéma exigé au Baccalauréat — thèse nuancée, argument freudien, exemple du dossier, objection sartrienne discutée, conclusion synthétique dépassant l'opposition.
\end{exoresolu}

\courssub{Bilan de l'activité}

Ce TP a permis de mettre en évidence trois acquis essentiels de la leçon.

Premièrement, la \cle{conscience} n'est ni transparente ni toute-puissante~: les documents 1 et 2 montrent concrètement qu'elle peut être trompée, fatiguée, influencée par le milieu, et que des contenus psychiques lui échappent. Les élèves ont ainsi vérifié sur des cas camerounais familiers les \emph{limites de la conscience} étudiées en cours.

Deuxièmement, l'\cle{inconscient} freudien offre un modèle d'interprétation puissant~: actes manqués et rêves cessent d'être des accidents absurdes pour devenir des messages déguisés. Mais la confrontation avec Sartre a montré les \emph{limites de la psychanalyse}~: ses interprétations sont difficiles à prouver et peuvent servir d'alibi.

Troisièmement, la question de la \cle{responsabilité} a obligé les élèves à dépasser le «~pour ou contre~» vers une position nuancée et justifiée~: l'inconscient explique sans forcément excuser. Sur le plan méthodologique, l'activité a entraîné les deux savoir-faire officiels de l'unité — appliquer les notions à des situations concrètes et rédiger une démarche argumentée — et a répété, en miniature, le geste de la dissertation~: analyser, confronter deux thèses d'auteurs, illustrer, conclure avec nuance.

\newpage

\coursec{Méthodes et exercices résolus}

\courssub{Fiches méthodes et exercices résolus}

\begin{methode}[Appliquer les notions à une situation concrète de la vie courante]
Pour traiter une situation concrète (un fait divers, une scène de la vie quotidienne, un cas professionnel) avec les notions de \cle{conscience} et d'\cle{inconscient} :
\begin{enumerate}
\item \textbf{Lire et décrire} la situation en une ou deux phrases, sans jugement : qui fait quoi, dans quelles circonstances ?
\item \textbf{Dégager le problème philosophique} : la situation met-elle en jeu la lucidité d'un sujet, un acte inexpliqué, un oubli, un rêve, une erreur de jugement ?
\item \textbf{Choisir la notion pertinente} : \cle{conscience} (présence à soi, \cle{intentionnalité}, volonté) ou \cle{inconscient} (pulsion, \cle{refoulement}, acte manqué, rêve).
\item \textbf{Mobiliser une thèse d'auteur} : Descartes (la conscience comme évidence du \emph{cogito}), Bergson (la conscience comme mémoire et durée), Freud (l'inconscient comme cause des symptômes), Sartre (la \cle{mauvaise foi}).
\item \textbf{Appliquer rigoureusement} : montrer en quoi la thèse éclaire la situation, puis discuter ses limites (la notion suffit-elle à tout expliquer ?).
\item \textbf{Conclure} par une réponse nuancée et justifiée, en revenant explicitement à la situation de départ.
\end{enumerate}
\textbf{Formule à retenir} : Situation $\rightarrow$ Problème $\rightarrow$ Notion $\rightarrow$ Thèse $\rightarrow$ Application $\rightarrow$ Conclusion nuancée.
\end{methode}

\begin{exoresolu}[Application : le lapsus au marché]
\textbf{Énoncé.} À Yaoundé, au marché Mokolo, une vendeuse accueille un client qu'elle soupçonne de marchander excessivement. Voulant dire « Bienvenue, monsieur », elle dit : « Malvenue, monsieur ». En vous appuyant sur la notion d'inconscient, analysez cette situation et montrez ce qu'elle révèle de la vie psychique.
\tcblower
\textbf{Solution détaillée.}
\begin{enumerate}
\item \emph{Description.} La vendeuse prononce un mot contraire à son intention consciente : elle voulait être polie, elle a été insultante. Il s'agit d'un \cle{lapsus}, c'est-à-dire d'un acte manqué de la parole.
\item \emph{Problème.} D'où vient cet écart entre l'intention consciente et l'acte réel ? La conscience est-elle maîtresse de nos paroles ?
\item \emph{Notion et thèse.} Freud, dans la \emph{Psychopathologie de la vie quotidienne} (1901), soutient que le lapsus n'est pas un hasard : il est la manifestation d'un désir refoulé. La vendeuse éprouve une hostilité réelle envers ce client, mais les convenances sociales la contraignent à la politesse ; le désir refoulé s'exprime alors déguisé, par la faute de langage.
\item \emph{Application.} Le mot « malvenue » trahit le contenu inconscient : l'irritation. Le lapsus fonctionne comme un compromis entre la pulsion (chasser le client) et la censure morale (rester courtoise). Il prouve que la conscience n'est pas transparente à elle-même : « le moi n'est pas maître dans sa propre maison », écrit Freud.
\item \emph{Discussion des limites.} On peut objecter que certains lapsus relèvent de la fatigue ou de la distraction purement mécaniques, sans signification psychique. La psychanalyse risque ici de surinterpréter.
\item \emph{Conclusion.} La situation illustre la thèse freudienne : nos actes apparemment insignifiants ont un sens caché, et l'inconscient limite le pouvoir de la conscience. Toutefois, toute erreur de parole n'est pas nécessairement un message de l'inconscient.
\end{enumerate}
\end{exoresolu}

\begin{methode}[Construire une problématique à partir d'un sujet de dissertation]
\begin{enumerate}
\item \textbf{Repérer les termes clés} du sujet et les définir séparément (\cle{conscience} : connaissance immédiate de soi et du monde ; \cle{inconscient} : ensemble des contenus psychiques échappant à la conscience).
\item \textbf{Identifier l'opposition ou la tension} entre les termes : ici, la conscience prétend à la \cle{transparence}, l'inconscient affirme l'\cle{opacité}.
\item \textbf{Formuler la question centrale} sous forme interrogative : la conscience se suffit-elle à elle-même ? Suis-je ce que je sais de moi ?
\item \textbf{Prévoir le plan dialectique} : thèse (la conscience est souveraine — Descartes), antithèse (l'inconscient gouverne — Freud), synthèse (une conscience lucide mais limitée).
\end{enumerate}
\textbf{Réflexe} : la problématique doit tenir en une phrase interrogative annoncée dès l'introduction.
\end{methode}

\begin{exoresolu}[Problématiser un sujet]
\textbf{Énoncé.} Sujet : « La conscience de soi suffit-elle à nous connaître ? » Dégagez la problématique et annoncez un plan en trois parties.
\tcblower
\textbf{Solution détaillée.}
\begin{enumerate}
\item \emph{Analyse des termes.} « Conscience de soi » : capacité réflexive par laquelle le sujet se saisit comme pensant et agissant. « Suffire » : être une condition nécessaire et complète. « Se connaître » : accéder à la vérité de soi-même.
\item \emph{Tension.} Le présupposé commun est que je suis l'être que je connais le mieux (évidence cartésienne du \emph{cogito}). Pourtant, l'expérience montre des zones d'ombre : oublis, rêves, actes impulsifs, maladies nerveuses.
\item \emph{Problématique.} « Si la conscience me donne un accès immédiat à moi-même, comment expliquer que je puisse me méconnaître ? La conscience de soi épuise-t-elle la réalité du sujet, ou faut-il admettre un inconscient qui en fixe les limites ? »
\item \emph{Plan.} I. La conscience comme connaissance de soi évidente et fondatrice (Descartes, \emph{Méditations métaphysiques} : je pense, donc je suis). II. Les limites de la conscience et l'hypothèse de l'inconscient (Freud : le refoulement, le rêve comme « voie royale » vers l'inconscient). III. Synthèse : la conscience reste précieuse comme lucidité morale, mais la connaissance de soi exige un travail d'interprétation (Ricoeur : le soi se comprend à travers des signes).
\end{enumerate}
\textbf{Conclusion.} La problématique oppose la transparence cartésienne à l'opacité freudienne et ouvre sur une connaissance de soi conquise, non donnée.
\end{exoresolu}

\begin{methode}[Rédiger et justifier une démarche argumentative]
Pour rédiger un développement convaincant :
\begin{enumerate}
\item \textbf{Annoncer} l'idée de la partie en une phrase claire (la « phrase de thèse »).
\item \textbf{Définir} les notions mobilisées dès leur première apparition.
\item \textbf{Exposer la thèse d'un auteur} : nom, œuvre, argument central, éventuellement une courte citation exacte.
\item \textbf{Illustrer} par un exemple concret, de préférence proche de la vie des élèves (école, famille, marché, atelier).
\item \textbf{Discuter} : présenter une objection, puis y répondre ; c'est ce qui distingue une dissertation d'un résumé de cours.
\item \textbf{Conclure la partie} et \textbf{transitionner} vers la suivante par une question.
\end{enumerate}
\textbf{Formule à retenir} : Idée $\rightarrow$ Définition $\rightarrow$ Auteur $\rightarrow$ Exemple $\rightarrow$ Objection/Réponse $\rightarrow$ Bilan $\rightarrow$ Transition.
\end{methode}

\begin{exoresolu}[Rédiger un paragraphe argumenté]
\textbf{Énoncé.} Rédigez un paragraphe (environ 15 lignes) justifiant la thèse suivante : « La conscience a une valeur morale : elle rend l'homme responsable de ses actes. »
\tcblower
\textbf{Solution détaillée (modèle de rédaction).}

La conscience possède d'abord une \cle{valeur morale}, car elle fonde la \cle{responsabilité}. Par conscience morale, on entend la capacité du sujet à juger ses propres actes à la lumière du bien et du mal. Kant, dans la \emph{Critique de la raison pratique}, montre que la conscience du devoir — « la loi morale en moi » — distingue l'homme de l'animal : seul un être conscient peut se donner une règle et s'y tenir librement. Ainsi, un élève de Terminale qui, tenté de tricher au Probatoire, renonce parce qu'il s'entend dire « ceci est mal », éprouve cette voix intérieure qu'est la conscience morale. Sans elle, aucun tribunal, aucune sanction scolaire n'aurait de sens : on ne punit que celui qu'on suppose avoir su ce qu'il faisait. On objectera que certains criminels agissent sans remords : la conscience serait donc inégalement répartie. Mais l'absence de remords ne supprime pas la capacité de juger ; elle signale seulement que la conscience morale peut être étouffée par l'habitude ou la passion, ce que Socrate résumait : « nul ne fait le mal volontairement ». La conscience est donc bien la condition de la responsabilité, même si elle peut être défaillante. Reste à savoir si cette conscience est toujours lucide, ou si des forces cachées la trompent : c'est la question de l'inconscient.

\emph{Justification de la démarche} : le paragraphe suit les sept étapes — phrase de thèse, définition, auteur (Kant), exemple camerounais contextualisé (triche à l'examen), objection (le criminel sans remords), réponse (Socrate), bilan et transition vers l'inconscient.
\end{exoresolu}

\begin{methode}[Expliquer et discuter un texte (commentaire)]
\begin{enumerate}
\item \textbf{Lire trois fois} le texte ; souligner les termes techniques et les articulations logiques (« donc », « mais », « en effet »).
\item \textbf{Identifier l'auteur et la thèse} : quelle position le texte défend-il ?
\item \textbf{Dégager la structure} : diviser le texte en deux ou trois moments et les formuler en phrases simples.
\item \textbf{Expliquer} chaque moment : définir les notions, reconstruire le raisonnement, donner un exemple personnel.
\item \textbf{Discuter} : évaluer la portée et les limites de la thèse (est-elle vraie ? universelle ? vérifiable ?).
\item \textbf{Conclure} : quel intérêt philosophique présente ce texte pour le problème de la conscience ?
\end{enumerate}
\textbf{Attention} : ne jamais paraphraser ligne par ligne ; expliquer, c'est reconstruire l'argumentation, pas répéter le texte.
\end{methode}

\begin{exoresolu}[Commentaire d'un extrait de Freud]
\textbf{Énoncé.} « Le moi n'est pas maître dans sa propre maison » (Freud). Expliquez cette formule et discutez sa portée.
\tcblower
\textbf{Solution détaillée.}
\begin{enumerate}
\item \emph{Compréhension.} Le « \cle{moi} » désigne la conscience, le sujet qui se dit « je ». La « maison » est la psyché tout entière. Freud affirme donc que la conscience ne gouverne pas la totalité de la vie psychique.
\item \emph{Contexte et thèse.} Freud présente la psychanalyse comme la troisième « blessure narcissique » infligée à l'homme : après Copernic (la Terre n'est pas le centre du monde) et Darwin (l'homme descend du singe), Freud montre que le moi n'est pas maître chez lui, car une grande partie de nos pensées, désirs et souvenirs est inconsciente.
\item \emph{Explicitation.} L'appareil psychique se compose du \cle{ça} (pulsions), du \cle{moi} (instance de réalité) et du \cle{surmoi} (interdits intériorisés). Le moi est tiraillé entre ces forces qu'il ne contrôle pas. Preuves cliniques : les rêves, les actes manqués, les symptômes névrotiques, qui disent ce que la conscience ignore.
\item \emph{Exemple.} Un apprenti menuisier d'Obala qui oublie systématiquement son rendez-vous chez un client qu'il craint ne subit pas un simple hasard : son oubli exprime un refoulement, une peur que sa conscience refuse d'assumer.
\item \emph{Discussion.} Portée : la formule détrône la conscience absolue de Descartes et invite à l'humilité — se connaître devient un travail, non une évidence. Limites : la notion d'inconscient est difficilement réfutable au sens scientifique (Popper reproche à la psychanalyse d'expliquer tout et son contraire) ; de plus, nier toute maîtrise au moi menace la responsabilité morale.
\item \emph{Conclusion.} La formule de Freud a une valeur critique décisive : elle oblige la philosophie à penser un sujet divisé. Mais elle doit être tempérée : reconnaître l'inconscient, ce n'est pas abolir la conscience, c'est lui assigner la tâche de s'élucider elle-même.
\end{enumerate}
\end{exoresolu}

\begin{methode}[Évaluer la valeur et les limites d'une théorie (la psychanalyse)]
Pour juger équilibrément une doctrine comme la \cle{psychanalyse} :
\begin{enumerate}
\item \textbf{Rappeler la thèse} en deux lignes : existence d'un inconscient, méthode d'investigation (association libre, interprétation des rêves), visée thérapeutique (faire advenir le refoulé à la conscience).
\item \textbf{Présenter ses apports} : découverte du sens caché des symptômes ; nouvelle image de l'homme ; efficacité thérapeutique dans certaines névroses ; influence sur la littérature et la culture.
\item \textbf{Présenter ses limites} : caractère non falsifiable des interprétations (Popper) ; difficulté de vérification expérimentale ; risque de tout ramener à la sexualité ; durée et coût des cures.
\item \textbf{Comparer} avec une approche rivale (psychologie expérimentale, neurosciences) en rappelant que, dans le cadre du programme, la conscience reste le point de départ de toute enquête psychologique.
\item \textbf{Porter un jugement personnel argumenté} : la psychanalyse est-elle une science, une \cle{herméneutique} de l'existence, ou les deux ?
\end{enumerate}
\end{methode}

\begin{exoresolu}[Dissertation courte : valeur et limites de la psychanalyse]
\textbf{Énoncé.} « La psychanalyse est-elle une science ? » Rédigez une introduction et l'annonce du plan.
\tcblower
\textbf{Solution détaillée.}

\emph{Accroche.} Depuis un siècle, la psychanalyse de Freud prétend guérir les souffrances de l'âme en explorant l'inconscient. \emph{Définitions.} La psychanalyse est à la fois une théorie de l'appareil psychique (inconscient, refoulement, pulsions) et une méthode thérapeutique fondée sur la parole ; une science est un savoir fondé sur l'observation, l'expérimentation et la vérification. \emph{Tension.} Or la psychanalyse repose sur des interprétations de cas singuliers, difficilement soumises à l'expérience : peut-on encore l'appeler science ? \emph{Problématique.} La psychanalyse satisfait-elle aux critères de la scientificité, ou sa valeur est-elle ailleurs, dans une herméneutique de l'existence humaine ? \emph{Annonce du plan.} Nous verrons d'abord que la psychanalyse se présente comme une science de l'inconscient, fondée sur la clinique ; nous montrerons ensuite qu'elle échappe aux critères de la science expérimentale (irréfutabilité dénoncée par Popper) ; nous soutiendrons enfin que sa valeur propre est celle d'une méthode d'interprétation du sens, irremplaçable pour comprendre le sujet humain, sans être une science au sens strict.
\end{exoresolu}

\begin{methode}[Construire une conclusion de dissertation]
\begin{enumerate}
\item \textbf{Répondre explicitement} à la question posée dans la problématique (oui, non, ou réponse nuancée — mais tranchée).
\item \textbf{Résumer le chemin parcouru} en une phrase par partie, sans répéter les exemples.
\item \textbf{Justifier la nuance} : montrer que la synthèse dépasse l'opposition initiale (conscience/inconscient ne s'excluent pas, ils se complètent).
\item \textbf{Ouvrir} sur une question voisine du programme (la liberté, le sujet, le temps) en une seule phrase, sans développer.
\end{enumerate}
\textbf{Interdit} : introduire une notion nouvelle, un nouvel auteur ou un nouvel exemple dans la conclusion.
\end{methode}

\begin{exoresolu}[Rédiger une conclusion]
\textbf{Énoncé.} Rédigez la conclusion de la dissertation : « Suis-je ce que j'ai conscience d'être ? »
\tcblower
\textbf{Solution détaillée.}

Nous avons d'abord reconnu, avec Descartes, que la conscience est le fondement indubitable de mon existence : je ne peux douter que je pense. Mais l'examen des rêves, des lapsus et des symptômes a montré, avec Freud, que cette conscience est traversée par un inconscient qui pense et désire à mon insu : je suis donc plus que ce que j'ai conscience d'être. La synthèse s'impose alors : je ne suis ni pure transparence ni pure opacité ; je suis un sujet divisé, dont la conscience a pour tâche d'élucider progressivement les profondeurs. Se connaître soi-même n'est pas un donné, c'est un travail — celui de toute une vie. Cette conquête de soi par soi n'est-elle pas, au fond, la condition même de la liberté ?
\end{exoresolu}

\begin{methode}[Utiliser correctement les auteurs et les exemples]
\begin{enumerate}
\item \textbf{Citer brièvement} : nom de l'auteur, thèse en une phrase, éventuellement une citation courte et exacte (« Je pense, donc je suis » de Descartes ; « le rêve est la voie royale vers l'inconscient » de Freud).
\item \textbf{Ne jamais} écrire « selon les philosophes » ou « un auteur a dit » : chaque thèse doit être attribuée nommément.
\item \textbf{Choisir des exemples simples et vécus} : le rêve d'un examen la veille du Bac, l'oubli d'un rendez-vous redouté, le remords après une faute, la mauvaise foi d'un élève qui se trouve des excuses.
\item \textbf{Faire travailler l'exemple} : après l'avoir raconté en deux lignes, expliciter ce qu'il démontre (« cet exemple montre que... »).
\item \textbf{Éviter l'accumulation} : un auteur et un exemple par argument suffisent ; la qualité de l'analyse prime sur la quantité de références.
\end{enumerate}
\end{methode}

\begin{exoresolu}[Corriger un paragraphe défectueux]
\textbf{Énoncé.} Un élève écrit : « Les philosophes ont dit que la conscience est très importante. Par exemple mon ami rêve beaucoup. Donc l'inconscient existe. » Relevez les défauts de ce paragraphe et proposez une version corrigée.
\tcblower
\textbf{Solution détaillée.}
\emph{Défauts relevés} : (1) « les philosophes » : aucune attribution précise ; (2) « très importante » : jugement vague, sans définition de la conscience ; (3) l'exemple est raconté sans être analysé : en quoi rêver « beaucoup » prouve-t-il l'inconscient ? (4) le « donc » prétend conclure sans raisonnement ; (5) confusion : l'importance de la conscience ne prouve pas l'existence de l'inconscient.

\emph{Version corrigée.} La conscience, définie comme la connaissance immédiate que le sujet a de lui-même, semble d'abord épuiser la réalité psychique : c'est la thèse de Descartes, pour qui « je pense, donc je suis » fonde toute certitude. Pourtant, l'expérience du rêve contredit cette transparence. Ainsi, un élève rêve la veille du Baccalauréat qu'il arrive en retard à l'épreuve : ce rêve, dont le contenu échappe à sa volonté, exprime selon Freud une angoisse refoulée que la conscience diurne refusait d'avouer. Le rêve est alors « la voie royale vers l'inconscient » : il révèle que la psyché contient des contenus que le moi ignore. L'existence de l'inconscient s'impose donc comme hypothèse explicative, sans pour autant abolir la conscience, qui demeure le lieu de la lucidité et de la décision.
\end{exoresolu}

\begin{attention}[Les 5 erreurs qui coûtent des points au Bac]
\begin{enumerate}
\item \textbf{Confondre \cle{inconscient} et \cle{non-conscient}.} Le non-conscient (fonctions vitales automatiques : respiration, digestion) n'est pas l'inconscient freudien, qui est un psychisme actif fait de désirs refoulés. Employer « inconscient » pour dire « inattentif » est une faute de notion.
\item \textbf{Présenter Freud comme niant la conscience.} Freud ne supprime pas la conscience : il montre qu'elle n'est qu'une partie de la psyché. Écrire « pour Freud, la conscience n'existe pas » est un contresens éliminatoire.
\item \textbf{Affirmer sans argumenter.} « La conscience a des limites » n'est pas un argument : il faut une preuve (rêve, lapsus, acte manqué, maladie nerveuse) et un auteur. Toute affirmation non justifiée est sanctionnée.
\item \textbf{Paraphraser le texte au lieu de l'expliquer.} Dans le commentaire, recopier les phrases de l'auteur avec d'autres mots ne vaut rien : il faut reconstruire le raisonnement, définir les termes et discuter la thèse.
\item \textbf{Négliger la conclusion et la réponse à la problématique.} Beaucoup de copies s'arrêtent sans répondre à la question posée, ou introduisent un nouvel auteur dans la conclusion. La conclusion doit trancher, résumer le chemin parcouru, et ne contenir aucune notion nouvelle.
\end{enumerate}
\end{attention}

\newpage

\coursec{Exercices}

\coursec{Leçon 04 : La conscience et l'inconscient}

\courssub{Problématique et définitions}

\begin{definition}[Conscience]
Du latin \emph{conscientia} (\emph{cum-scire}, « savoir avec »). La conscience est la connaissance immédiate et intuitive que le sujet a de ses propres états psychiques (pensées, sentiments, volontés) et de son existence. Elle implique une \cle{dualité sujet-objet} : le sujet se sait sujet.
\end{definition}

\begin{definition}[Conscience spontanée (ou immédiate)]
Conscience \emph{de} quelque chose (objet, perception, émotion). Elle est intentionnelle (Husserl) : toute conscience est conscience \emph{de}. Exemple : « Je vois cet arbre », « J'ai faim ». Pas de retour réflexif sur le « je ».
\end{definition}

\begin{definition}[Conscience réflexive]
Conscience \emph{de soi}. Le sujet se prend lui-même pour objet : « Je sais que je vois cet arbre », « Je pense que j'ai faim ». C'est la marque de l'humanité chez Descartes (\emph{cogito}).
\end{definition}

\begin{aretenir}[Distinction fondamentale]
\begin{itemize}
\item \textbf{Conscience spontanée} : vécue, pré-réflexive, tournée vers le monde (extériorité).
\item \textbf{Conscience réflexive} : pensée, retournement sur soi, fondement de l'identité personnelle et de la responsabilité (intériorité).
\end{itemize}
\end{aretenir}

\begin{exemplebox}[Exemple camerounais]
Au marché Mokolo (Yaoundé), une marchande négocie le prix du plantain (\emph{conscience spontanée} : attention aux prix, aux clients). Soudain, elle se dit : « Je suis en train de m'énerver pour 100 FCFA » (\emph{conscience réflexive} : elle observe sa propre colère).
\end{exemplebox}

\begin{propriete}[Problématique de la leçon]
La conscience est-elle une fenêtre transparente sur nous-mêmes ? \textbf{Thèse classique (Descartes, Alain)} : oui, elle est certitude, liberté, essence de l'homme. \textbf{Antithèse (Nietzsche, Freud, Marx)} : non, elle est partielle, illusoire, déterminée par l'inconscient, le corps, la société. \textbf{Synthèse} : la conscience est une \cle{tâche} (lucidité, examen de soi) plus qu'un donné.
\end{propriete}

\courssub{Nature de la conscience}

\begin{definition}[Intentionalité]
Caractère de toute conscience d'être tournée vers un objet. « Toute conscience est conscience de quelque chose » (Brentano, Husserl). La conscience n'est pas un contenant vide, elle est \emph{relation}.
\end{definition}

\begin{aretenir}[Trois figures de la conscience (programme)]
\begin{enumerate}
\item \textbf{Conscience perceptive / spontanée} : contact direct avec le monde (sensations, perceptions).
\item \textbf{Conscience réflexive / de soi} : le « je pense » accompagnant toute représentation (Descartes, Kant : \emph{Je pense} doit pouvoir accompagner toutes mes représentations).
\item \textbf{Conscience morale} : jugement de valeur sur ses actes (bien/mal), voix intérieure (Kant : \emph{fait de la raison}, Alain : \emph{jugement}).
\end{enumerate}
\end{aretenir}

\begin{savaistu}[Descartes et la transparence]
Dans les \emph{Méditations métaphysiques}, Descartes découvre dans le \emph{cogito} (« Je pense, donc je suis ») une vérité indubitable. La conscience est \cle{transparente à elle-même} : je ne peux pas me tromper sur le fait que je pense (même si je me trompe sur l'objet). C'est le fondement de la subjectivité moderne.
\end{savaistu}

\begin{savaistu}[Alain : la conscience comme jugement]
Pour Alain (\emph{Propos sur les pouvoirs}), « Penser, c'est dire non ». La conscience n'est pas passivité, elle est \cle{puissance de jugement} et de volonté. La lucidité est un effort, une vigilance contre les passions, les préjugés, l'endormissement.
\end{savaistu}

\courssub{Valeur de la conscience}

\begin{propriete}[Valeur épistémologique : la certitude]
Seule la conscience réflexive donne une \cle{certitude subjective} (Descartes). Elle est le point d'Archimède de la connaissance : je doute de tout, mais je ne doute pas que je doute.
\end{propriete}

\begin{propriete}[Valeur éthique et juridique : condition de la responsabilité]
\begin{itemize}
\item \textbf{Kant} : L'autonomie (se donner sa loi) suppose la conscience rationnelle. Sans conscience, pas de liberté, pas de devoir, pas de responsabilité.
\item \textbf{Droit camerounais (Code pénal, art. 64-65)} : L'irresponsabilité pénale pour trouble psychique ou neuropsychique \emph{abolissant} le discernement ou le contrôle des actes. La conscience (discernement) est la condition de l'imputabilité.
\end{itemize}
\end{propriete}

\begin{exemplebox}[Responsabilité et ivresse]
Un chauffeur de taxi à Yaoundé cause un accident en état d'ivresse. Il plaide : « Je ne savais plus ce que je faisais. » Le droit répond : l'ivresse \emph{volontaire} ne supprime pas la responsabilité (on est responsable de s'être mis dans cet état), sauf cas de force majeure (boisson frelatée à son insu). La conscience \emph{antérieure} à l'acte engage.
\end{exemplebox}

\begin{aretenir}[Valeur anthropologique]
« Connais-toi toi-même » (Delphes / Socrate). La conscience réflexive est ce qui permet à l'homme de se constituer comme \cle{sujet moral}, capable de projets, de promesses, de conversion.
\end{aretenir}

\courssub{Limites de la conscience}

\begin{savaistu}[Nietzsche : la conscience comme surface et mensonge]
Dans \emph{Le Gai Savoir} et \emph{Généalogie de la morale}, Nietzsche affirme :
\begin{itemize}
\item La conscience est un « organe » tardif, imparfait, « le plus dangereux » car il simplifie, généralise, ment.
\item « La conscience n'est que la surface de l'esprit » : les vrais moteurs sont les \cle{volontés de puissance}, les instincts, les affects inconscients.
\item La conscience morale est un outil de domestication du troupeau (moralité des esclaves).
\end{itemize}
\end{savaistu}

\begin{savaistu}[Freud : le moi n'est pas maître dans sa propre maison]
Trois blessures à l'amour-propre (Copernic, Darwin, Freud). La psychanalyse montre que la conscience est un \cle{petit îlot} au sommet de l'appareil psychique. La vie psychique est \emph{essentiellement inconsciente}.
\end{savaistu}

\begin{aretenir}[Déterminismes limites de la conscience]
\begin{itemize}
\item \textbf{Inconscient psychique (Freud)} : refoulement, pulsions, désirs interdits.
\item \textbf{Corps / Biologie} : hormones, fatigue, maladies, neurosciences (décisions prises avant conscience, Libet).
\item \textbf{Social / Culturel (Marx, Bourdieu, Durkheim)} : idéologie, habitus, représentations collectives. « Ce n'est pas la conscience des hommes qui détermine leur existence, c'est leur existence sociale qui détermine leur conscience » (Marx).
\end{itemize}
\end{aretenir}

\begin{attention}[Piège fréquent]
Ne pas confondre \emph{inconscient freudien} (dynamique, refoulé, pulsionnel, structuré comme un langage) et \emph{inconscient cognitif/neuroscientifique} (traitements automatiques, priming) ni \emph{préconscient} (disponible à la conscience, souvenirs accessibles).
\end{attention}

\courssub{L'inconscient freudien}

\begin{definition}[Première topique (1900) : Les trois instances topographiques]
\begin{itemize}
\item \textbf{Conscient (Cs)} : perceptions, pensées actuelles. Petite partie.
\item \textbf{Préconscient (Pcs)} : souvenirs, savoirs accessibles à la conscience par un effort d'attention. « Mémoire disponible ».
\item \textbf{Inconscient (Ics)} : \cle{Refoulé}. Représentations chargées d'affect, incompatibles avec le Moi, maintenues hors conscience par la \cle{répression} (refoulement primaire) et le \cle{refoulement proprement dit} (secondaire). Ne connaît ni temps, ni négation, ni contradiction. Régi par les \cle{procédés primaires} (déplacement, condensation).
\end{itemize}
\end{definition}

\begin{definition}[Deuxième topique (1923) : Les trois instances structurales]
\begin{itemize}
\item \textbf{Ça (Es)} : Réservoir des pulsions (vie/Eros, mort/Thanatos), énergie pure, principe de plaisir, amoral, intemporel. « Ça » = l'inconnu en nous.
\item \textbf{Moi (Ich)} : Interface réalité. Principe de réalité. Médiatise entre Ça, Surmoi, Monde extérieur. Siège de la conscience et du préconscient. \textbf{Mais une partie du Moi est inconsciente} (mécanismes de défense).
\item \textbf{Surmoi (Über-Ich)} : Intériorisation des interdits parentaux/sociaux (idéal du moi, conscience morale). Juge sévère, source de culpabilité. En grande partie inconscient.
\end{itemize}
\end{definition}

\begin{aretenir}[Le refoulement (mécanisme fondamental)]
Ce n'est pas l'oubli volontaire. C'est un \cle{processus inconscient} qui repousse dans l'Inconscient/Ça des représentations (idées, souvenirs, désirs) sources de déplaisir (angoisse, culpabilité, conflit moral). Le refoulé \emph{revient} (retour du refoulé) sous formes déguisées.
\end{aretenir}

\begin{exemplebox}[Manifestations de l'inconscient (voie royale)]
\begin{itemize}
\item \textbf{Rêve} : « Accomplissement déguisé d'un désir refoulé ». Travail du rêve : condensation, déplacement, figuration, élaboration secondaire. Contenu latent (désir) $\to$ contenu manifeste (rêve raconté).
\item \textbf{Actes manqués (lapsus, oubli, méprise, bévue)} : Satisfaction substitutive d'un désir refoulé. Ex : Lapsus « La conscience est l'ennemi... l'essence de l'homme ».
\item \textbf{Symptômes névrotiques} : Compromis entre désir refoulé et défense. Phobies, TOC, conversion (hystérie).
\item \textbf{Transfert} : Revécu infantile sur la personne de l'analyste.
\end{itemize}
\end{exemplebox}

\courssub{Valeur et limites de la psychanalyse}

\begin{propriete}[Valeur thérapeutique et herméneutique]
\begin{itemize}
\item \textbf{Cure par la parole} : Association libre, levée du refoulement, \cle{insight} (prise de conscience). « Là où était le Ça, le Moi doit advenir » (Freud).
\item \textbf{Interprétation} : Décryptage du sens caché (herméneutique du soupçon, Ricoeur). Redonne du sens à l'histoire du sujet.
\item \textbf{Autonomie} : Le sujet reprend possession de son histoire, réduit l'emprise de la répétition compulsive.
\end{itemize}
\end{propriete}

\begin{propriete}[Limites et débat épistémologique : La psychanalyse est-elle une science ?]
\begin{itemize}
\item \textbf{Thèse (Pro-science)} : Observation clinique rigoureuse, concepts opératoires (refoulement, transfert), prédictions (structure des névroses), efficacité thérapeutique vérifiable. Popper lui-même a nuancé : « métaphysique » ne veut pas dire « sans valeur ».
\item \textbf{Antithèse (Popper, critère de \cle{réfutabilité})} : Une théorie scientifique doit pouvoir être \emph{réfutée} par l'expérience. La psychanalyse explique \emph{tout} a posteriori : le patient accepte $\to$ confirmation ; il refuse $\to$ « résistance » = confirmation. \textbf{Irréfutable = non scientifique}.
\item \textbf{Grünbaum} : Effet placebo, suggestion, pas de groupe témoin.
\item \textbf{Synthèse (Ricoeur, Habermas)} : La psychanalyse n'est pas une science de la nature (nomologique) mais une \cle{science de l'interprétation} (herméneutique), une « archéologie du sujet ». Valide sur le plan clinique et anthropologique, discutable sur le plan épistémologique poppérien.
\end{itemize}
\end{propriete}

\begin{attention}[Erreurs à ne pas commettre au Bac]
\begin{itemize}
\item Dire « L'inconscient, c'est ce dont on n'a pas conscience » (trop vague : inclut le préconscient, l'inconscient cognitif).
\item Confondre \emph{refoulement} (inconscient, pathogène) et \emph{répression} (consciente, volontaire, saine).
\item Dire « Freud a découvert l'inconscient » (Leibniz, Schopenhauer, Janet, Charcot avant lui). Freud a \emph{théorisé l'inconscient dynamique et refoulé}.
\item Oublier la \textbf{deuxième topique} (Ça, Moi, Surmoi) qui complexifie le rapport conscience/inconscient (Moi partiellement inconscient).
\end{itemize}
\end{attention}

\courssub{Synthèse, méthode et entraînement au Bac}

\begin{methode}[Réussir une dissertation de philosophie (Bac Tech)]
\begin{enumerate}
\item \textbf{Analyse du sujet} (15 min) : Définir \emph{tous} les termes, repérer le présupposé, formuler le \textbf{problème} (question tension), annoncer le \textbf{plan dialectique} (Thèse / Antithèse / Synthèse).
\item \textbf{Recherche d'idées} (15 min) : Fiche auteurs/notions/exemples par partie. Au moins 3 auteurs précis.
\item \textbf{Rédaction} (1h15) :
\begin{itemize}
\item \textbf{Introduction} : Accroche $\to$ Définitions $\to$ Problème $\to$ Plan.
\item \textbf{Développement} : 3 parties équilibrées. \textbf{Une thèse = un argument principal + auteur + exemple + transition}. Connecteurs logiques (\emph{en effet, toutefois, par conséquent}).
\item \textbf{Conclusion} : Réponse directe au problème $\to$ Résumé serré $\to$ Ouverture (autre champ, actualité, paradoxe).
\end{itemize}
\item \textbf{Relecture} (15 min) : Orthographe, syntaxe, respect du sujet, noms propres exacts.
\end{enumerate}
\end{methode}

\begin{methode}[Réussir une explication de texte]
\begin{enumerate}
\item \textbf{Lecture active} : Auteur, œuvre, date, contexte. Souligner thèse, arguments, concepts-clés, connecteurs.
\item \textbf{Introduction} : Présentation (auteur, œuvre, thèse générale), problème du texte, annonce du plan (suivre le mouvement du texte).
\item \textbf{Développement} : Expliquer \emph{pas à pas} (par paragraphe ou idée) : \textbf{Citer} $\to$ \textbf{Expliquer} (vocabulaire, référence, contexte) $\to$ \textbf{Analyser} (portée, enjeu).
\item \textbf{Conclusion} : Résumé de l'explication + \textbf{Discussion critique} (argumentée, nuancée, un auteur du programme).
\end{enumerate}
\end{methode}

\courssub{Je vérifie mes connaissances}

\begin{exercice}[QCM : la conscience et l'inconscient]
Pour chaque question, une seule réponse est exacte. Recopie la lettre correspondante.
\begin{enumerate}
\item La conscience de soi (conscience réflexive) est :
\begin{enumerate}
\item[a)] la simple perception d'un objet extérieur ;
\item[b)] la capacité de l'esprit à se prendre lui-même pour objet ;
\item[c)] l'ensemble des souvenirs oubliés.
\end{enumerate}
\item Pour Descartes, le \emph{cogito} (« je pense, donc je suis ») exprime :
\begin{enumerate}
\item[a)] une vérité indubitable saisie par la conscience ;
\item[b)] une illusion des sens ;
\item[c)] un produit de l'inconscient.
\end{enumerate}
\item Selon Freud, le refoulement est :
\begin{enumerate}
\item[a)] l'oubli volontaire d'un mauvais souvenir ;
\item[b)] le mécanisme par lequel un contenu psychique inacceptable est maintenu hors de la conscience ;
\item[c)] la guérison définitive d'un traumatisme.
\end{enumerate}
\item Un lapsus est, pour la psychanalyse :
\begin{enumerate}
\item[a)] une simple erreur de langue sans signification ;
\item[b)] un acte manqué révélant un désir inconscient ;
\item[c)] une preuve de fatigue intellectuelle.
\end{enumerate}
\end{enumerate}
\end{exercice}

\begin{exercice}[Vrai ou faux justifié]
Pour chaque affirmation, indique si elle est vraie ou fausse, puis justifie ta réponse en deux ou trois phrases.
\begin{enumerate}
\item « La conscience nous donne une connaissance parfaite et totale de nous-mêmes. »
\item « Pour Freud, le moi n'est pas maître dans sa propre maison. »
\item « L'inconscient freudien est un lieu de la mémoire où dorment les souvenirs oubliés. »
\item « La psychanalyse prétend soigner certains troubles psychiques par la parole. »
\end{enumerate}
\end{exercice}

\begin{exercice}[Définitions]
Définis en trois à cinq lignes chacune des notions suivantes, en donnant pour chacune un exemple concret tiré de la vie courante :
\begin{enumerate}
\item la conscience spontanée (conscience immédiate) ;
\item la conscience réflexive ;
\item l'inconscient (freudien) ;
\item le refoulement ;
\item l'acte manqué.
\end{enumerate}
\end{exercice}

\begin{exercice}[Auteurs et thèses]
Relie chaque auteur (colonne de gauche) à la thèse qui lui correspond (colonne de droite, mélangées), puis reformule chaque thèse avec tes propres mots.
\begin{center}
\begin{tabularx}{\linewidth}{|l|X|}
\hline
\rowcolor{popblueL} \textbf{Auteur} & \textbf{Thèse (à relier)} \\
\hline
Descartes & 1. « Le moi n'est pas maître dans sa propre maison » : l'homme est dominé par des forces inconscientes. \\
\hline
Freud & 2. La conscience est la marque certaine de l'existence du sujet pensant. \\
\hline
Nietzsche & 3. La conscience n'est que la surface de l'esprit ; les instincts et les forces profondes décident. \\
\hline
Alain & 4. La conscience est puissance de jugement et de volonté : penser, c'est dire non. \\
\hline
\end{tabularx}
\end{center}
\end{exercice}

\courssub{Je m'entraîne}

\begin{exercice}[Analyse de situations camerounaises]
Pour chacune des trois situations suivantes, identifie le rôle joué par la conscience, par les déterminismes sociaux et par l'inconscient, puis justifie ta réponse par écrit en un paragraphe argumenté (cinq à huit lignes).
\begin{enumerate}
\item \textbf{Situation 1.} Amina, élève de Terminale à Douala, s'inscrit en série C alors qu'elle rêve de devenir couturière : sa famille et son entourage répètent qu'« un bon élève doit faire les sciences ».
\item \textbf{Situation 2.} Après avoir accusé à tort son camarade d'avoir volé son téléphone, Tchoupo éprouve un remords durable qui l'empêche de dormir.
\item \textbf{Situation 3.} Depuis plusieurs semaines, Ngo Bell rêve chaque nuit qu'elle arrive en retard à l'examen et que sa feuille reste blanche.
\end{enumerate}
\end{exercice}

\begin{exercice}[Le lapsus révélateur]
Lors d'un exposé de philosophie devant toute la classe, un élève déclare : « La conscience est l'ennemi... je veux dire, l'essence de l'homme ! » La classe rit ; l'élève, gêné, affirme que ce n'est « qu'une erreur ».
\begin{enumerate}
\item Rappelle ce que Freud entend par \cle{acte manqué} et par \cle{refoulement}.
\item Rédige un paragraphe argumenté (dix à quinze lignes) montrant comment la psychanalyse interpréterait ce lapsus, puis indique ce qu'une analyse purement consciente (simple fatigue, stress de l'exposé) pourrait objecter.
\end{enumerate}
\end{exercice}

\begin{exercice}[Conscience et responsabilité]
Un chauffeur de taxi à Yaoundé, surpris en état d'ivresse après un accident, se défend en disant : « Je ne savais plus ce que je faisais. »
\begin{enumerate}
\item Explique pourquoi la conscience est traditionnellement considérée comme la condition de la responsabilité morale et juridique.
\item La perte de conscience (ivresse, trouble psychique) excuse-t-elle totalement l'agent ? Argumente en distinguant la responsabilité de l'acte et la responsabilité de s'être mis dans cet état.
\end{enumerate}
\end{exercice}

\begin{exercice}[La psychanalyse est-elle une science ?]
On te demande de construire le plan dialectique (thèse / antithèse / synthèse) d'une réflexion sur la question : « La psychanalyse est-elle une science ? »
\begin{enumerate}
\item Rappelle le critère de \cle{réfutabilité} proposé par Karl Popper : une théorie est scientifique si elle peut être mise à l'épreuve et éventuellement réfutée par les faits.
\item Rédige la thèse (la psychanalyse est une science : méthode d'observation clinique, interprétation des rêves, résultats thérapeutiques).
\item Rédige l'antithèse (la psychanalyse n'est pas une science : ses concepts — inconscient, refoulement — échappent à toute vérification expérimentale ; toute objection peut être interprétée comme une « résistance », ce qui rend la théorie irréfutable).
\item Propose une synthèse nuancée en cinq lignes.
\end{enumerate}
\end{exercice}

\courssub{Je me prépare au Bac}

\begin{exercice}[Dissertation : la conscience suffit-elle à définir l'homme ?]
\textbf{Sujet :} \emph{La conscience suffit-elle à définir l'homme ?}

\textbf{Consignes (durée : 2 heures).} Rédige une dissertation complète en suivant la méthode du Baccalauréat technique :
\begin{enumerate}
\item \textbf{Introduction} : accroche, définition des termes du sujet (conscience, définir, homme), formulation du problème (la conscience est-elle le propre de l'homme et son essence tout entière ?), annonce du plan.
\item \textbf{Développement en trois parties} :
\begin{enumerate}
\item[(a)] Thèse : la conscience comme essence de l'homme (Descartes, le \emph{cogito} ; la conscience, condition de la liberté et de la responsabilité).
\item[(b)] Antithèse : les limites de la conscience (illusions de la conscience chez Nietzsche ; découverte de l'inconscient chez Freud ; déterminismes sociaux et corporels).
\item[(c)] Synthèse : l'homme est un être de conscience \emph{et} d'inconscient ; la conscience reste une tâche à accomplir (lucidité, examen de soi) plutôt qu'une évidence.
\end{enumerate}
\item \textbf{Conclusion} : réponse claire à la question posée et ouverture (par exemple : connais-toi toi-même).
\end{enumerate}
\textbf{Exigences :} au moins trois auteurs cités avec exactitude, deux exemples concrets dont un tiré du contexte camerounais, transitions explicites entre les parties.
\end{exercice}

\begin{exercice}[Dissertation : peut-on être tenu responsable de ses rêves ?]
\textbf{Sujet :} \emph{Peut-on être tenu responsable de ses rêves ?}

\textbf{Consignes (durée : 2 heures).}
\begin{enumerate}
\item Analyse du sujet : définis \cle{responsabilité} et \cle{rêve} ; dégage le présupposé de la question (le rêve exprimerait un désir qui nous appartient).
\item Première partie : nous ne sommes pas responsables de nos rêves, car ils échappent à la volonté et à la conscience (le rêve se subit ; la responsabilité suppose la conscience et la liberté).
\item Deuxième partie : pourtant, selon Freud, le rêve est « l'accomplissement d'un désir » refoulé ; il révèle notre inconscient, donc une part de nous-mêmes. Peut-on être responsable de ce que l'on est sans l'avoir choisi ?
\item Troisième partie : synthèse — distinguer la responsabilité juridique et morale des actes (qui exige la conscience) de la responsabilité de se connaître soi-même (devoir de lucidité face à ses désirs).
\item Conclusion argumentée.
\end{enumerate}
\textbf{Exigences :} problématique explicite, mobilisation des notions d'\cle{inconscient} et de \cle{refoulement}, au moins un contre-exemple discuté.
\end{exercice}

\begin{exercice}[Explication de texte : Freud, \emph{Une difficulté de la psychanalyse}]
\textbf{Texte.} « L'humanité, au cours de sa longue évolution, a dû subir de la part de la science deux grandes blessures infligées à son amour-propre naïf. La première, lorsqu'elle apprit que notre terre n'était pas le centre de l'univers [...]. La seconde, lorsque la recherche biologique ravit à l'homme sa prétendue supériorité sur les autres espèces [...]. Mais la troisième blessure, la plus sensible, c'est la psychologie moderne qui la lui inflige, en montrant que le moi n'est pas même maître dans sa propre maison, qu'il se contente de renseignements rares et fragmentaires sur ce qui se passe, à son insu, dans sa vie psychique. » (Freud, adapté)

\textbf{Questions :}
\begin{enumerate}
\item Dégage la \cle{thèse} du texte en une phrase.
\item Explique les deux premières « blessures » évoquées par Freud (Copernic, Darwin) et montre en quoi la psychanalyse constitue la troisième.
\item Analyse la métaphore « le moi n'est pas maître dans sa propre maison » : que signifient le « moi », la « maison » et l'absence de maîtrise ?
\item Quelle est la portée philosophique de ce texte sur la question de la conscience et de la connaissance de soi ?
\item Discussion : cette conception ruine-t-elle toute possibilité de liberté et de responsabilité humaines ? Rédige un paragraphe argumenté (quinze à vingt lignes) en mobilisant au moins un auteur du programme.
\end{enumerate}
\end{exercice}

\newpage

\coursec{Corrigés des exercices}

\coursec{Corrigés des exercices : La conscience et l'inconscient}

\begin{corrige}[Exercice 1 — QCM : la conscience et l'inconscient]
\textbf{Réponses :}
\begin{enumerate}
\item \textbf{b)} La conscience réflexive est la capacité de l'esprit à se prendre lui-même pour objet : « Je sais que je pense ». La réponse a) décrit la conscience spontanée (tournée vers l'objet) ; la réponse c) évoque le préconscient ou l'inconscient, non la conscience.
\item \textbf{a)} Le \emph{cogito} est une vérité indubitable : même en doutant de tout, je ne peux douter que je doute, donc que je pense. La conscience est ici certitude immédiate de soi.
\item \textbf{b)} Le refoulement est un mécanisme \emph{inconscient} (et non un oubli volontaire, réponse a) qui maintient hors de la conscience un contenu inacceptable (désir, souvenir traumatique). Il ne guérit rien (réponse c) : le refoulé revient sous formes déguisées.
\item \textbf{b)} Le lapsus est un acte manqué : il révèle, selon Freud, un désir ou une pensée refoulée qui s'exprime de manière déguisée. Ce n'est ni un hasard insignifiant (a) ni un simple effet de fatigue (c).
\end{enumerate}
\textbf{Points de vigilance :} bien distinguer \emph{oubli volontaire} (conscient) et \emph{refoulement} (inconscient) ; rappeler que pour Freud, « rien n'arrive par hasard » dans la vie psychique.
\end{corrige}

\begin{corrige}[Exercice 2 — Vrai ou faux justifié]
\begin{enumerate}
\item \textbf{Faux.} La conscience ne nous donne qu'une connaissance partielle de nous-mêmes. Nietzsche montre qu'elle n'est que « la surface de l'esprit » et Freud qu'elle n'est qu'un « petit îlot » au sommet de la vie psychique, essentiellement inconsciente. Nos actes ont des motifs cachés (pulsions, déterminismes sociaux) qui échappent à la conscience.
\item \textbf{Vrai.} C'est la formule même de Freud (\emph{Une difficulté de la psychanalyse}) : le moi ne reçoit que des « renseignements rares et fragmentaires » sur sa vie psychique. La psychanalyse inflige ainsi à l'homme la troisième « blessure narcissique », après Copernic et Darwin.
\item \textbf{Faux.} Cette définition correspond au \emph{préconscient} (souvenirs accessibles par un effort d'attention). L'inconscient freudien est dynamique : il contient des représentations \emph{refoulées}, incompatibles avec le moi, maintenues hors conscience par une force active (le refoulement), et il s'exprime par les rêves, les lapsus, les symptômes.
\item \textbf{Vrai.} La psychanalyse est une « cure par la parole » : par l'association libre et l'interprétation, elle vise à lever le refoulement et à rendre conscient ce qui était inconscient (« Là où était le Ça, le Moi doit advenir »), soulageant ainsi névroses et symptômes.
\end{enumerate}
\textbf{Points de vigilance :} la justification doit mobiliser une notion précise du cours (préconscient, refoulement, blessure narcissique) et non rester générale.
\end{corrige}

\begin{corrige}[Exercice 3 — Définitions]
\begin{enumerate}
\item \textbf{Conscience spontanée (immédiate)} : connaissance directe et non réfléchie que le sujet a d'un objet ou de son état présent, sans retour sur le « je ». Elle est intentionnelle : conscience \emph{de} quelque chose. \emph{Exemple :} un élève de Ngaoundéré sent la chaleur du soleil en marchant vers le lycée, sans se dire « je me perçois en train de sentir ».
\item \textbf{Conscience réflexive} : retour de l'esprit sur lui-même ; le sujet se prend lui-même pour objet (« je sais que je pense »). Fondement de l'identité, de la liberté et de la responsabilité (Descartes). \emph{Exemple :} après m'être emporté contre mon petit frère, je me dis : « Pourquoi ai-je réagi ainsi ? »
\item \textbf{Inconscient (freudien)} : ensemble des contenus psychiques (désirs, souvenirs, pulsions) \emph{refoulés} hors de la conscience parce qu'inacceptables pour le moi, mais qui demeurent actifs et s'expriment de manière déguisée. Il ignore le temps, la négation, la contradiction. \emph{Exemple :} une phobie inexpliquée des chiens remontant à un traumatisme d'enfance oublié.
\item \textbf{Refoulement} : mécanisme de défense \emph{inconscient} par lequel le moi repousse hors de la conscience une représentation source de déplaisir (angoisse, culpabilité). Ce n'est pas un oubli volontaire ; le refoulé fait retour (symptômes, rêves). \emph{Exemple :} un témoin d'un grave accident sur l'axe Yaoundé--Douala n'en garde aucun souvenir, mais fait des cauchemars récurrents.
\item \textbf{Acte manqué} : geste raté en apparence accidentel (lapsus, oubli, méprise, bévue) qui, selon Freud, réalise de façon déguisée un désir ou une intention refoulée. \emph{Exemple :} oublier « par hasard » le rendez-vous chez un dentiste que l'on redoute.
\end{enumerate}
\textbf{Points de vigilance :} chaque définition doit contenir un trait distinctif précis (réflexivité, refoulement, caractère inconscient du mécanisme) et un exemple \emph{différent} de ceux du cours si possible.
\end{corrige}

\begin{corrige}[Exercice 4 — Auteurs et thèses]
\textbf{Correspondances :} Descartes $\to$ 2 ; Freud $\to$ 1 ; Nietzsche $\to$ 3 ; Alain $\to$ 4.

\textbf{Reformulations attendues :}
\begin{itemize}
\item \textbf{Descartes (2)} : dans le doute universel, une seule certitude résiste : le fait même que je pense prouve que j'existe comme sujet pensant. La conscience de soi est le fondement indubitable de toute connaissance.
\item \textbf{Freud (1)} : la conscience n'est qu'une petite partie de la vie psychique ; nos pensées et nos actes sont largement déterminés par des désirs et des conflits inconscients que nous ignorons.
\item \textbf{Nietzsche (3)} : la conscience est un phénomène tardif et superficiel qui simplifie et ment ; les véritables moteurs de nos actions sont les instincts, les affects et la volonté de puissance.
\item \textbf{Alain (4)} : la conscience n'est pas une contemplation passive mais une activité de jugement : penser, c'est refuser les évidences, les passions et les préjugés ; la lucidité est un effort de vigilance.
\end{itemize}
\textbf{Points de vigilance :} ne pas attribuer à Freud la « découverte » de l'inconscient (il l'a \emph{théorisé}) ; ne pas confondre la critique nietzschéenne (généalogique) et la critique freudienne (clinique).
\end{corrige}

\begin{corrige}[Exercice 5 — Analyse de situations camerounaises]
\textbf{Situation 1 (Amina).} Le facteur dominant est le \emph{déterminisme social} : la représentation collective « un bon élève fait les sciences » fonctionne comme une idéologie intériorisée (Marx : l'existence sociale détermine la conscience ; Durkheim : les représentations collectives s'imposent à l'individu). La conscience d'Amina est partiellement aliénée : elle croit choisir librement alors que son choix est orienté par la famille et le prestige social. Mais la conscience réflexive reste possible : en s'examinant (« que désiré-je vraiment ? »), Amina peut reprendre la maîtrise de son projet. La situation illustre donc à la fois les limites de la conscience et sa valeur comme tâche de lucidité.

\textbf{Situation 2 (Tchoupo).} C'est la \emph{conscience morale} qui est à l'œuvre : le remords est la voix intérieure qui juge l'acte passé (Kant : fait de la raison ; Alain : jugement). Tchoupo se prend lui-même pour objet (conscience réflexive) et s'évalue. Ce remords prouve la valeur éthique de la conscience : elle fonde la responsabilité et appelle la réparation (avouer l'erreur, s'excuser). On peut noter que l'accusation hâtive initiale pouvait elle-même être influencée par un préjugé inconscient (par exemple une rivalité refoulée), ce qui montre l'entrelacs de la conscience et de l'inconscient.

\textbf{Situation 3 (Ngo Bell).} Le rêve répétitif relève de l'\emph{inconscient freudien} : le rêve est « l'accomplissement déguisé d'un désir refoulé » ou l'expression d'une angoisse non élaborée. Ici, l'angoisse de l'échec à l'examen, peut-être refoulée le jour (« tout va bien, je gère »), revient la nuit sous forme figurée (retard, feuille blanche = peur de l'échec). La conscience diurne maîtrise en apparence, mais l'inconscient révèle le conflit. La prise de conscience de cette angoisse (en parler, se préparer) peut en réduire l'emprise : « Là où était le Ça, le Moi doit advenir ».

\textbf{Points de vigilance :} pour chaque situation, nommer explicitement la notion mobilisée, citer un auteur, et éviter de tout ramener à l'inconscient (la conscience et le social ont aussi leur rôle).
\end{corrige}

\begin{corrige}[Exercice 6 — Le lapsus révélateur]
\textbf{1. Définitions.} Le \emph{refoulement} est le mécanisme inconscient par lequel le moi maintient hors de la conscience un contenu psychique (désir, pensée, souvenir) inacceptable, source de déplaisir ou de conflit. L'\emph{acte manqué} (lapsus, oubli, méprise) est une action ratée en apparence accidentelle, que Freud interprète comme la satisfaction déguisée d'un désir ou d'une pensée refoulée : le refoulé « fait retour » en profitant d'une baisse de la censure.

\textbf{2. Paragraphe argumenté (modèle).} Pour la psychanalyse, le lapsus de l'élève n'est pas un hasard : « rien n'arrive par hasard » dans la vie psychique. En disant « la conscience est l'\emph{ennemi} de l'homme » au lieu de « l'\emph{essence} de l'homme », l'élève a laissé s'exprimer une pensée refoulée. L'interprétation peut prendre plusieurs directions : peut-être l'élève éprouve-t-il, sans se l'avouer, une hostilité envers la philosophie ou envers l'exigence de lucidité que la matière représente ; peut-être a-t-il été marqué par la lecture de Nietzsche, pour qui la conscience est effectivement présentée comme un organe dangereux et menteur, et cette thèse, séduisante mais inavouable devant l'enseignant, s'est imposée malgré lui ; peut-être enfin l'angoisse de l'exposé a-t-elle affaibli la censure, laissant passer un doute refoulé sur sa propre maîtrise de soi. Le lapsus serait ainsi un compromis entre le désir inconscient de dire « ennemi » et l'intention consciente de dire « essence ». Cependant, une analyse purement consciente objecterait à juste titre que la fatigue, le trac devant la classe ou la simple proximité phonique et conceptuelle des deux mots (la leçon oppose justement la conscience à son « ennemi », l'inconscient) suffisent à expliquer l'erreur, sans hypothèse inconsciente. Cette objection illustre la limite épistémologique de la psychanalyse : toute explication alternative peut être réinterprétée comme « résistance », ce qui rend la théorie difficilement réfutable (Popper).

\textbf{Points de vigilance :} articuler les deux interprétations sans trancher dogmatiquement ; mobiliser les notions de refoulement, de censure, de compromis ; conclure sur la portée et la limite de l'interprétation psychanalytique.
\end{corrige}

\begin{corrige}[Exercice 7 — Conscience et responsabilité]
\textbf{1. La conscience, condition de la responsabilité.} Être responsable, c'est pouvoir répondre de ses actes, c'est-à-dire en être reconnu l'auteur et en assumer les conséquences (louange, blâme, sanction). Or cela suppose deux conditions : la \emph{connaissance} de ce que l'on fait (discernement du bien et du mal) et la \emph{liberté} (contrôle de ses actes). Sans conscience réflexive, l'acte n'est plus imputable au sujet : on ne juge pas un somnambule. C'est pourquoi Kant fait de l'autonomie rationnelle le fondement du devoir, et pourquoi le droit (Code pénal camerounais, art. 64-65) prévoit l'irresponsabilité pénale lorsqu'un trouble psychique \emph{abolit} le discernement ou le contrôle des actes au moment des faits.

\textbf{2. L'ivresse excuse-t-elle ? (modèle de réponse).} Il faut distinguer deux moments. Au moment de l'acte, l'ivresse a effectivement diminué ou supprimé le discernement : à strictement parler, l'agent ivre n'était pas pleinement conscient, donc pas pleinement imputable de l'accident lui-même. Mais la responsabilité se déplace alors en amont : le chauffeur était conscient et libre \emph{au moment de boire}, alors qu'il savait qu'il conduirait. Il est donc responsable de s'être volontairement mis dans un état d'incapacité : c'est une faute de prévision et d'imprudence. Le droit retient d'ailleurs ce raisonnement : l'ivresse volontaire est en général une circonstance aggravante et non excusante ; seule une perte de conscience \emph{involontaire et imprévisible} (boisson frelatée à son insu, trouble psychique soudain) exonère totalement. La conscience antérieure à l'acte engage donc la responsabilité : on répond non seulement de ce que l'on fait, mais de ce que l'on a rendu possible en connaissance de cause.

\textbf{Points de vigilance :} citer la distinction discernement/contrôle ; articuler responsabilité de l'acte et responsabilité de la mise en état ; éviter la réponse manichéenne (« il est coupable » / « il est innocent »).
\end{corrige}

\begin{corrige}[Exercice 8 — La psychanalyse est-elle une science ?]
\textbf{1. Critère de Popper.} Pour Karl Popper, une théorie est scientifique si et seulement si elle est \emph{réfutable} (falsifiable) : elle doit formuler des prédictions précises qui pourraient être démenties par l'expérience. Une théorie compatible avec \emph{tous} les faits possibles n'explique rien : l'irréfutabilité est un défaut, non une force.

\textbf{2. Thèse (la psychanalyse est une science).} La psychanalyse repose sur une méthode d'observation rigoureuse : la clinique, l'écoute des patients, l'association libre, l'étude des rêves et des actes manqués. Freud a construit des concepts opératoires (refoulement, transfert, topiques) permettant de classer les troubles (névroses, psychoses) et d'anticiper leur évolution. Surtout, la cure produit des effets vérifiables : disparition de symptômes, levée de phobies, reprise d'autonomie. Comme toute science, elle part du réel (les discours des patients) pour construire un savoir explicatif cohérent.

\textbf{3. Antithèse (la psychanalyse n'est pas une science).} Ses concepts fondamentaux — inconscient, refoulement, pulsions — sont par définition inobservables et échappent à toute mesure expérimentale : on ne peut ni isoler l'inconscient en laboratoire, ni constituer de groupe témoin fiable (Grünbaum dénonce l'effet placebo et la suggestion). Surtout, la théorie est \emph{irréfutable} au sens de Popper : si le patient accepte l'interprétation, elle est confirmée ; s'il la refuse, c'est une « résistance », donc elle est encore confirmée. Aucun fait ne peut la démentir : elle explique tout a posteriori, ce qui est le contraire de la démarche scientifique.

\textbf{4. Synthèse nuancée (modèle, cinq lignes).} La psychanalyse n'est pas une science de la nature au sens poppérien : elle ne produit ni lois mesurables ni prédictions réfutables. Mais elle n'est pas pour autant dénuée de valeur : c'est une \emph{science de l'interprétation} (Ricœur : herméneutique du soupçon), une méthode de déchiffrement du sens caché des conduites humaines, validée par sa pratique clinique. Sa scientificité stricte est discutable, sa fécondité thérapeutique et anthropologique ne l'est pas.

\textbf{Points de vigilance :} définir correctement la réfutabilité avant de l'appliquer ; ne pas conclure « la psychanalyse ne sert à rien » (confondre non-scientifique et sans valeur) ; la synthèse doit déplacer la question (quel \emph{type} de savoir ?).
\end{corrige}

\begin{corrige}[Exercice 9 — Dissertation : la conscience suffit-elle à définir l'homme ?]
\textbf{Corrigé détaillé (plan et contenus attendus).}

\textbf{Introduction.} \emph{Accroche :} depuis Socrate, « connais-toi toi-même » fait de la conscience le chemin de la sagesse. \emph{Définitions :} la conscience est la connaissance immédiate que le sujet a de ses états et de son existence (spontanée et réflexive) ; « définir » signifie saisir l'essence, le propre ; « l'homme » désigne l'être pensant et agissant. \emph{Problème :} la conscience est-elle le propre et l'essence entière de l'homme, ou n'est-elle que la partie émergée d'un être traversé par l'inconscient et les déterminismes ? \emph{Annonce du plan :} nous verrons d'abord que la conscience fonde l'humanité de l'homme, puis qu'elle est limitée et trompeuse, enfin qu'elle demeure une tâche plutôt qu'un donné.

\textbf{I. Thèse : la conscience, essence de l'homme.}
\begin{itemize}
\item \emph{Argument 1 (Descartes)} : le \emph{cogito} est la première certitude ; je peux douter de tout sauf que je pense. La conscience réflexive est le fondement indubitable de l'existence du sujet.
\item \emph{Argument 2 (Kant, Alain)} : la conscience est condition de la liberté, de la moralité et de la responsabilité ; « penser, c'est dire non » (Alain) : juger, c'est s'arracher aux passions.
\item \emph{Exemple :} le remords de Tchoupo après sa fausse accusation montre la conscience morale en acte ; sans elle, ni promesse, ni projet, ni dignité.
\end{itemize}

\textbf{II. Antithèse : les limites de la conscience.}
\begin{itemize}
\item \emph{Argument 1 (Nietzsche)} : la conscience n'est que « la surface de l'esprit », un organe tardif qui simplifie et ment ; les vrais moteurs sont instincts et volonté de puissance.
\item \emph{Argument 2 (Freud)} : « le moi n'est pas maître dans sa propre maison » ; rêves, lapsus, symptômes révèlent un inconscient dynamique qui détermine nos conduites.
\item \emph{Argument 3 (Marx)} : « ce n'est pas la conscience des hommes qui détermine leur existence, c'est leur existence sociale qui détermine leur conscience ».
\item \emph{Exemple camerounais :} Amina croit choisir librement la série C, mais son choix est orienté par la pression familiale et le prestige social des sciences.
\end{itemize}

\textbf{III. Synthèse : la conscience comme tâche.}
\begin{itemize}
\item Ni essence transparente (Descartes se trompe sur la toute-puissance du moi), ni pure illusion (Nietzsche et Freud \emph{pensent} encore : la critique de la conscience est un acte de conscience).
\item L'homme est un être de conscience \emph{et} d'inconscient ; la lucidité n'est pas un donné mais un travail : examen de soi, écoute de ses désirs, vigilance critique. « Là où était le Ça, le Moi doit advenir » (Freud) : la conscience se définit comme conquête.
\end{itemize}

\textbf{Conclusion.} La conscience ne suffit pas à définir l'homme, car elle n'en est qu'une part, limitée par l'inconscient, le corps et le social ; mais elle reste ce par quoi l'homme peut se saisir de ses limites : moins son essence que sa vocation. \emph{Ouverture :} les neurosciences contemporaines prolongent-elles ou ruinent-elles cette exigence de lucidité ?

\textbf{Barème indicatif (sur 20) :} introduction (problème, définitions, plan) : 4 pts ; thèse argumentée avec auteurs : 4 pts ; antithèse argumentée avec auteurs : 4 pts ; synthèse personnelle et nuancée : 4 pts ; exemples (dont un camerounais), transitions, langue : 4 pts.

\textbf{Points de vigilance :} au moins trois auteurs exacts (Descartes, Nietzsche, Freud, Marx, Alain) ; ne pas opposer conscience et inconscient comme deux « choses » mais comme deux dimensions ; la synthèse ne doit pas être un compromis mou mais un dépassement (la conscience comme tâche).
\end{corrige}

\begin{corrige}[Exercice 10 — Dissertation : peut-on être tenu responsable de ses rêves ?]
\textbf{Corrigé détaillé.}

\textbf{Analyse du sujet.} \emph{Responsabilité} : obligation de répondre de ses actes, supposant conscience et liberté. \emph{Rêve} : production psychique involontaire du sommeil. \emph{Présupposé de la question :} le rêve exprimerait un désir qui nous appartient — sinon la question de la responsabilité ne se poserait même pas. \emph{Problème :} peut-on être responsable de ce qui échappe à notre volonté, si ce qui échappe à notre volonté vient pourtant de nous ?

\textbf{I. Nous ne sommes pas responsables de nos rêves.}
\begin{itemize}
\item Le rêve se \emph{subit} : il survient sans décision, hors de tout contrôle de la conscience et de la volonté. Or la responsabilité morale et juridique exige le discernement et la liberté (Kant) : on ne juge pas un somnambule.
\item Le droit ne sanctionne que les actes : nul ne peut être condamné pour ses pensées ou ses songes.
\item \emph{Exemple :} Ngo Bell rêve chaque nuit d'arriver en retard à l'examen : lui reprocher ce rêve serait absurde, elle ne l'a ni voulu ni choisi.
\end{itemize}

\textbf{II. Pourtant, le rêve nous révèle (Freud).}
\begin{itemize}
\item Pour Freud, le rêve est « l'accomplissement déguisé d'un désir refoulé » : le contenu manifeste (le rêve raconté) cache un contenu latent (le désir inconscient) par le travail de condensation et de déplacement.
\item Le rêve est donc bien \emph{à nous} : il exprime notre inconscient, c'est-à-dire une part de nous-mêmes que nous ignorons. \emph{Contre-exemple discuté :} le rêveur qui rêve d'un acte violent n'est pas coupable de violence, mais le rêve signale un conflit ou une agressivité refoulée qui le concernent.
\item Question ouverte : peut-on être responsable de ce que l'on \emph{est} sans l'avoir choisi ?
\end{itemize}

\textbf{III. Synthèse : deux sens de la responsabilité.}
\begin{itemize}
\item Il faut distinguer la \emph{responsabilité juridique et morale des actes}, qui exige conscience et liberté — et qui ne s'applique donc pas aux rêves —, et la \emph{responsabilité de se connaître soi-même} : un devoir de lucidité face à ses désirs.
\item Nous ne répondons pas de nos rêves devant autrui, mais nous répondons de ce que nous faisons de ce qu'ils révèlent : ignorer délibérément ses désirs refoulés, c'est risquer de les laisser agir à notre insu (retour du refoulé). « Connais-toi toi-même » devient une tâche éthique.
\end{itemize}

\textbf{Conclusion.} Non, nous ne sommes pas responsables de nos rêves au sens où nous répondons de nos actes ; mais oui, nous sommes responsables de l'attention que nous portons à ce qu'ils nous apprennent de nous-mêmes. \emph{Ouverture :} la cure analytique réalise précisément ce passage de la responsabilité subie à la responsabilité assumée.

\textbf{Barème indicatif (sur 20) :} analyse du sujet et problématique : 4 pts ; première partie (irresponsabilité) : 4 pts ; deuxième partie (Freud, refoulement) : 4 pts ; synthèse (distinction des deux responsabilités) : 4 pts ; contre-exemple, langue, cohérence : 4 pts.

\textbf{Points de vigilance :} mobiliser explicitement \emph{inconscient} et \emph{refoulement} (exigence du sujet) ; traiter le contre-exemple au lieu de l'esquiver ; ne jamais dire que le rêveur est coupable de ses rêves (confondre responsabilité morale et contenu psychique).
\end{corrige}

\begin{corrige}[Exercice 11 — Explication de texte : Freud, \emph{Une difficulté de la psychanalyse}]
\textbf{1. Thèse du texte.} La psychanalyse inflige à l'amour-propre humain une troisième et plus grave blessure : elle démontre que le moi n'est pas maître chez lui, c'est-à-dire que la conscience ne contrôle pas l'essentiel de la vie psychique, dominée par l'inconscient.

\textbf{2. Les trois blessures.} La \emph{première blessure} est cosmologique : Copernic (et Galilée) montrent que la Terre n'est pas le centre de l'univers ; l'homme perd sa place centrale dans le cosmos. La \emph{deuxième blessure} est biologique : Darwin établit que l'homme descend d'autres espèces animales ; il perd sa prétendue supériorité radicale sur la nature. La \emph{troisième blessure}, « la plus sensible », est psychologique : Freud montre que l'homme n'est même pas maître dans sa propre vie psychique — il perd la maîtrise de lui-même, ce qui le touche au plus intime.

\textbf{3. Analyse de la métaphore.} Le « \emph{moi} » désigne l'instance consciente, ce par quoi le sujet se dit « je » et se croit autonome. La « \emph{maison} » figure l'appareil psychique tout entier, y compris ses régions obscures (l'inconscient, le Ça). L'« absence de maîtrise » signifie que le moi ne reçoit que des « renseignements rares et fragmentaires » sur ce qui se passe « à son insu » : désirs refoulés, conflits, pulsions décident largement de nos pensées et de nos actes. Le moi ressemble à un locataire qui croit être propriétaire alors que d'autres forces habitent et gouvernent la maison.

\textbf{4. Portée philosophique.} Le texte ruine la thèse cartésienne de la \emph{transparence de la conscience à elle-même} : contre le \emph{cogito} (« je pense, donc je suis »), Freud soutient que « je » ne suis pas ce que je pense être. La connaissance de soi n'est plus un donné immédiat mais une enquête indirecte, passant par l'interprétation des rêves, des lapsus et des symptômes. La subjectivité devient un problème, non une évidence.

\textbf{5. Discussion (modèle de paragraphe, quinze à vingt lignes).} Cette conception semble d'abord ruiner la liberté : si mes actes sont déterminés par des forces inconscientes, comment pourrais-je en répondre ? Nietzsche renforce ce soupçon en décrivant la conscience comme une surface trompeuse masquant les instincts. Pourtant, cette conclusion serait hâtive. D'abord, déterminisme n'est pas fatalité : connaître les déterminations qui pèsent sur moi, c'est déjà commencer à m'en libérer — c'est le sens de la formule freudienne « Là où était le Ça, le Moi doit advenir » : la cure vise précisément à \emph{étendre} le domaine de la conscience et donc de la liberté. Ensuite, la psychanalyse elle-même suppose la liberté : elle s'adresse à un sujet capable de parler, d'associer, de s'interroger et de choisir de changer ; un être totalement déterminé ne pourrait ni mentir sur ses motifs ni entreprendre une analyse. Enfin, Kant rappelle que la responsabilité ne suppose pas une liberté absolue mais la capacité de se juger soi-même : la conscience morale, même limitée, suffit à fonder l'imputabilité — le droit ne déclare irresponsable que celui dont le discernement est \emph{aboli}, non celui qui est simplement influencé. La découverte de l'inconscient ne supprime donc pas la liberté ; elle la redéfinit : non plus une transparence souveraine, mais une conquête permanente de lucidité. La liberté cesse d'être un donné pour devenir une tâche, ce qui la rend plus exigeante, non moins réelle.

\textbf{Barème indicatif (sur 20) :} thèse dégagée : 3 pts ; explication des trois blessures : 4 pts ; analyse de la métaphore : 4 pts ; portée philosophique (opposition à Descartes) : 4 pts ; discussion argumentée avec un auteur du programme : 5 pts.

\textbf{Points de vigilance :} suivre le mouvement du texte sans le paraphraser ; citer le texte pour l'expliquer ; dans la discussion, éviter les deux écueils symétriques — « Freud détruit toute liberté » et « la liberté reste intacte » : la réponse attendue est une redéfinition de la liberté comme tâche de lucidité.
\end{corrige}

\newpage

\coursec{Fiche bilan}

\begin{popfigure}
\begin{tikzpicture}[mindmap, grow cyclic, every node/.style={concept, text=white, font=\small\bfseries},
root concept/.append style={concept color=popdark, font=\large\bfseries},
level 1/.append style={level distance=3.6cm, sibling angle=51.5},
level 2/.append style={level distance=2.2cm, sibling angle=45, font=\scriptsize}]
\node[root concept]{La conscience et l'inconscient}
  child[concept color=popblue]{ node{Définitions}
    child{ node{Conscience : présence à soi et au monde} }
    child{ node{Inconscient : psychisme caché} }
  }
  child[concept color=popgreen]{ node{Nature de la conscience}
    child{ node{Conscience spontanée / réfléchie} }
    child{ node{Descartes : \emph{cogito}} }
  }
  child[concept color=poporange]{ node{Valeur de la conscience}
    child{ node{Liberté, morale, connaissance} }
    child{ node{Kant, Sartre} }
  }
  child[concept color=poppink]{ node{Limites de la conscience}
    child{ node{Illusions, oublis, passions} }
    child{ node{Spinoza, Alain} }
  }
  child[concept color=poppurple]{ node{L'inconscient freudien}
    child{ node{Ça, Moi, Surmoi} }
    child{ node{Pulsions, refoulement} }
  }
  child[concept color=popgold]{ node{Psychanalyse}
    child{ node{Valeur : cure, interprétation} }
    child{ node{Limites : scientificité contestée} }
  }
  child[concept color=popmuted]{ node{Méthode Bac}
    child{ node{Problématiser : la conscience suffit-elle ?} }
    child{ node{Dissertation en 3 parties} }
  };
\end{tikzpicture}
\legende{Carte mentale de la leçon : les sept branches essentielles à maîtriser pour l'examen.}
\end{popfigure}

\begin{bilan}
\textbf{Définitions clés}
\begin{itemize}
\item \cle{Conscience} : faculté par laquelle l'homme se saisit lui-même et saisit le monde ; elle est à la fois \emph{connaissance} (je sais ce que je fais) et \emph{présence à soi} (je sais que j'existe).
\item \cle{Conscience spontanée} : conscience immédiate, prise dans l'action. \cle{Conscience réfléchie} : retour de l'esprit sur lui-même (réflexion).
\item \cle{Inconscient} (Freud) : ensemble des contenus psychiques (pulsions, désirs refoulés, souvenirs) qui échappent à la conscience mais influencent la conduite.
\item \cle{Refoulement} : mécanisme qui maintient hors de la conscience les désirs inacceptables.
\item \cle{Psychanalyse} : méthode fondée par Freud pour explorer l'inconscient (libre association, interprétation des rêves, des lapsus et des actes manqués).
\end{itemize}

\textbf{Thèses et auteurs à connaître par c\oe ur}
\begin{center}
\begin{tabularx}{\linewidth}{|l|X|}
\hline
\rowcolor{popblueL}\textbf{Auteur} & \textbf{Thèse} \\
\hline
Descartes & \og Je pense, donc je suis \fg{} : la conscience est le fondement certain de toute connaissance et de l'existence. \\
\hline
Spinoza & Les hommes se croient libres parce qu'ils ont conscience de leurs désirs mais ignorent les causes qui les déterminent : la conscience peut être une illusion. \\
\hline
Kant & La conscience morale rend l'homme responsable et autonome. \\
\hline
Sartre & La conscience est liberté : l'homme est condamné à être libre. \\
\hline
Freud & Le Moi \og n'est pas maître dans sa propre maison \fg{} : l'inconscient gouverne une grande partie de nos actes. \\
\hline
Alain & La psychanalyse risque d'excuser l'homme au lieu de le responsabiliser. \\
\hline
\end{tabularx}
\end{center}

\textbf{Structure freudienne de la personnalité} : le \cle{Ça} (pulsions), le \cle{Moi} (instance de réalité, médiateur), le \cle{Surmoi} (interdits moraux et sociaux intériorisés).

\textbf{Méthodes express}
\begin{enumerate}
\item Pour problématiser : opposer la thèse cartésienne (toute-puissance de la conscience) à la thèse freudienne (l'inconscient nous détermine), puis chercher une synthèse.
\item En dissertation : I. valeur de la conscience ; II. ses limites et la découverte de l'inconscient ; III. bilan critique de la psychanalyse.
\item Toujours illustrer par un exemple concret de la vie courante (lapsus, oubli, rêve, acte impulsif).
\end{enumerate}
\end{bilan}

\begin{aretenir}[Checklist avant le Bac]
\begin{itemize}
\item[$\square$] Je sais définir la conscience et distinguer conscience spontanée et conscience réfléchie.
\item[$\square$] Je connais la thèse du \emph{cogito} de Descartes et sa portée.
\item[$\square$] Je sais expliquer en quoi la conscience fonde la liberté, la morale et la connaissance.
\item[$\square$] Je connais les limites de la conscience (illusions, oublis, déterminismes) avec l'exemple de Spinoza.
\item[$\square$] Je sais définir l'inconscient freudien et le refoulement.
\item[$\square$] Je connais les trois instances : Ça, Moi, Surmoi.
\item[$\square$] Je sais citer les manifestations de l'inconscient : rêves, lapsus, actes manqués.
\item[$\square$] Je connais la valeur de la psychanalyse (méthode d'exploration et de guérison).
\item[$\square$] Je connais ses limites (scientificité contestée, risque de déresponsabilisation).
\item[$\square$] Je sais construire une problématique : \og la conscience suffit-elle à expliquer l'homme ? \fg
\item[$\square$] Je sais rédiger une dissertation en trois parties avec exemples camerounais et vie courante.
\item[$\square$] Je sais justifier ma conclusion personnelle de manière argumentée.
\end{itemize}
\end{aretenir}

\findecours

\end{document}
